\chapter{镜像对称的 Strominger--Yau--Zaslow 图景}\label{chap:6}

\SourceParagraph{ch06-intro-001}{431}
回顾一下，在作为概览的\BookChapterRef{1}中，已经非正式地介绍了 Strominger--Yau--Zaslow（SYZ）猜想的物理背景，并在\BookRef{conj:1.5}{猜想1.5}中陈述了这一猜想。现在，将系统地考察从 SYZ 猜想发展起来的镜像对称图景。首先，研究 McLean 关于特殊 Lagrange 子流形模空间的理论。这将导出 SYZ 纤维化底空间上的某些关键结构；最终会看到，底空间上的这些结构，实际上比猜想中的特殊 Lagrange 纤维化本身更重要。将解释如何从这些结构导出一个镜像对称的简化模型，并从这一观点说明镜像对称的各个方面。特别地，复结构与 Kähler 结构的交换会变得清楚，而后可以将其推广为 Hitchin 广义复结构的交换。还将在这一半平坦背景下，考察 A 膜与 B 膜的相互交换。最后，考虑一些带有奇异纤维的环面纤维化，它们提供更有意思的例子；要把镜像对称中最有意思的例子纳入讨论，就必须考虑这些纤维化。特别地，将在拓扑层面详细描述 SYZ 对偶。

\section{特殊 Lagrange 子流形的模空间}\label{sec:6.1}

\SourceParagraph{ch06-6.1-intro-001}{431}
正如\BookSectionRef{sec:1.3}{1.3}所说明的，D 膜的物理导出了 Strominger--Yau--Zaslow 猜想；如今，它已经是一个相当精确的数学陈述。要开始从数学角度思考这一猜想，需要考察特殊 Lagrange 子流形模空间上出现的自然结构。McLean\BookCite{ref356}与 Hitchin\BookCite{ref232}的奠基性工作首先研究了这些结构。以下，$X$ 为复 $n$ 维 Calabi--Yau 流形，带有处处非零的全纯 $n$ 形式 $\Omega$，以及对应于 Ricci 平坦度量的 Kähler 形式 $\omega$。回顾一下，子流形 $M\subseteq X$ 为特殊 Lagrange 子流形，是指 $\dim_{\mathbb R}M=\dim_{\mathbb C}X$，即 $M$ 的实维数为 $X$ 实维数的一半；$\omega|_M=0$，即 $M$ 为 Lagrange 子流形；并且 $\operatorname{Im}\Omega|_M=0$。采用 $\Omega$ 与 $\omega$ 相容的通常归一化，并以 $\operatorname{Re}\Omega|_M$ 为 $M$ 的体积形式来选取定向。

\subsection{McLean 定理}\label{sec:6.1.1}

\SourceParagraph{ch06-6.1.1-001}{431}
在文献\BookCite{ref356}中，McLean 证明了特殊 Lagrange 子流形的模空间没有障碍，并确定了它的切空间。具体地，令 $X$ 为 Calabi--Yau 流形，$M\subseteq X$ 为特殊 Lagrange 子流形。$M$ 在 $X$ 内的全部形变所构成空间的切空间，由 $M$ 在 $X$ 中的法丛的整体截面 $\Gamma(M,N_{M/X})$ 给出。由于 $M$ 为 Lagrange 子流形，给定法向量场 $\nu\in\Gamma(M,N_{M/X})$，就得到 $M$ 上良定义的1形式 $\iota(\nu)\omega$；其中，$\iota$ 表示将向量场 $\nu$ 与形式 $\omega$ 作内乘。具体地，对于 $M$ 的切向量 $w$，有
\[
\bigl(\iota(\nu)\omega\bigr)(w):=\omega(\nu,w).
\]
它不依赖于把 $\nu$ 提升为 $X$ 中向量场时的选择。同样，由于 $\operatorname{Im}\Omega|_M=0$，$\iota(\nu)\operatorname{Im}\Omega$ 是 $M$ 上良定义的 $(n-1)$ 形式。

\SourceParagraph{ch06-6.1.1-002}{432}
特殊 Lagrange 子流形的一个重要性质，可以由下面这个简单的线性代数事实导出。

\SourceParagraph{ch06-6.1.1-003}{432}
\BookStatementNumber{1}
\begin{exercise}\label{exer:6.1}
考虑 $\mathbb C^n$，赋予辛形式
\[
\omega=\frac{\mathrm i}{2}\sum_{i=1}^n\mathrm dz_i\wedge\mathrm d\bar z_i
\]
与全纯 $n$ 形式 $\Omega=\mathrm dz_1\wedge\cdots\wedge\mathrm dz_n$。注意，$\operatorname{SU}(n)$ 是 $\operatorname{GL}(n,\mathbb C)$ 中保持 $\omega$ 与 $\Omega$ 的子群。证明：若 $V\subseteq\mathbb C^n$ 是 $n$ 维实子空间，且 $\omega|_V=\operatorname{Im}\Omega|_V=0$，则 $V$ 是子空间 $\mathbb R^n\subseteq\mathbb C^n$，即由 $\operatorname{Im}z_i=0$，$i=1,\ldots,n$ 定义的子空间，在 $\operatorname{SU}(n)$ 的某个元素作用下的像。
\end{exercise}

\SourceParagraph{ch06-6.1.1-004}{432}
由此可以得到一个关键计算。若 $M\subseteq X$ 为特殊 Lagrange 子流形，则对于任意点 $m\in M$，都能在 $X$ 于 $m$ 点的切空间与 $\mathbb C^n$ 之间找到同构 $u\colon T_mX\to\mathbb C^n$，使得在 $m$ 点有
\[
u^*\!\left(\frac{\mathrm i}{2}\sum_{i=1}^n\mathrm dz_i\wedge\mathrm d\bar z_i\right)=\omega,
\qquad
u^*(\mathrm dz_1\wedge\cdots\wedge\mathrm dz_n)=\Omega,
\]
并且 $u(T_mM)=\mathbb R^n\subseteq\mathbb C^n$。写成 $z_i=x_i+\mathrm i y_i$，并将 $u(T_mM)$ 的法向量写成
\[
\nu=\sum_{i=1}^na_i\frac{\partial}{\partial y_i},
\]
则在 $m$ 点有
\[
\bigl(\iota(\nu)\omega\bigr)|_M=-\sum_{i=1}^na_i\,\mathrm dx_i,
\]
而
\[
\begin{aligned}
\bigl(\iota(\nu)\operatorname{Im}\Omega\bigr)|_M
&=\sum_{i=1}^n(-1)^{i+1}a_i\,
\mathrm dx_1\wedge\cdots\wedge\widehat{\mathrm dx_i}\wedge\cdots\wedge\mathrm dx_n\\
&=\sum_{i=1}^na_i*\mathrm dx_i,
\end{aligned}
\]
其中，$*$ 为 Hodge 星算子。因此，得到
\[
*\iota(\nu)\omega=-\iota(\nu)\operatorname{Im}\Omega.
\]

\SourceParagraph{ch06-6.1.1-005}{432}
现在，把 $M$ 在 $X$ 中的一个小管状邻域，与 $N_{M/X}$ 零截面的一个管状邻域识别。这样，在 $M$ 的法向量场空间 $\Gamma(M,N_{M/X})$ 中，存在零截面的一个开邻域 $U$，对于其中每个法向量场 $\nu$，都可以定义映射 $\exp_\nu\colon M\to X$。在上述管状邻域的识别下，这个映射把 $m\in M$ 映到 $\nu(m)$，也就是沿 $\nu$ 指定的方向形变 $M$。于是可以定义
\[
F\colon U\longrightarrow\Gamma(M,\Omega_M^n)\oplus\Gamma(M,\Omega_M^2),
\qquad
F(\nu)=\bigl(\exp_\nu^*(\operatorname{Im}\Omega),\exp_\nu^*\omega\bigr).
\]
这里，值域分别是 $n$ 形式空间与2形式空间的直和。$M$ 附近的特殊 Lagrange 子流形的局部模空间，由 $F^{-1}(0)$ 给出。

\SourceParagraph{ch06-6.1.1-006}{433}
通过计算 $F'(0)$，可以求出 $F^{-1}(0)$ 的切空间。$F'(0)$ 是线性映射
\[
F'(0)\colon\Gamma(M,N_{M/X})
\longrightarrow\Gamma(M,\Omega_M^n)\oplus\Gamma(M,\Omega_M^2),
\]
并且
\[
\begin{aligned}
F'(0)(\nu)
&=\left.\frac{\mathrm d}{\mathrm dt}F(t\nu)\right|_{t=0}\\
&=\bigl((\mathcal L_\nu\operatorname{Im}\Omega)|_M,
          (\mathcal L_\nu\omega)|_M\bigr)\\
&=\bigl((\iota(\nu)\mathrm d\operatorname{Im}\Omega
        +\mathrm d(\iota(\nu)\operatorname{Im}\Omega))|_M,\;
        (\iota(\nu)\mathrm d\omega+\mathrm d(\iota(\nu)\omega))|_M\bigr)\\
&=\bigl(\mathrm d(\iota(\nu)\operatorname{Im}\Omega)|_M,
         \mathrm d(\iota(\nu)\omega)|_M\bigr).
\end{aligned}
\]
这里，$\mathcal L_\nu$ 为关于 $\nu$ 的 Lie 导数，第三行使用的是 Lie 导数的 Cartan 公式。因此，$\nu\in\ker F'(0)$，当且仅当在 $M$ 上有
\[
\mathrm d(\iota(\nu)\omega)=0,
\qquad
\mathrm d(\iota(\nu)\operatorname{Im}\Omega)=0.
\]
后一条件等价于 $\mathrm d*\bigl(\iota(\nu)\omega\bigr)=0$。因此，$\iota(\nu)\omega$ 必须是 $M$ 上的调和1形式。于是，$F^{-1}(0)$ 的切空间为 $\mathcal H^1(M,\mathbb R)$，即 $M$ 上的调和1形式空间。这里不展开分析细节；应用隐函数定理，便得到 McLean 定理\BookCite{ref356}。

\SourceParagraph{ch06-6.1.1-007}{433}
\BookStatementNumber{2}
\begin{theorem}\label{thm:6.2}
令 $M\subseteq X$ 为紧、无边界的特殊 Lagrange 子流形。$M$ 在 $X$ 中的特殊 Lagrange 形变空间，是一个流形 $B$；在对应于 $M$ 的点 $[M]\in B$，其切空间同构于 $\mathcal H^1(M,\mathbb R)$，即 $M$ 上的调和1形式空间。
\end{theorem}

\SourceParagraph{ch06-6.1.1-008}{433}
现在，按如下方式考虑由流形 $B$ 参数化的特殊 Lagrange 子流形族。令 $\mathcal U\subseteq B\times X$ 为子流形，令 $p_1\colon\mathcal U\to B$、$p_2\colon\mathcal U\to X$ 为投影。假定对所有 $y\in B$，$p_2(p_1^{-1}(y))$ 都是 $X$ 的特殊 Lagrange 子流形。对于每个 $y\in B$，都有映射
\[
T_yB\longrightarrow\mathcal H^1(p_1^{-1}(y),\mathbb R),
\]
其构造本质上就是前面所描述的。具体地，给定 $\nu\in T_yB$，把它提升为 $p_1^{-1}(y)$ 的法向量场，仍记为 $\nu$。于是，$p_{2*}\nu$ 是 $p_2(p_1^{-1}(y))$ 的法向量场，对应于形变方向 $\nu$。此时，$\iota(p_{2*}\nu)\omega=\iota(\nu)p_2^*\omega$ 是对应于 $\nu$ 的调和1形式。现在假定，对于每个 $y\in B$，这个映射都是同构。这样，就得到 $B$ 上的下列结构。

\SourceParagraph{ch06-6.1.1-009}{434}
\textbf{McLean 度量。}$\mathcal H^1(p_1^{-1}(y),\mathbb R)$ 上的 Hodge 度量
\[
g(\alpha,\beta)=\int_{p_1^{-1}(y)}\alpha\wedge*\beta
\]
在 $B$ 上诱导一个度量，可以写成
\[
g(\nu_1,\nu_2)=-\int_{p_1^{-1}(y)}
\iota(\nu_1)p_2^*\omega\wedge\iota(\nu_2)p_2^*\operatorname{Im}\Omega.
\]
Hodge 度量正定，因此得到了 $B$ 上的 Riemann 度量。

\SourceParagraph{ch06-6.1.1-010}{434}
\textbf{第一套仿射坐标。}令 $U\subseteq B$ 为一个小开集，并假定存在子流形 $\gamma_1,\ldots,\gamma_s\subseteq\mathcal U$，它们是 $U$ 上的1维闭链族；也就是说，对 $y\in U$，$\gamma_i\cap p_1^{-1}(y)$ 是 $p_1^{-1}(y)$ 中的1维闭链。进一步假定，$\{\gamma_i\cap p_1^{-1}(y)\}$ 构成 $H_1(p_1^{-1}(y),\mathbb Z)/\mathrm{tors}$ 的一组基。于是，可以考虑 $U$ 上的1形式
\[
\omega_i=p_{1*}(p_2^*\omega|_{\gamma_i}).
\]
它由沿纤维积分得到；对 $\nu\in T_yB$，有
\[
\omega_i(\nu)=\int_{p_1^{-1}(y)\cap\gamma_i}\iota(\nu)p_2^*\omega.
\]
由于 $\omega$ 为闭形式，每个 $\omega_i$ 也为闭形式。因此，局部地，可以写成 $\omega_i=\mathrm dy_i$，其中 $y_1,\ldots,y_s$ 为函数。它们具有以下性质。
\begin{itemize}
\item 在每个 $y\in U$，形式 $\omega_1,\ldots,\omega_s$ 都线性无关。否则，存在不全为零的 $a_1,\ldots,a_s\in\mathbb R$，使得对所有 $\nu\in T_yB$，有
\[
\int_{p_1^{-1}(y)\cap\sum_i a_i\gamma_i}\iota(\nu)p_2^*\omega=0.
\]
但是，$\iota(\nu)p_2^*\omega|_{p_1^{-1}(y)}$ 遍历 $\mathcal H^1(p_1^{-1}(y),\mathbb R)$ 的全部元素，因此，$\sum_i a_i\gamma_i=0$ 于 $H_1(p_1^{-1}(y),\mathbb R)$ 中。按照假设，$\gamma_i$ 线性无关，这就给出矛盾。因此，函数 $y_1,\ldots,y_s$ 在 $y$ 附近给出一套局部坐标。
\item $y_1,\ldots,y_s$ 在选定积分常数后良定义。如果改用另一组基 $\gamma_1,\ldots,\gamma_s$，那么 $y_1,\ldots,y_s$ 相应地变为
\[
y_i'=\sum_j a_{ij}y_j+b_i,
\qquad (a_{ij})_{1\leq i,j\leq s}\in\operatorname{GL}(s,\mathbb Z),
\quad b_i\in\mathbb R.
\]
因此，得到了良定义的仿射坐标，见\BookRef{def:6.3}{定义6.3}。
\end{itemize}

\SourceParagraph{ch06-6.1.1-011}{434}
\textbf{第二套仿射坐标。}对于 $U$ 上的 $(n-1)$ 维闭链族
\[
\Gamma_1,\ldots,\Gamma_s\subseteq\mathcal U,
\]
重复同样的过程，其中，$\{\Gamma_i\cap p_1^{-1}(y)\}$ 构成
\[
H_{n-1}(p_1^{-1}(y),\mathbb Z)/\mathrm{tors}
\]
的一组基。这样得到1形式 $\lambda_i=-p_{1*}(p_2^*\operatorname{Im}\Omega|_{\Gamma_i})$，等价地，
\[
\lambda_i(\nu)=-\int_{p_1^{-1}(y)\cap\Gamma_i}
\iota(\nu)p_2^*\operatorname{Im}\Omega.
\]
同样，$\lambda_1,\ldots,\lambda_s$ 是闭1形式，并且在每点都线性无关。因此，得到构成局部坐标系的函数 $\check y_1,\ldots,\check y_s$，满足 $\mathrm d\check y_i=\lambda_i$。

\subsection{仿射流形与 Legendre 变换}\label{sec:6.1.2}

\SourceParagraph{ch06-6.1.2-001}{435}
现在，以更抽象的方式考察前一小节出现的仿射坐标所给出的结构。固定
\[
M=\mathbb Z^n,\qquad N=\operatorname{Hom}_{\mathbb Z}(M,\mathbb Z),
\qquad M_{\mathbb R}=M\otimes_{\mathbb Z}\mathbb R,
\qquad N_{\mathbb R}=N\otimes_{\mathbb Z}\mathbb R.
\]
记
\[
\operatorname{Aff}(M_{\mathbb R})
=M_{\mathbb R}\rtimes\operatorname{GL}_n(\mathbb R)
\]
为 $M_{\mathbb R}$ 的仿射变换群，也就是形如 $T(x)=Ax+b$ 的映射 $T\colon M_{\mathbb R}\to M_{\mathbb R}$，其中，$A\colon M_{\mathbb R}\to M_{\mathbb R}$ 为可逆线性映射，$b\in M_{\mathbb R}$。它有子群
\[
\operatorname{Aff}(M)=M\rtimes\operatorname{GL}_n(\mathbb Z).
\]

\SourceParagraph{ch06-6.1.2-002}{435}
\BookStatementNumber{3}
\begin{definition}\label{def:6.3}
令 $B$ 为 $n$ 维流形。$B$ 上的仿射结构，由开覆盖 $\{U_i\}$ 及坐标图 $\psi_i\colon U_i\to M_{\mathbb R}$ 给出，要求转移函数 $\psi_i\circ\psi_j^{-1}$ 属于 $\operatorname{Aff}(M_{\mathbb R})$。若转移函数属于 $\operatorname{Aff}(M)$，则称这一仿射结构为整仿射结构。若 $B$ 与 $B'$ 为同维的仿射流形，$f\colon B\to B'$ 为浸入，则当 $f$ 在每个（整）仿射坐标图中都由（整）仿射映射给出时，称 $f$ 为（整）仿射映射。
\end{definition}

\SourceParagraph{ch06-6.1.2-003}{435}
由此看到，在\BookSectionRef{sec:6.1.1}{6.1.1}的情形中，$B$ 实际上带有两个不同的仿射结构。分别称它们为辛仿射结构与复仿射结构，因为它们分别由 $X$ 上的辛形式 $\omega$ 与全纯 $n$ 形式 $\Omega$ 诱导。

\SourceParagraph{ch06-6.1.2-004}{435}
关于文献\BookCite{ref232}，这里最后要说明的是，这两个仿射结构与 McLean 度量在下述意义下相容；特别地，两个仿射结构由 Legendre 变换联系起来。

\SourceParagraph{ch06-6.1.2-005}{435}
\BookStatementNumber{4}
\begin{proposition}\label{prop:6.4}
如\BookSectionRef{sec:6.1.1}{6.1.1}所述，令 $B$ 为参数化某个特殊 Lagrange 子流形 $M\subseteq X$ 的特殊 Lagrange 形变的模空间。令 $y\in B$。则 $y$ 有开邻域 $U$，带辛仿射坐标 $y_1,\ldots,y_s$，并且 $U$ 上存在凸函数 $K$，使得
\[
g\!\left(\frac{\partial}{\partial y_i},\frac{\partial}{\partial y_j}\right)
=\frac{\partial^2K}{\partial y_i\partial y_j}.
\]
此外，$\check y_i=\partial K/\partial y_i$ 构成一套复仿射坐标。如果
\[
\check K(\check y_1,\ldots,\check y_s)
=\sum_i\check y_i y_i-K(y_1,\ldots,y_s)
\]
是 $K$ 的 Legendre 变换，则
\[
y_i=\frac{\partial\check K}{\partial\check y_i},
\qquad
\frac{\partial^2\check K}{\partial\check y_i\partial\check y_j}
=g\!\left(\frac{\partial}{\partial\check y_i},\frac{\partial}{\partial\check y_j}\right).
\]
\end{proposition}

\begin{proof}
\SourceParagraph{ch06-6.1.2-006}{436}
在开邻域 $U$ 上，取\BookSectionRef{sec:6.1.1}{6.1.1}中的族 $\gamma_1,\ldots,\gamma_s$、$\Gamma_1,\ldots,\Gamma_s$，使得两组基互为 Poincaré 对偶，即对 $y\in U$，有 $\gamma_i\cdot\Gamma_j=\delta_{ij}$。令 $\gamma_1^*,\ldots,\gamma_s^*$ 与 $\Gamma_1^*,\ldots,\Gamma_s^*$，分别为 $H^1(p_1^{-1}(y),\mathbb R)$ 与 $H^{n-1}(p_1^{-1}(y),\mathbb R)$ 的对偶基。这样得到局部坐标 $y_1,\ldots,y_s$，满足 $\mathrm dy_i=\omega_i$。特别地，
\[
\delta_{ij}=\omega_i\!\left(\frac{\partial}{\partial y_j}\right)
=\int_{\gamma_i\cap p_1^{-1}(y)}
\iota\!\left(\frac{\partial}{\partial y_j}\right)\omega,
\]
所以，$\iota(\partial/\partial y_j)\omega$ 定义 $H^1(p_1^{-1}(y),\mathbb R)$ 中的上同调类 $\gamma_j^*$。这里及以下，省略形式沿 $p_2$ 的拉回记号。同样，令
\[
g_{ij}=-\int_{\Gamma_i\cap p_1^{-1}(y)}
\iota\!\left(\frac{\partial}{\partial y_j}\right)\operatorname{Im}\Omega;
\]
则 $-\iota(\partial/\partial y_j)\operatorname{Im}\Omega$ 定义上同调类
\[
\sum_i g_{ij}\Gamma_i^*\in H^{n-1}(p_1^{-1}(y),\mathbb R),
\]
并且 $\lambda_i=\sum_jg_{ij}\,\mathrm dy_j$。因此，
\[
\begin{aligned}
g\!\left(\frac{\partial}{\partial y_j},\frac{\partial}{\partial y_k}\right)
&=-\int_{p_1^{-1}(y)}
\iota\!\left(\frac{\partial}{\partial y_j}\right)\omega
\wedge\iota\!\left(\frac{\partial}{\partial y_k}\right)\operatorname{Im}\Omega\\
&=g_{jk}.
\end{aligned}
\]

\SourceParagraph{ch06-6.1.2-007}{436}
另一方面，令 $\check y_1,\ldots,\check y_s$ 为满足 $\mathrm d\check y_i=\lambda_i$ 的坐标。于是，
\[
\frac{\partial\check y_i}{\partial y_j}
=g_{ij}=g_{ji}=\frac{\partial\check y_j}{\partial y_i},
\]
所以，$\sum_i\check y_i\,\mathrm dy_i$ 是闭1形式。因此，局部地存在函数 $K$，使得 $\partial K/\partial y_i=\check y_i$，并且
\[
\frac{\partial^2K}{\partial y_i\partial y_j}
=g\!\left(\frac{\partial}{\partial y_i},\frac{\partial}{\partial y_j}\right).
\]
一个简单计算便可确认，$\partial\check K/\partial\check y_i=y_i$。另一方面，
\[
\begin{aligned}
g\!\left(\frac{\partial}{\partial\check y_i},\frac{\partial}{\partial\check y_j}\right)
&=g\!\left(\sum_k\frac{\partial y_k}{\partial\check y_i}\frac{\partial}{\partial y_k},
\sum_l\frac{\partial y_l}{\partial\check y_j}\frac{\partial}{\partial y_l}\right)\\
&=\sum_{k,l}\frac{\partial y_k}{\partial\check y_i}
\frac{\partial y_l}{\partial\check y_j}
g\!\left(\frac{\partial}{\partial y_k},\frac{\partial}{\partial y_l}\right)\\
&=\sum_{k,l}\frac{\partial y_k}{\partial\check y_i}
\frac{\partial y_l}{\partial\check y_j}\frac{\partial\check y_k}{\partial y_l}\\
&=\frac{\partial y_j}{\partial\check y_i}
=\frac{\partial^2\check K}{\partial\check y_i\partial\check y_j}.
\end{aligned}
\]
\end{proof}

\SourceParagraph{ch06-6.1.2-008}{436}
能够以更加不依赖坐标的方式进行这些构造，会很有用。下面给出这样的表述。

\SourceParagraph{ch06-6.1.2-009}{436}
\BookStatementNumber{5}
\begin{proposition}\label{prop:6.5}
令 $\pi\colon\widetilde B\to B$ 为（整）仿射流形 $B$ 的万有覆盖，因而在 $\widetilde B$ 上诱导一个（整）仿射结构。则存在（整）仿射映射 $\delta\colon\widetilde B\to M_{\mathbb R}$，称为展开映射；任意两个这样的映射只相差一个（整）仿射变换。
\end{proposition}
\begin{proof}
\SourceParagraph{ch06-6.1.2-010}{437}
这是一个标准结果，证明见文献\BookCite{ref184}第641页。只需把仿射坐标图粘合起来。
\end{proof}

\SourceParagraph{ch06-6.1.2-011}{437}
注意，展开映射不必是单射，也不必是覆盖映射；一般而言，它只是浸入。

\SourceParagraph{ch06-6.1.2-012}{437}
\BookStatementNumber{6}
\begin{definition}\label{def:6.6}
基本群 $\pi_1(B)$ 通过覆叠变换作用于 $\widetilde B$；对 $\gamma\in\pi_1(B)$，令 $\Psi_\gamma\colon\widetilde B\to\widetilde B$ 为相应的覆叠变换。由展开映射的唯一性，存在 $\rho(\gamma)\in\operatorname{Aff}(M_{\mathbb R})$，使得
\[
\rho(\gamma)\circ\delta\circ\Psi_\gamma=\delta.
\]
映射 $\rho\colon\pi_1(B)\to\operatorname{Aff}(M_{\mathbb R})$，称为完整表示。若仿射结构为整仿射结构，则 $\operatorname{im}\rho\subseteq\operatorname{Aff}(M)$。
\end{definition}

\SourceParagraph{ch06-6.1.2-013}{437}
采用约定 $\Psi_{\gamma_1\gamma_2}=\Psi_{\gamma_2}\circ\Psi_{\gamma_1}$，于是，上面定义的 $\rho$ 是群同态。完整表示还可以用稍有不同的方式理解。$TB$ 上有一个平坦联络，使得只要 $y_1,\ldots,y_n$ 是仿射坐标系，$\partial/\partial y_1,\ldots,\partial/\partial y_n$ 就是平坦截面。那么，$\rho$ 的线性部分就是这一仿射联络的完整表示；也就是说，$\rho(\gamma)$ 的线性部分，由 $TB$ 中沿 $\gamma^{-1}$ 的平行移动给出。表示 $\rho$ 本身，可以看作 $TB$ 上某个自然仿射联络的完整表示，细节见文献\BookCite{ref184}。注意，在坐标变换的线性部分为整矩阵时，$\rho$ 的线性部分正是切丛格局部系统 $\Lambda$ 的单值变换；这一局部系统将在\BookSectionRef{sec:6.2.1}{6.2.1}中明确构造。

\SourceParagraph{ch06-6.1.2-014}{437}
\BookStatementNumber{7}
\begin{example}\label{ex:6.7}
\begin{enumerate}[label=(\arabic*)]
\item 对于格 $\Lambda\subseteq M_{\mathbb R}$，环面 $T^n=M_{\mathbb R}/\Lambda$ 有一个自然的仿射结构，其展开映射由恒等映射 $\operatorname{id}\colon M_{\mathbb R}\to M_{\mathbb R}$ 给出。这里，对 $\lambda\in\pi_1(T^n)=\Lambda$，$\rho(\lambda)$ 是平移 $-\lambda$。因此，这一仿射结构为整仿射结构，当且仅当 $\Lambda\subseteq M$。
\end{enumerate}

\SourceParagraph{ch06-6.1.2-015}{437}
\begin{enumerate}[label=(\arabic*),start=2]
\item 这个例子来自文献\BookCite{ref033}。取 $M=\mathbb Z^2$，考虑如下定义的子群 $G\subseteq\operatorname{Aff}(M_{\mathbb R})$：
\[
G=\left\{A\in\operatorname{Aff}(M_{\mathbb R})\;\middle|\;
\begin{aligned}
A(m_1,m_2)&=\bigl(m_1+vm_2+u+v(v-1)/2,\;m_2+v\bigr),\\[-2pt]
&\hspace{40mm}u,v\in\mathbb R
\end{aligned}\right\}.
\]
这并非文献\BookCite{ref033}所给出的原有形式；这里将 $G$ 用平移 $(0,-1/2)$ 共轭，以获得更好的整性。$G$ 同构于 $\mathbb R^2$。如果选取任意格 $\Gamma\subseteq G$，则 $\Gamma$ 的作用是正常不连续的，因此，$M_{\mathbb R}/\Gamma$ 是仿射流形，在拓扑上为二维环面。对于由 $M_{\mathbb R}$ 对正常不连续的群作用取商得到的二维环面，除第一个例子之外，这是唯一的另一种仿射结构，见文献\BookCite{ref033}的定理4.5。注意，相对于整结构 $M\subseteq M_{\mathbb R}$，这个仿射结构为整仿射结构，当且仅当
\[
\Gamma\subseteq\left\{A\in\operatorname{Aff}(M_{\mathbb R})\;\middle|\;
\begin{aligned}
A(m_1,m_2)&=\bigl(m_1+vm_2+u+v(v-1)/2,\;m_2+v\bigr),\\[-2pt]
&\hspace{40mm}u\in\mathbb Z,\quad v\in\mathbb Z
\end{aligned}\right\}.
\]
\end{enumerate}
\end{example}

\SourceParagraph{ch06-6.1.2-016}{438}
\BookStatementNumber{8}
\begin{note}[numbered]\label{note:6.8}
事实上，可以通过给出展开映射，指定 $B$ 上的仿射结构。换言之，若有浸入 $\delta\colon\widetilde B\to M_{\mathbb R}$，以及某个表示
\[
\rho\colon\pi_1(B)\longrightarrow\operatorname{Aff}(M_{\mathbb R}),
\]
使得对所有 $\gamma\in\pi_1(B)$，都有
\[
\rho(\gamma)\circ\delta\circ\Psi_\gamma=\delta,
\]
则 $\delta$ 决定 $B$ 上的一个仿射结构。此外，若 $\rho$ 的像包含于 $\operatorname{Aff}(M)$，则这一仿射结构为整仿射结构。
\end{note}

\SourceParagraph{ch06-6.1.2-017}{438}
现在，以不依赖坐标的方式描述 Legendre 变换。假定 $B$ 为仿射流形，带有 Hessian 形式的度量；也就是说，局部地存在函数 $K$，使得 $g_{ij}=\partial^2K/\partial y_i\partial y_j$，其中，$y_1,\ldots,y_n$ 为局部仿射坐标。注意，$K$ 不依赖于具体的坐标选择。不过，它只在相差一个仿射线性函数的意义下良定义，即允许加上形如 $\sum_i a_i y_i+b$ 的函数。任意两个这样的函数 $K$，都相差一个仿射线性函数。因此，可以把它们粘合，得到良定义的 $K\colon\widetilde B\to\mathbb R$。注意，
\[
K-K\circ\Psi_\gamma=:\alpha(\gamma)
\]
是仿射线性函数 $\alpha(\gamma)\colon\widetilde B\to\mathbb R$；换言之，$\alpha(\gamma)$ 是展开映射与某个仿射线性函数 $M_{\mathbb R}\to\mathbb R$ 的复合。

\SourceParagraph{ch06-6.1.2-018}{438}
展开映射 $\delta\colon\widetilde B\to M_{\mathbb R}$ 的微分，给出余切丛之间的同构
\[
\delta^*\colon\delta^*T^*M_{\mathbb R}\longrightarrow T^*\widetilde B.
\]
由于 $T^*M_{\mathbb R}=M_{\mathbb R}\times N_{\mathbb R}$，$\delta^*$ 给出 $\widetilde B\times N_{\mathbb R}$ 与 $T^*\widetilde B$ 的同构。令 $q\colon\widetilde B\times N_{\mathbb R}\to N_{\mathbb R}$ 为投影。$K$ 的微分 $\mathrm dK$ 是 $T^*\widetilde B$ 的截面，即映射 $\mathrm dK\colon\widetilde B\to T^*\widetilde B$。因此，可以把 $(\delta^*)^{-1}(\mathrm dK)$ 看作 $\delta^*T^*M_{\mathbb R}$ 的截面。于是，可以把
\[
\check\delta:=q\circ(\delta^*)^{-1}(\mathrm dK)
\]
看作函数 $\check\delta\colon\widetilde B\to N_{\mathbb R}$。在这些识别下，它就是微分 $\mathrm dK$。由于 $K$ 的 Hessian 矩阵正定，$\check\delta$ 是浸入。

\SourceParagraph{ch06-6.1.2-019}{438}
\BookStatementNumber{9}
\begin{proposition}\label{prop:6.9}
令 $\widetilde\rho$ 为 $\rho$ 与自然投影 $\operatorname{Aff}(M_{\mathbb R})\to\operatorname{GL}(M_{\mathbb R})$ 的复合，即 $\rho$ 的线性部分。那么，
\[
{}^t\widetilde\rho(\gamma^{-1})\circ\check\delta\circ\Psi_\gamma
+\mathrm d\alpha(\gamma)=\check\delta,
\]
并且，$\pi_1(B)$ 通过仿射变换作用于 $\widetilde B$ 上的新仿射结构。对这一作用取商，就得到 $B$ 上的新仿射结构，记为 $\check B$。因此，由 $\check\delta$ 给出的仿射结构的完整表示 $\check\rho\colon\pi_1(\check B)\to\operatorname{Aff}(N_{\mathbb R})$，其线性部分为 $\widetilde\rho^\vee$，即表示 $\widetilde\rho$ 的对偶。
\end{proposition}
\begin{proof}
\SourceParagraph{ch06-6.1.2-020}{439}
由于 $\delta\circ\Psi_\gamma=\rho(\gamma^{-1})\circ\delta$，有交换图
\[
\BookImage{02_Unnumbered_Diagrams/Chapter_06/Ch06_Diagram_PDFp439_01_Bookp429.pdf}
\]
但是，$\rho(\gamma^{-1})^*$ 在 $N_{\mathbb R}$ 上的作用，就是 $ {}^t\widetilde\rho(\gamma^{-1})$。因此，对于 $\widetilde B$ 上的任意1形式 $\omega$，在点 $b\in\widetilde B$ 有
\[
\begin{aligned}
\bigl[q\circ(\delta^*)^{-1}(\Psi_\gamma^*\omega)\bigr](b)
&=\bigl[q\circ\rho(\gamma^{-1})^*\circ(\delta^*)^{-1}(\omega)\bigr](\Psi_\gamma(b))\\
&={}^t\widetilde\rho(\gamma^{-1})
\bigl[q\circ(\delta^*)^{-1}(\omega)\bigr](\Psi_\gamma(b)).
\end{aligned}
\]
由链式法则，得到形式的等式
\[
\Psi_\gamma^*(\mathrm dK)=\mathrm d(K\circ\Psi_\gamma).
\]
因此，点 $\Psi_\gamma(b)$ 处的形式 $\mathrm dK$，等于将拉回 $\Psi_{\gamma^{-1}}^*$ 应用于点 $b\in\widetilde B$ 处的形式 $\mathrm d(K\circ\Psi_\gamma)$。这转化为恒等式
\[
\begin{aligned}
\check\delta(\Psi_\gamma(b))
&=\bigl[q\circ(\delta^*)^{-1}
\bigl(\Psi_{\gamma^{-1}}^*(\mathrm d(K\circ\Psi_\gamma))\bigr)\bigr](\Psi_\gamma(b))\\
&={}^t\widetilde\rho(\gamma)
\bigl[q\circ(\delta^*)^{-1}(\mathrm dK-\mathrm d\alpha(\gamma))\bigr](b)\\
&={}^t\widetilde\rho(\gamma)
\bigl(\check\delta(b)-\mathrm d\alpha(\gamma)\bigr).
\end{aligned}
\]
这里，$\mathrm d\alpha(\gamma)$ 自然地识别为 $N_{\mathbb R}$ 的一个元素。
\end{proof}

\SourceParagraph{ch06-6.1.2-021}{439}
还可以定义 $K$ 的 Legendre 变换：对 $b\in\widetilde B$，令
\[
\check K(b)=\langle\check\delta(b),\delta(b)\rangle-K(b).
\]
于是，$\check K$ 在 $\check B$ 上的 Hessian 矩阵，定义的度量与 $K$ 在 $B$ 上定义的度量相同，并且在相应识别下有 $\mathrm d\check K=\delta$。称 $(\check B,\check K)$ 为 $(B,K)$ 的 Legendre 变换。

\SourceParagraph{ch06-6.1.2-022}{439}
显然，若知道 $B$、$\check B$ 和 $g$ 三者中的任意两个，就也知道第三个。

\SourceParagraph{ch06-6.1.2-023}{439}
\BookStatementNumber{10}
\begin{example}\label{ex:6.10}
设 $B=M_{\mathbb R}/\Gamma$，如\BookRef{ex:6.7}{例题6.7}(1)中所述。势函数 $K\colon M_{\mathbb R}\to\mathbb R$ 的一种选择，是一个凸二次函数，由 $M_{\mathbb R}$ 上的内积 $\langle\cdot,\cdot\rangle$ 决定：$K(x)=\tfrac12\langle x,x\rangle$。注意，对 $\gamma\in\Gamma$，有
\[
\begin{aligned}
K(x+\gamma)
&=\tfrac12\langle x+\gamma,x+\gamma\rangle\\
&=\tfrac12\langle x,x\rangle+\langle x,\gamma\rangle
  +\tfrac12\langle\gamma,\gamma\rangle\\
&=K(x)+\langle x,\gamma\rangle+\tfrac12\langle\gamma,\gamma\rangle.
\end{aligned}
\]
在 $\Psi_\gamma(x)=x+\gamma$ 与 $\alpha(\gamma)=K-K\circ\Psi_\gamma$ 的约定下，
\[
\alpha(\gamma)(x)=-\langle x,\gamma\rangle-\tfrac12\langle\gamma,\gamma\rangle,
\]
而 $\mathrm d\alpha(\gamma)$ 是 $N_{\mathbb R}$ 中对应于 $M_{\mathbb R}$ 上泛函 $-\langle\gamma,\cdot\rangle$ 的元素。换言之，这个内积使得可以将 $M_{\mathbb R}$ 与 $N_{\mathbb R}$ 识别；事实上，容易看到，$\check\delta\colon M_{\mathbb R}\to N_{\mathbb R}$ 就是这一识别。它把 $\Gamma\subseteq M_{\mathbb R}$ 识别为
\[
\check\Gamma=\{\mathrm d\alpha(\gamma)\in N_{\mathbb R}\mid\gamma\in\Gamma\},
\]
这里，整组格向量取负号，并不改变所得格。
\end{example}

\SourceParagraph{ch06-6.1.2-024}{440}
\BookStatementNumber{11}
\begin{exercise}\label{exer:6.11}
证明：在\BookRef{ex:6.7}{例题6.7}(2)中，不存在 Hessian 矩阵处处正定的凸函数 $K\colon M_{\mathbb R}\to\mathbb R$，使得对所有 $\gamma\in\Gamma$，都有 $K-K\circ\Psi_\gamma=\alpha(\gamma)$，其中，$\alpha(\gamma)$ 为某个仿射线性函数。因此，$B=M_{\mathbb R}/\Gamma$ 不可能是某个 Calabi--Yau 流形上特殊 Lagrange 环面的模空间。
\end{exercise}

\section{半平坦 SYZ 图景}\label{sec:6.2}

\SourceParagraph{ch06-6.2-intro-001}{440}
现在，将利用前一节详细讨论的、特殊 Lagrange 纤维化底空间上的结构，描述一种简单形式的镜像对称。关于半平坦镜像对称的更多细节，见文献\BookCite{ref331}。

\subsection{基本形式}\label{sec:6.2.1}

\SourceParagraph{ch06-6.2.1-001}{440}
现在，可以利用\BookSectionRef{sec:6.1}{6.1}讨论的结构，定义一个镜像对称的简化模型。在本节中，始终固定仿射流形 $B$，并假定所有转移映射都属于 $M_{\mathbb R}\rtimes\operatorname{GL}(n,\mathbb Z)$，而非一般的 $\operatorname{Aff}(M_{\mathbb R})$。

\SourceParagraph{ch06-6.2.1-002}{440}
给定这些数据，令 $y_1,\ldots,y_n$ 为局部仿射坐标。于是，在 $TB$ 中得到一族格，由 $\partial/\partial y_1,\ldots,\partial/\partial y_n$ 生成。由于整性假设，它在整体上良定义：更换坐标会产生同一个格的另一组基，两组基由 $\operatorname{GL}(n,\mathbb Z)$ 的一个元素联系。这定义了局部系统 $\Lambda\subseteq TB$。同样，局部地，$\mathrm dy_1,\ldots,\mathrm dy_n$ 生成局部系统 $\check\Lambda\subseteq T^*B$。还记 $\Lambda_{\mathbb R}=\Lambda\otimes_{\mathbb Z}\mathbb R$、$\check\Lambda_{\mathbb R}=\check\Lambda\otimes_{\mathbb Z}\mathbb R$。它们也是分别包含在 $TB$ 与 $T^*B$ 中的局部系统，只是现在允许 $\partial/\partial y_1,\ldots,\partial/\partial y_n$ 或 $\mathrm dy_1,\ldots,\mathrm dy_n$ 的实线性组合作为局部截面。

\SourceParagraph{ch06-6.2.1-003}{440}
现在，可以定义两个环面丛：
\[
X(B):=TB/\Lambda,\qquad
\check X(B):=T^*B/\check\Lambda.
\]
它们都是 $B$ 上的 $T^n$ 丛。记相应投影为
\[
f\colon X(B)\longrightarrow B,
\qquad
\check f\colon\check X(B)\longrightarrow B,
\]
并称它们为对偶环面丛。

\SourceParagraph{ch06-6.2.1-004}{440}
这些丛 $X(B)$ 与 $\check X(B)$ 还带有一些附加结构。首先，$\check X(B)$ 自然地为辛流形：$T^*B$ 上的典范辛形式，下降为 $\check X(B)$ 上的辛形式 $\omega$。其次，$X(B)$ 带有复结构。局部地，可以用全纯坐标描述它。令 $U\subseteq B$ 为带有仿射坐标 $y_1,\ldots,y_n$ 的开集。于是，$TU$ 有坐标函数 $y_1,\ldots,y_n$ 与 $x_1=\mathrm dy_1,\ldots,x_n=\mathrm dy_n$；这里，将 $\mathrm dy_i$ 看作切向量上的线性函数。那么，
\[
q_j=\exp\bigl(2\pi\mathrm i(x_j+\mathrm i y_j)\bigr)
\]
给出 $TU/\Lambda$ 上的一套全纯坐标。

\SourceParagraph{ch06-6.2.1-005}{441}
\BookStatementNumber{12}
\begin{exercise}\label{exer:6.12}
考察坐标 $y_j$ 作整仿射变换时，坐标 $q_j$ 如何变换。注意，这个坐标变换是全纯的。
\end{exercise}

\SourceParagraph{ch06-6.2.1-006}{441}
完成这个练习，就说明这些坐标给出 $X(B)$ 上良定义的复结构。注意，还可以如下以不依赖坐标的方式描述它。展开映射的微分，将 $T\widetilde B=\widetilde B\times M_{\mathbb R}$ 映到 $TM_{\mathbb R}=M_{\mathbb R}\times M_{\mathbb R}$，从而可以拉回
\[
M_{\mathbb R}\times M_{\mathbb R}
=M_{\mathbb R}\oplus\mathrm iM_{\mathbb R}
=M_{\mathbb R}\otimes\mathbb C
\]
上的自然复结构。在这里，以纤维坐标 $x$ 与底坐标 $y$ 的顺序写成 $x+\mathrm i y$。上述坐标所描述的复结构，就是由 $M_{\mathbb R}\otimes\mathbb C$ 上典范复结构的拉回诱导的。

\SourceParagraph{ch06-6.2.1-007}{441}
此外，$X(B)$ 还带有一个自然的局部全纯 $n$ 形式 $\Omega$：
\[
\Omega=\frac{\mathrm dq_1\wedge\cdots\wedge\mathrm dq_n}{q_1\cdots q_n}.
\]
它未必在整体上良定义，不过，$\Omega\wedge\overline\Omega$ 在整体上良定义，如下面的练习所说明的。

\SourceParagraph{ch06-6.2.1-008}{441}
\BookStatementNumber{13}
\begin{exercise}\label{exer:6.13}
考察并证明：$\Omega$ 在属于 $M_{\mathbb R}\rtimes\operatorname{SL}_n(\mathbb Z)$ 的坐标变换 $y_1,\ldots,y_n$ 下保持不变，而在属于 $M_{\mathbb R}\rtimes\operatorname{GL}(n,\mathbb Z)$ 的坐标变换下，只在相差一个符号的意义下保持不变。因此，若完整表示的像包含于 $M_{\mathbb R}\rtimes\operatorname{SL}_n(\mathbb Z)$，则得到整体全纯 $n$ 形式 $\Omega$。
\end{exercise}

\SourceParagraph{ch06-6.2.1-009}{441}
现在，进一步假定 $B$ 上有一个 Hessian 形式的度量 $g$，势函数为 $K\colon\widetilde B\to\mathbb R$。那么，$X(B)$ 与 $\check X(B)$ 实际上都成为 Kähler 流形。

\SourceParagraph{ch06-6.2.1-010}{441}
\BookStatementNumber{14}
\begin{proposition}\label{prop:6.14}
$K\circ f$ 是 $X(B)$ 上的一个多值 Kähler 势，定义 Kähler 形式
\[
\omega=2\mathrm i\,\partial\bar\partial(K\circ f).
\]
记 $g_{ij}=\partial^2K/\partial y_i\partial y_j$。这一度量 Ricci 平坦，当且仅当 $\log\det(g_{ij})$ 在局部仿射坐标中为仿射线性函数。若要求自然全纯体积形式 $\Omega$ 平行，即采用相对于这一体积形式的 Calabi--Yau 归一化，则其条件为实 Monge--Ampère 方程
\[
\det\!\left(\frac{\partial^2K}{\partial y_i\partial y_j}\right)=\text{常数}.
\]
\end{proposition}
\begin{proof}
\SourceParagraph{ch06-6.2.1-011}{441}
局部地，使用仿射坐标 $(y_i)$ 与多值复坐标
\[
z_i=\frac{1}{2\pi\mathrm i}\log q_i=x_i+\mathrm i y_i,
\]
计算得
\[
\omega=2\mathrm i\,\partial\bar\partial(K\circ f)
=\frac{\mathrm i}{2}\sum_{i,j}
\frac{\partial^2K}{\partial y_i\partial y_j}\,
\mathrm dz_i\wedge\mathrm d\bar z_j,
\]
它显然正定。Ricci 形式为 $-\mathrm i\,\partial\bar\partial\log\det(g_{ij})$；由于行列式只依赖于 $y$，它为零当且仅当 $\log\det(g_{ij})$ 关于 $y$ 的 Hessian 矩阵为零。此外，$\omega^n$ 为 $\Omega\wedge\overline\Omega$ 的常数倍，当且仅当 $\det(\partial^2K/\partial y_i\partial y_j)$ 为常数。
\end{proof}

\SourceParagraph{ch06-6.2.1-012}{441}
把这个 Kähler 流形记为 $X(B,K)$。

\SourceParagraph{ch06-6.2.1-013}{441}
对偶地，有下述命题。

\SourceParagraph{ch06-6.2.1-014}{442}
\BookStatementNumber{15}
\begin{proposition}\label{prop:6.15}
在 $T^*B$ 的局部典范坐标 $y_i,\check x_i$ 中，函数
\[
z_i=\check x_i+\mathrm i\frac{\partial K}{\partial y_i}
\]
在 $\check X(B)$ 上诱导良定义的复结构。相对于这一复结构，典范辛形式 $\omega$ 是一个度量的 Kähler 形式。记 $\check y_i=\partial K/\partial y_i$ 与 $g_{ij}=\partial^2K/\partial y_i\partial y_j$。度量 Ricci 平坦，当且仅当 $-\log\det(g_{ij})$ 表示为 $\check y$ 的函数后为仿射线性函数。若采用自然全纯体积形式的 Calabi--Yau 归一化，则条件为实 Monge--Ampère 方程
\[
\det\!\left(\frac{\partial^2K}{\partial y_i\partial y_j}\right)=\text{常数}.
\]
\end{proposition}
\begin{proof}
\SourceParagraph{ch06-6.2.1-015}{442}
与\BookRef{exer:6.12}{练习6.12}中一样，容易看到，坐标 $y_i$ 的仿射线性变换，以及坐标 $\check x_i$ 相应的变换，会使坐标 $z_i$ 作仿射线性变换。因此，它们诱导良定义的复结构；这一结构在 $\check x_i\mapsto\check x_i+1$ 下保持不变，因而在 $\check X(B)$ 上诱导复结构。于是，计算得
\[
\omega=\sum_i\mathrm d\check x_i\wedge\mathrm dy_i
=\frac{\mathrm i}{2}\sum_{i,j}g^{ij}\,\mathrm dz_i\wedge\mathrm d\bar z_j,
\]
其中，$g_{ij}=\partial^2K/\partial y_i\partial y_j$，$(g^{ij})$ 为 $(g_{ij})$ 的逆矩阵。Ricci 形式由 $-\mathrm i\,\partial\bar\partial\log\det(g^{ij})$ 给出，而度量系数只依赖于 $\operatorname{Im}z_i=\check y_i$，这给出所述的一般 Ricci 平坦条件。在自然体积形式的归一化下，要求 $\det(g^{ij})$ 为常数，等价于 $\det(g_{ij})$ 为常数。
\end{proof}

\SourceParagraph{ch06-6.2.1-016}{442}
与前面一样，称这个 Kähler 流形为 $\check X(B,K)$。现在，注意以下事实。

\SourceParagraph{ch06-6.2.1-017}{442}
\BookStatementNumber{16}
\begin{proposition}\label{prop:6.16}
存在 Kähler 流形之间的典范同构
\[
X(B,K)\cong\check X(\check B,\check K),
\]
其中，$(\check B,\check K)$ 为 $(B,K)$ 的 Legendre 变换。
\end{proposition}
\begin{proof}
\SourceParagraph{ch06-6.2.1-018}{442}
当然，作为流形，$B=\check B$，但它们带有不同的仿射结构。此外，由\BookRef{prop:6.4}{命题6.4}，$K$ 与 $\check K$ 诱导的度量 $g$ 相同。现在，利用这个度量，将 $TB$ 与 $T^*B=T^*\check B$ 识别。在局部坐标 $(y_i)$ 中，$\partial/\partial y_i$ 与 $\sum_jg_{ij}\,\mathrm dy_j$ 识别。但是，
\[
\mathrm d\check y_i
=\sum_j\frac{\partial^2K}{\partial y_i\partial y_j}\,\mathrm dy_j
=\sum_jg_{ij}\,\mathrm dy_j,
\]
所以，$\partial/\partial y_i$ 与 $\mathrm d\check y_i$ 识别。因此，这一识别下降为 $X(B)$ 与 $\check X(\check B)$ 的典范识别。

\SourceParagraph{ch06-6.2.1-019}{442}
只需确认，这个识别给出 Kähler 流形的同构。但是，在这一识别下，$T^*\check B$ 上的复坐标
\[
\check z_i=\check x_i+\mathrm i\frac{\partial\check K}{\partial\check y_i}
\]
与 $TB$ 上的坐标 $z_i=x_i+\mathrm i y_i$ 识别，所以复结构相同。最后，Kähler 形式分别为
\[
\frac{\mathrm i}{2}\sum_{i,j}g_{ij}\,\mathrm dz_i\wedge\mathrm d\bar z_j,
\qquad
\frac{\mathrm i}{2}\sum_{i,j}\check g^{ij}\,\mathrm dz_i\wedge\mathrm d\bar z_j,
\]
其中，$\check g_{ij}=g(\partial/\partial\check y_i,\partial/\partial\check y_j)$。但是，
\[
\begin{aligned}
\delta_{ik}
&=\sum_jg_{ij}g^{jk}
=\sum_jg^{jk}g\!\left(\frac{\partial}{\partial y_i},\frac{\partial}{\partial y_j}\right)\\
&=\sum_jg^{jk}g\!\left(
\sum_a g_{ia}\frac{\partial}{\partial\check y_a},
\sum_l g_{jl}\frac{\partial}{\partial\check y_l}\right)\\
&=\sum_{j,a,l}g^{jk}g_{ia}g_{jl}\check g_{al}\\
&=\sum_a g_{ia}\check g_{ak}.
\end{aligned}
\]
所以，$\check g^{ij}=g_{ij}$。因此，两个 Kähler 形式也相同。
\end{proof}

\SourceParagraph{ch06-6.2.1-020}{443}
现在，可以更明确地陈述镜像对称最简单的形式。

\SourceParagraph{ch06-6.2.1-021}{443}
\BookStatementNumber{17}
\begin{definition}\label{def:6.17}
若 $B$ 为仿射流形，转移函数属于 $\mathbb R^n\rtimes\operatorname{GL}(n,\mathbb Z)$，则称 $X(B)$ 与 $\check X(B)$ 互为 SYZ 对偶。若进一步有凸函数 $K\colon\widetilde B\to\mathbb R$，其 Hessian 矩阵正定并定义 $B$ 上的度量，则称
\[
X(B,K)\cong\check X(\check B,\check K)
\]
与
\[
\check X(B,K)\cong X(\check B,\check K)
\]
互为 SYZ 对偶。
\end{definition}

\SourceParagraph{ch06-6.2.1-022}{443}
前一种情形，是复流形与辛流形之间的对偶；后一种情形，是 Kähler 流形之间的对偶。将两种情形都看作镜像对称的简单形式。

\subsection{半平坦微分形式}\label{sec:6.2.2}

\SourceParagraph{ch06-6.2.2-001}{443}
本小节将讨论 $X(B)$ 与 $\check X(B)$ 上的形式，以及它们之间的关系。因此，采用 $B$ 上的局部仿射坐标 $y_1,\ldots,y_n$，以及切丛上相应的坐标 $x_1,\ldots,x_n$ 和余切丛上的坐标 $\check x_1,\ldots,\check x_n$，会很有用。此外，本小节假定 $B$ 的转移映射属于 $\mathbb R^n\rtimes\operatorname{SL}_n(\mathbb Z)$，从而 $X(B)$ 带有处处非零的全纯 $n$ 形式。

\SourceParagraph{ch06-6.2.2-002}{443}
\BookStatementNumber{18}
\begin{definition}\label{def:6.18}
$X(B)$ 上的 $(p,q)$ 型半平坦微分形式，是一个 $(p+q)$ 形式，在 $B$ 上局部地写为
\[
\sum_{\substack{\#I=p\\\#J=q}}\alpha_{IJ}(y)\,\mathrm dy_I\wedge\mathrm dx_J.
\]
相应地，$\check X(B)$ 上的 $(p,q)$ 型半平坦微分形式，局部地写为
\[
\sum_{\substack{\#I=p\\\#J=q}}\alpha_{IJ}(y)\,\mathrm dy_I\wedge\mathrm d\check x_J.
\]
这里，$I,J$ 为多重指标集，$\alpha_{IJ}$ 为 $B$ 上的函数，$\#I,\#J$ 分别为 $I,J$ 的元素个数。注意，这里的类型不依赖于仿射坐标的选择。用 $S^{p,q}$ 与 $\check S^{p,q}$，分别表示 $X(B)$ 与 $\check X(B)$ 上的 $(p,q)$ 型半平坦形式空间。
\end{definition}

\SourceParagraph{ch06-6.2.2-003}{444}
\begin{example*}
$\check X(B)$ 上的辛形式 $\omega$，是 $(1,1)$ 型半平坦形式。
\end{example*}

\SourceParagraph{ch06-6.2.2-004}{444}
关键的一点是，存在自然同构
\[
\mu\colon S^{p,q}\xrightarrow{\ \cong\ }\check S^{p,n-q},
\]
现在来定义它。首先，只考虑一个环面 $T^n\cong V/\Lambda$，其中，$V$ 为向量空间，$\Lambda$ 为格。那么，$H^p(V/\Lambda,\mathbb R)$ 自然地与环面上的常系数 $p$ 形式空间，即 $\bigwedge^pV^\vee$，识别。另一方面，对偶环面为 $V^\vee/\Lambda^\vee$，其中，$\Lambda^\vee\subseteq V^\vee$ 是在 $\Lambda$ 上取整数值的线性泛函所构成的格。于是，$H^p(V^\vee/\Lambda^\vee,\mathbb R)\cong\bigwedge^pV$。此外，楔积给出自然配对
\[
\bigwedge^pV^\vee\times\bigwedge^{n-p}V^\vee
\longrightarrow\bigwedge^nV^\vee.
\]
因此，选择同构 $\bigwedge^nV^\vee\cong\mathbb R$，这等价于选择一个非零 $n$ 形式，就得到 $\bigwedge^{n-p}V^\vee$ 与 $\bigwedge^pV$ 的识别。具体地，给定 $V$ 上的坐标 $x_1,\ldots,x_n$，可以把 $\bigwedge^pV$ 的一个元素写为
\[
\sum_{\#J=p}a_J\frac{\partial}{\partial x_J},
\]
其中，$J=\{j_1,\ldots,j_p\}$ 为多重指标集，且
\[
\frac{\partial}{\partial x_J}
=\frac{\partial}{\partial x_{j_1}}\wedge\cdots\wedge\frac{\partial}{\partial x_{j_p}}.
\]
若所选的 $n$ 形式为 $\mathrm dx_1\wedge\cdots\wedge\mathrm dx_n$，则相应的 $\bigwedge^{n-p}V^\vee$ 中的元素，可以写为内乘
\[
\iota\!\left(\sum_Ja_J\frac{\partial}{\partial x_J}\right)
(\mathrm dx_1\wedge\cdots\wedge\mathrm dx_n).
\]
这里采用如下约定：多重向量 $v_1\wedge\cdots\wedge v_p$ 与 $q$ 形式 $\alpha$ 的内乘，是 $(q-p)$ 形式 $\iota(v_1\wedge\cdots\wedge v_p)\alpha$，在向量 $w_1,\ldots,w_{q-p}$ 上取值为
\[
\alpha(v_1,\ldots,v_p,w_1,\ldots,w_{q-p}).
\]
这主要是为了明确符号约定；虽然繁琐，但遗憾的是，它非常重要。

\SourceParagraph{ch06-6.2.2-005}{444}
现在来看整体情形。由于对 $B$ 的转移映射所作的假设，流形 $X(B)$ 带有整体全纯 $n$ 形式 $\Omega$。将它乘以固定的归一化常数后，局部地有
\[
\Omega=\bigwedge_{i=1}^n(\mathrm dx_i+\mathrm i\,\mathrm dy_i).
\]
因此，在 $X(B)$ 上，可以把 $\Omega$ 限制到任意纤维，得到一个典范选择的 $n$ 形式 $\mathrm dx_1\wedge\cdots\wedge\mathrm dx_n$。然后，局部地定义同构 $\mu\colon S^{p,q}\to\check S^{p,n-q}$ 为
\[
\begin{aligned}
&\mu\!\left(\sum_{I,J}a_{IJ}(y)\,\mathrm dy_I\wedge\mathrm dx_J\right)\\
&\quad=(-1)^{p+q(q-1)/2+n(n-1)/2}
\sum_{I,J}a_{IJ}(y)\,\mathrm dy_I\wedge
\iota\!\left(\frac{\partial}{\partial\check x_J}\right)
(\mathrm d\check x_1\wedge\cdots\wedge\mathrm d\check x_n).
\end{aligned}
\]

\SourceParagraph{ch06-6.2.2-006}{445}
\BookStatementNumber{19}
\begin{exercise}\label{exer:6.19}
\begin{enumerate}[label=(\arabic*)]
\item 证明 $\mu$ 良定义，不依赖于仿射坐标的选择。
\item 证明 $\mu$ 把闭形式映到闭形式。
\end{enumerate}
\end{exercise}

\SourceParagraph{ch06-6.2.2-007}{445}
最后，讨论局部系统的 de Rham 上同调。

\SourceParagraph{ch06-6.2.2-008}{445}
令 $E$ 为流形 $M$ 上的局部系统，纤维为 $\mathbb R^n$。令 $\mathcal E=E\otimes_{\mathbb R}C^\infty(M)$；这里，$C^\infty(M)$ 表示 $M$ 上的 $C^\infty$ 函数层。那么，$\mathcal E$ 是 $M$ 上某个向量丛的截面层。局部地，若 $U\subseteq M$ 为开集，$E|_U$ 平凡，并有截面基 $s_1,\ldots,s_n\in\Gamma(U,E)$，则
\[
\Gamma(U,\mathcal E)=\left\{\sum_i f_i s_i\;\middle|\;
f_i\colon U\to\mathbb R\text{ 为 }C^\infty\text{ 函数}\right\}.
\]
$\mathcal E$ 有由 $E$ 决定的联络 $\nabla\colon\mathcal E\to T^*M\otimes\mathcal E$。局部地，它的形式为
\[
\nabla\!\left(\sum_i f_i s_i\right)=\sum_i\mathrm df_i\otimes s_i.
\]
这个联络平坦，因为 $\nabla s_i=0$，$i=1,\ldots,n$。因此，得到映射序列
\[
\mathcal E\xrightarrow{\nabla}T^*M\otimes\mathcal E
\xrightarrow{\nabla}\bigwedge^2T^*M\otimes\mathcal E
\xrightarrow{\nabla}\cdots.
\]
由于联络平坦，$\nabla^2=0$；换言之，$\nabla$ 定义一个复形。此外，
\[
\ker\bigl(\mathcal E\xrightarrow{\nabla}T^*M\otimes\mathcal E\bigr)=E,
\]
所以，这个复形是 $E$ 的解消。它推广了通常的 de Rham 复形
\[
\bigwedge^0T^*M\xrightarrow{\mathrm d}\bigwedge^1T^*M
\xrightarrow{\mathrm d}\cdots,
\]
后者是常值层 $\mathbb R$ 的解消。此外，每个层 $\bigwedge^pT^*M\otimes\mathcal E$，都是所谓的细层，参见例如文献\BookCite{ref094}。因此，对所有 $q>0$，有
\[
H^q\!\left(M,\bigwedge^pT^*M\otimes\mathcal E\right)=0.
\]
由标准同调论证，得到
\[
H^p(M,E)=
\frac{\ker\!\left(\nabla\colon\Gamma(M,\bigwedge^pT^*M\otimes\mathcal E)
\to\Gamma(M,\bigwedge^{p+1}T^*M\otimes\mathcal E)\right)}
{\operatorname{im}\!\left(\nabla\colon\Gamma(M,\bigwedge^{p-1}T^*M\otimes\mathcal E)
\to\Gamma(M,\bigwedge^pT^*M\otimes\mathcal E)\right)}.
\]
这就是 $H^p(M,E)$ 的 de Rham 实现。也可以用 Čech 上同调计算这个上同调群，见\BookSectionRef{sec:4.5.1}{4.5.1}，这也将很有用。

\SourceParagraph{ch06-6.2.2-009}{445}
现在，考虑仿射流形 $B$ 上的局部系统 $\bigwedge^q\check\Lambda_{\mathbb R}$。给定纤维 $\bigwedge^q\check\Lambda_{\mathbb R,b}$，可以将其看作环面 $\Lambda_{\mathbb R,b}/\Lambda_b$ 上的常系数 $q$ 形式空间。在这一识别下，若仿射坐标图 $U$ 上的仿射坐标为 $y_1,\ldots,y_n$，则局部截面基可以识别为
\[
\{\mathrm dx_{i_1}\wedge\cdots\wedge\mathrm dx_{i_q}
\mid1\leq i_1<\cdots<i_q\leq n\}.
\]
令 $\mathcal E_q=\bigwedge^q\check\Lambda_{\mathbb R}\otimes_{\mathbb R}C^\infty(B)$。于是看到，$f^{-1}(U)$ 上的 $(p,q)$ 型半平坦微分形式空间，可以与 $\bigwedge^pT^*B\otimes\mathcal E_q$ 在 $U$ 上的截面空间识别，其中，将 $\sum_{I,J}\alpha_{IJ}\,\mathrm dy_I\wedge\mathrm dx_J$ 与 $\sum_{I,J}\alpha_{IJ}\,\mathrm dy_I\otimes\mathrm dx_J$ 识别。此外，联络
\[
\nabla\colon\bigwedge^pT^*B\otimes\mathcal E_q
\longrightarrow\bigwedge^{p+1}T^*B\otimes\mathcal E_q
\]
与微分 $\mathrm d\colon S^{p,q}\to S^{p+1,q}$ 相同。因此，已经证明以下命题。

\SourceParagraph{ch06-6.2.2-010}{446}
\BookStatementNumber{20}
\begin{proposition}\label{prop:6.20}
存在自然映射
\[
\{S^{p,q}\text{ 中的闭形式}\}
\longrightarrow H^p\!\left(B,\bigwedge^q\check\Lambda_{\mathbb R}\right),
\]
以及对偶的自然映射
\[
\{\check S^{p,q}\text{ 中的闭形式}\}
\longrightarrow H^p\!\left(B,\bigwedge^q\Lambda_{\mathbb R}\right).
\]
对偶映射采用整体符号 $(-1)^p$：将 $\mathrm dy_I\wedge\mathrm d\check x_J$ 对应于 $(-1)^p\mathrm dy_I\otimes\partial/\partial y_J$。特别地，$\sum_{i,j}b_{ij}\,\mathrm d\check x_j\wedge\mathrm dy_i$ 对应于 $\sum_{i,j}b_{ij}\,\mathrm dy_i\otimes\partial/\partial y_j$。
\end{proposition}

\SourceParagraph{ch06-6.2.2-011}{446}
还有必要说明，给定 $H^p(B,\bigwedge^q\check\Lambda_{\mathbb R})$ 中某元素的 de Rham 代表，如何关联一个 Čech 代表；对 $H^p(B,\bigwedge^q\Lambda_{\mathbb R})$ 的对偶构造完全相同。为此，从一个 $(p,q)$ 型闭半平坦形式 $\alpha$ 出发。进一步，选择 $B$ 的开覆盖 $\{U_i\}_{i\in I}$，使得对于所有指标，其非空有限交集都可缩。因此，在每个 $U_i$ 上，存在半平坦 $(p-1,q)$ 形式 $\alpha_i^1$，使得在 $f^{-1}(U_i)$ 上有 $\mathrm d\alpha_i^1=\alpha$。这里先取 $p>0$。事实上，在 $U_i$ 上，
\[
\alpha=\sum_{I,J}\alpha_{IJ}(y)\,\mathrm dy_I\wedge\mathrm dx_J,
\]
而 $\mathrm d\alpha=0$ 蕴含，对于每个 $J$，有 $\mathrm d(\sum_I\alpha_{IJ}\,\mathrm dy_I)=0$。由 $U_i$ 可缩，$\sum_I\alpha_{IJ}\,\mathrm dy_I$ 为恰当形式。现在，在 $U_i\cap U_j$ 上，$\alpha_i^1$ 与 $\alpha_j^1$ 未必相同，但 $\mathrm d(\alpha_j^1-\alpha_i^1)=0$。令 $\beta_{ij}^1=\alpha_j^1-\alpha_i^1$ 于 $U_i\cap U_j$ 上。若 $p\geq2$，由同一论证，可以找到半平坦 $(p-2,q)$ 形式 $\alpha_{ij}^2$，使得 $\mathrm d\alpha_{ij}^2=\alpha_j^1-\alpha_i^1$。于是，对于 $i_0,i_1,i_2$，在 $U_{i_0}\cap U_{i_1}\cap U_{i_2}$ 上有
\[
\begin{aligned}
\mathrm d(\alpha_{i_0i_1}^2+\alpha_{i_1i_2}^2-\alpha_{i_0i_2}^2)
&=\alpha_{i_1}^1-\alpha_{i_0}^1+\alpha_{i_2}^1-\alpha_{i_1}^1
-\alpha_{i_2}^1+\alpha_{i_0}^1\\
&=0.
\end{aligned}
\]
令 $\beta_{i_0i_1i_2}^2=\alpha_{i_0i_1}^2+\alpha_{i_1i_2}^2-\alpha_{i_0i_2}^2$。归纳地继续这一过程，直到在 $U_{i_0}\cap\cdots\cap U_{i_p}$ 上得到闭半平坦 $(0,q)$ 形式 $\beta_{i_0\cdots i_p}^p$。由于这些形式为闭形式，它们实际上是 $\bigwedge^q\check\Lambda_{\mathbb R}$ 的平坦截面，并定义这个局部系统的 Čech $p$ 上闭链。对于 $p=0$，原来的闭半平坦形式本身就是这样的平坦截面。

\SourceParagraph{ch06-6.2.2-012}{446}
\BookStatementNumber{21}
\begin{exercise}\label{exer:6.21}
按如下方式写出这一归纳过程的细节。令
\[
\begin{aligned}
Z^{p,q}
&=\ker\!\left(\nabla\colon\bigwedge^pT^*B\otimes\mathcal E_q
\longrightarrow\bigwedge^{p+1}T^*B\otimes\mathcal E_q\right)\\
&=\operatorname{im}\!\left(\nabla\colon\bigwedge^{p-1}T^*B\otimes\mathcal E_q
\longrightarrow\bigwedge^pT^*B\otimes\mathcal E_q\right),
\end{aligned}
\]
其中，后一等式对 $p>0$ 成立，右侧均指相应的光滑截面层。这给出层的正合列
\[
0\longrightarrow Z^{p,q}\longrightarrow\bigwedge^pT^*B\otimes\mathcal E_q
\xrightarrow{\nabla}Z^{p+1,q}\longrightarrow0,
\]
并且 $Z^{0,q}=\bigwedge^q\check\Lambda_{\mathbb R}$。从闭半平坦 $(p_0,q)$ 形式 $\alpha$，即 $Z^{p_0,q}$ 的截面出发，利用 Čech 上边界映射，考察这一正合列在 $p=p_0-1,p_0-2,\ldots$ 时给出的连接同态的像的 Čech 代表：
\[
\begin{aligned}
\alpha\in H^0(B,Z^{p_0,q})
&\longrightarrow H^1(B,Z^{p_0-1,q})\longrightarrow\cdots\\
&\longrightarrow H^{p_0}(B,Z^{0,q})
=H^{p_0}\!\left(B,\bigwedge^q\check\Lambda_{\mathbb R}\right).
\end{aligned}
\]
描述这些 Čech 代表，并证明它们与前面所描述的 $\beta^i$ 相同。
\end{exercise}

\subsection{带 B 场的 SYZ}\label{sec:6.2.3}

\SourceParagraph{ch06-6.2.3-001}{447}
不考虑 B 场，镜像对称就不完整。传统上，B 场应当取值于 $H^2(X,\mathbb R/\mathbb Z)$。不过，暂时将在一个较小的空间中考虑 B 场，而把一般情形留到\BookSectionRef{sec:6.2.5}{6.2.5}讨论。

\SourceParagraph{ch06-6.2.3-002}{447}
与通常一样，令 $B$ 为仿射流形，转移函数属于 $\mathbb R^n\rtimes\operatorname{GL}(n,\mathbb Z)$。$B$ 上有局部系统 $\Lambda$，以及 $\Lambda_{\mathbb R}=\Lambda\otimes_{\mathbb Z}\mathbb R$。局部系统 $\Lambda_{\mathbb R}$，就是 $TB$ 中相对于 $B$ 上诱导的平坦联络为平坦的截面所构成的层。

\SourceParagraph{ch06-6.2.3-003}{447}
考虑 $f\colon X(B)\to B$。若 $U\subseteq B$ 为开集，则 $\Gamma(U,\Lambda_{\mathbb R}/\Lambda)$ 全纯地作用于 $f^{-1}(U)$。事实上，若 $s\in\Gamma(U,\Lambda_{\mathbb R}/\Lambda)$，则 $s$ 通过沿纤维平移作用于 $X(U)=TU/\Lambda|_U$，在纤维 $T_bU/\Lambda_b$ 上的作用为 $x\mapsto x+s(b)$。这一作用是全纯的：对于 $\Lambda_{\mathbb R}/\Lambda$ 的截面 $\sum_i\theta_i\partial/\partial y_i$，其中 $\theta_1,\ldots,\theta_n$ 局部为常数，在全纯坐标 $q_i$ 中，作用为
\[
(q_1,\ldots,q_n)\longmapsto
(e^{2\pi\mathrm i\theta_1}q_1,\ldots,e^{2\pi\mathrm i\theta_n}q_n).
\]
此外，若 $K$ 是 $B$ 上一个度量的势函数，则这一作用保持所诱导的 Kähler 度量。

\SourceParagraph{ch06-6.2.3-004}{447}
因此，给定 $\Lambda_{\mathbb R}/\Lambda$ 的 Čech 1上闭链 $(b_{ij})$，即对于 $B$ 的某个开覆盖 $\{U_i\}$，有 $b_{ij}\in\Gamma(U_i\cap U_j,\Lambda_{\mathbb R}/\Lambda)$，且在 $U_i\cap U_j\cap U_k$ 上满足 $b_{ij}+b_{jk}=b_{ik}$，就可以如下构造新的复流形 $X(B,\mathbf B)$，或新的 Kähler 流形 $X(B,\mathbf B,K)$。沿 $X(U_i\cap U_j)$，通过识别
\[
X(U_i)\supseteq X(U_i\cap U_j)
\xrightarrow{\ b_{ij}\ }X(U_i\cap U_j)\subseteq X(U_j),
\]
把流形 $X(U_i)$，或 $X(U_i,K)$，粘合起来；其中，$b_{ij}$ 通过平移作用。Čech 上闭链条件保证这些粘合相容，因而产生复流形 $X(B,\mathbf B)$；若有势函数 $K$，则产生 Kähler 流形 $X(B,\mathbf B,K)$。

\SourceParagraph{ch06-6.2.3-005}{447}
因此，完整的半平坦 SYZ 图景如下。数据由仿射流形 $B$、势 $K$ 和 B 场
\[
\mathbf B\in H^1(B,\Lambda_{\mathbb R}/\Lambda),
\qquad
\check{\mathbf B}\in H^1(B,\check\Lambda_{\mathbb R}/\check\Lambda)
\]
组成。现在，正如\BookRef{prop:6.16}{命题6.16}的证明所说明的，利用 $B$ 的仿射结构定义的局部系统 $\check\Lambda$，与利用 $\check B$ 的仿射结构定义的局部系统 $\Lambda$ 相同。因此，称二元组
\[
\bigl(X(B,\mathbf B,K),\check{\mathbf B}\bigr)
\]
与
\[
\bigl(X(\check B,\check{\mathbf B},\check K),\mathbf B\bigr)
\]
互为 SYZ 对偶。

\SourceParagraph{ch06-6.2.3-006}{448}
这一表述，可能看起来不像 B 场给出复化 Kähler 类的传统形式。因此，稍后将给出这一构造的另一种形式，使这种关系变得清楚。

\SourceParagraph{ch06-6.2.3-007}{448}
\BookStatementNumber{22}
\begin{example}\label{ex:6.22}
现在，在椭圆曲线的情形中考察所有这些概念。令 $B=\mathbb R/\tau_2\mathbb Z$，其中，$\tau_2>0$ 为实数；$B$ 上的仿射结构，由 $\mathbb R$ 上的自然仿射结构诱导。那么，$X(B)=\mathbb R/\tau_2\mathbb Z\times\mathbb R/\mathbb Z$。用 $y$ 表示 $\mathbb R/\tau_2\mathbb Z$ 上的坐标，$x$ 表示 $\mathbb R/\mathbb Z$ 上的坐标，其复结构由复坐标 $z=x+\mathrm i y$ 给出。等价地，$X(B)$ 是椭圆曲线 $\mathbb C/\langle1,\mathrm i\tau_2\rangle$，其中，$\langle1,\mathrm i\tau_2\rangle$ 表示由这两个数生成的 $\mathbb C$ 中的格。

\SourceParagraph{ch06-6.2.3-008}{448}
要计算 $H^1(B,\Lambda_{\mathbb R}/\Lambda)$，注意，$\Lambda$ 为常值层 $\mathbb Z$，所以
\[
H^1(B,\Lambda_{\mathbb R}/\Lambda)
=H^1(S^1,\mathbb R/\mathbb Z)=\mathbb R/\mathbb Z.
\]
将如下结论留给读者作为练习：若 $\mathbf B\in H^1(B,\mathbb R/\mathbb Z)$ 由 $\tau_1\in\mathbb R$ 代表，则 $X(B,\mathbf B)$ 同构于椭圆曲线 $\mathbb C/\langle1,\tau\rangle$，其中，$\tau=\tau_1+\mathrm i\tau_2$。当然，将 $\tau_1$ 换成 $\tau_1+1$，对格没有影响。

\SourceParagraph{ch06-6.2.3-009}{448}
另一方面，$\check X(B)=\mathbb R/\tau_2\mathbb Z\times\mathbb R/\mathbb Z$ 是辛流形，辛形式为 $\mathrm d\check x\wedge\mathrm dy$，体积为 $\tau_2$。

\SourceParagraph{ch06-6.2.3-010}{448}
若允许使用势，可以取 $K$ 为 $y$ 的任意严格凸二次函数。由于度量只在相差一个仿射线性函数的意义下确定 $K$，可以假定
\[
K=\frac{\check\tau_2}{\tau_2}\frac{y^2}{2},
\qquad
\check B=\mathbb R/\check\tau_2\mathbb Z,
\qquad
\check K=\frac{\tau_2}{\check\tau_2}\frac{\check y^2}{2}.
\]
那么，在 $\mathbb C/\langle1,\tau_1+\mathrm i\tau_2\rangle$ 上使用复坐标 $z$，由 $K\circ f$ 定义的 Kähler 度量，其 Kähler 形式为
\[
\omega=\frac{\mathrm i}{2}\frac{\check\tau_2}{\tau_2}\,
\mathrm dz\wedge\mathrm d\bar z,
\]
体积为 $\check\tau_2$。再取 B 场 $\mathbf B=\tau_1$ 与 $\check{\mathbf B}=\check\tau_1$，则 $(X(B,\tau_1,K),\check\tau_1)$ 与 $(X(\check B,\check\tau_1,\check K),\tau_1)$ 互为 SYZ 对偶。这就是说，复化 Kähler 类为 $\check\tau_1+\mathrm i\check\tau_2$ 的椭圆曲线 $\mathbb C/\langle1,\tau_1+\mathrm i\tau_2\rangle$，与复化 Kähler 类为 $\tau_1+\mathrm i\tau_2$ 的椭圆曲线 $\mathbb C/\langle1,\check\tau_1+\mathrm i\check\tau_2\rangle$ 互为镜像。
\end{example}

\SourceParagraph{ch06-6.2.3-011}{448}
\BookStatementNumber{23}
\begin{exercise}\label{exer:6.23}
在\BookRef{ex:6.7}{例题6.7}(2)中，考虑由 $T_1,T_2$ 生成的格 $\Gamma\subseteq G$，其中
\[
\begin{aligned}
T_1(m_1,m_2)&=(m_1+1,m_2),\\
T_2(m_1,m_2)&=(m_1+em_2+e(e-1)/2,\;m_2+e),
\end{aligned}
\]
且 $e>0$ 为整数。令 $B=\mathbb R^2/\Gamma$。证明
\[
H^1(B,\Lambda_{\mathbb R}/\Lambda)
=(\mathbb R/\mathbb Z)^2\oplus(\mathbb Z/e\mathbb Z).
\]
提示：$\Lambda$ 为局部系统。在选取对应环路方向后，它沿 $H_1(B,\mathbb Z)$ 中对应于 $T_1,T_2$ 的生成元的单值变换，分别为
\[
\begin{pmatrix}1&0\\0&1\end{pmatrix},
\qquad
\begin{pmatrix}1&e\\0&1\end{pmatrix}.
\]
局部系统 $\Lambda_{\mathbb R}/\Lambda$ 的上同调，与群上同调 $H^1(\mathbb Z^2,(\mathbb R/\mathbb Z)^2)$ 相同；其中，$\mathbb Z^2=H_1(B,\mathbb Z)$ 通过上述矩阵作用。然后，用标准方法计算这一群上同调，参见例如文献\BookCite{ref068}。
\end{exercise}

\SourceParagraph{ch06-6.2.3-012}{449}
还有一种不同的方式，可以更直接地用微分形式理解 B 场；\BookSectionRef{sec:6.2.5}{6.2.5}将把这一方法推广到广义复结构。首先，回顾 Hitchin\BookCite{ref232}的一些观察。

\SourceParagraph{ch06-6.2.3-013}{449}
\BookStatementNumber{24}
\begin{proposition}\label{prop:6.24}
令 $X$ 为实 $2n$ 维流形。若 $\Omega$ 是 $X$ 上的复值 $C^\infty$ $n$ 形式，满足以下三个性质：
\begin{enumerate}[label=(\arabic*)]
\item $\mathrm d\Omega=0$；
\item $\Omega$ 局部可分解，即局部地可以写为 $\theta_1\wedge\cdots\wedge\theta_n$，其中，$\theta_1,\ldots,\theta_n$ 为1形式；
\item 在 $X$ 上处处有 $(-1)^{n(n-1)/2}(\mathrm i/2)^n\Omega\wedge\overline\Omega>0$；
\end{enumerate}
则 $\Omega$ 决定 $X$ 上的一个复结构，相对于这一复结构，$\Omega$ 为全纯 $n$ 形式。
\end{proposition}
\begin{proof}
\SourceParagraph{ch06-6.2.3-014}{449}
$\Omega$ 在每点定义 $T^*X\otimes\mathbb C$ 的一个 $n$ 维子空间 $\Omega_X^{1,0}$：取分解 $\Omega=\theta_1\wedge\cdots\wedge\theta_n$，再取 $T^*X\otimes\mathbb C$ 中由 $\theta_1,\ldots,\theta_n$ 张成的子空间。外代数中的一个标准事实说明，这样得到的子空间良定义，实际上不依赖于具体的分解选择。条件(3)蕴含，这个子空间与其复共轭互补，从而得到分裂
\[
T^*X\otimes_{\mathbb R}\mathbb C=\Omega_X^{1,0}\oplus\Omega_X^{0,1}.
\]
于是，可以把这两个子空间分别作为 $J$ 在余切空间上的 $\pm\mathrm i$ 特征空间，定义一个概复结构 $J$。

\SourceParagraph{ch06-6.2.3-015}{449}
要证明 $J$ 为可积复结构，利用第一个条件。写成 $\Omega=\theta_1\wedge\cdots\wedge\theta_n$，并注意 $0=\Omega\wedge\theta_i$，所以
\[
\begin{aligned}
0=\mathrm d(\Omega\wedge\theta_i)
&=\mathrm d\Omega\wedge\theta_i+(-1)^n\Omega\wedge\mathrm d\theta_i\\
&=(-1)^n\Omega\wedge\mathrm d\theta_i.
\end{aligned}
\]
现在，回顾 Newlander--Nirenberg 定理。若 $\theta_1,\ldots,\theta_n$ 为关于 $J$ 的 $(1,0)$ 形式的一组基，则可以写为
\[
\mathrm d\theta_i
=\sum_{j<k}A_i^{jk}\theta_j\wedge\theta_k
+\sum_{j,k}B_i^{jk}\theta_j\wedge\bar\theta_k
+\sum_{j<k}C_i^{jk}\bar\theta_j\wedge\bar\theta_k.
\]
Newlander--Nirenberg 定理说，$J$ 可积，当且仅当对所有 $i,j,k$，$C_i^{jk}=0$。在这里，这由 $0=\Omega\wedge\mathrm d\theta_i$ 蕴含。因此，$J$ 可积。
\end{proof}

\SourceParagraph{ch06-6.2.3-016}{449}
为了供后面使用，还指出指定 Ricci 平坦 Kähler 流形所需的数据。

\SourceParagraph{ch06-6.2.3-017}{449}
\BookStatementNumber{25}
\begin{proposition}\label{prop:6.25}
令 $X$ 为实 $2n$ 维流形。设 $\omega$ 为 $X$ 上的辛形式，$\Omega$ 为 $X$ 上的复值 $n$ 形式，并满足：
\begin{enumerate}[label=(\arabic*)]
\item $\Omega$ 满足\BookRef{prop:6.24}{命题6.24}的条件；
\item 在\BookRef{prop:6.24}{命题6.24}的复结构中，$\omega$ 为正的 $(1,1)$ 形式；
\item $(-1)^{n(n-1)/2}(\mathrm i/2)^n\Omega\wedge\overline\Omega=\omega^n/n!$。
\end{enumerate}
则 $\Omega$ 在 $X$ 上诱导复结构，使得 $\omega$ 为 $X$ 上的 Kähler 形式，且其对应度量 Ricci 平坦。
\end{proposition}
\begin{proof}
\SourceParagraph{ch06-6.2.3-018}{450}
结论直接成立：条件(3)使所选全纯 $n$ 形式的长度为常数，因而相应的 Ricci 形式为零。
\end{proof}

\SourceParagraph{ch06-6.2.3-019}{450}
现在，如\BookSectionRef{sec:6.2.2}{6.2.2}中一样，假定 $B$ 的转移映射属于 $\mathbb R^n\rtimes\operatorname{SL}_n(\mathbb Z)$。不再直接把 B 场看作 $H^1(B,\Lambda_{\mathbb R}/\Lambda)$ 中的元素，而将其定义为 $\check X(B)$ 上的一个 $(1,1)$ 型半平坦闭形式 $\mathbf B$。特别地，它代表 $H^1(B,\Lambda_{\mathbb R})$ 中的一个元素，进而代表 $H^1(B,\Lambda_{\mathbb R}/\Lambda)$ 中的元素。注意，未必每个 $H^1(B,\Lambda_{\mathbb R}/\Lambda)$ 的元素，都在 $H^1(B,\Lambda_{\mathbb R})$ 的像中；例如，$H^2(B,\Lambda)$ 可能有挠元。对于不在这一像中的 B 场 $\mathbf B$，$X(B,\mathbf B)$ 甚至在拓扑上也不与 $X(B)$ 相同，下面的构造不适用。

\SourceParagraph{ch06-6.2.3-020}{450}
考虑
\[
e^{\mathbf B+\mathrm i\omega}\in\bigoplus_{p=0}^n\check S^{p,p},
\]
其中
\[
e^{\mathbf B+\mathrm i\omega}
=1+\mathbf B+\mathrm i\omega
+\frac{(\mathbf B+\mathrm i\omega)\wedge(\mathbf B+\mathrm i\omega)}{2}+\cdots.
\]
可以把它形式地理解，也可以将其看作 Hitchin 意义下的广义 Calabi--Yau 结构；\BookSectionRef{sec:6.2.5}{6.2.5}将讨论这一点。

\SourceParagraph{ch06-6.2.3-021}{450}
那么，$\mu^{-1}(e^{\mathbf B+\mathrm i\omega})\in\bigoplus_{p=0}^nS^{p,n-p}$，实际上是一个 $n$ 形式。

\SourceParagraph{ch06-6.2.3-022}{450}
\BookStatementNumber{26}
\begin{proposition}\label{prop:6.26}
$\mu^{-1}(e^{\mathbf B+\mathrm i\omega})$ 是相对于 $X(B)$ 上某个复结构的全纯 $n$ 形式。赋予这一复结构以后，$X(B)$ 同构于 $X(B,\mathbf B)$。这里，利用\BookRef{prop:6.20}{命题6.20}，把 $\mathbf B$ 看作定义 $H^1(B,\Lambda_{\mathbb R}/\Lambda)$ 中元素的形式。
\end{proposition}
\begin{proof}
\SourceParagraph{ch06-6.2.3-023}{450}
需要证明 $\mu^{-1}(e^{\mathbf B+\mathrm i\omega})$ 满足\BookRef{prop:6.24}{命题6.24}的条件。它显然为闭形式，因为 $e^{\mathbf B+\mathrm i\omega}$ 的每个分量都是闭形式，而 $\mu^{-1}$ 把闭形式映到闭形式。另一方面，仔细处理符号约定可知，若
\[
\mathbf B=\sum_{i,j}b_{ij}\,\mathrm d\check x_j\wedge\mathrm dy_i,
\qquad
\omega=\sum_i\mathrm d\check x_i\wedge\mathrm dy_i,
\]
则
\[
\mu^{-1}(e^{\mathbf B+\mathrm i\omega})
=\bigwedge_{i=1}^n\left(\mathrm dx_i
+\sum_{j=1}^n(b_{ji}+\mathrm i\delta_{ji})\,\mathrm dy_j\right).
\]
因此，$\mu^{-1}(e^{\mathbf B+\mathrm i\omega})$ 满足\BookRef{prop:6.24}{命题6.24}的假设，得到环面丛 $X(B)$ 上的一个复结构。需要将其与 $X(B,\mathbf B)$ 上的复结构比较；不过，为此，首先要写出 $\mathbf B$ 的一个 Čech 代表。所幸，前一小节的结尾已经说明了如何做到这一点。选择 $B$ 的开覆盖 $\{U_i\}$。在 $U_i$ 上，使用坐标 $y_j$，写成
\[
\mathbf B=\sum_{j,k}b_{jk}\,\mathrm d\check x_k\wedge\mathrm dy_j.
\]
然后，可以找到函数 $c_i^k$，使得 $\mathrm dc_i^k=\sum_jb_{jk}\,\mathrm dy_j$。这里，$(c_i^1,\ldots,c_i^n)$ 应当看作 $TB$ 的截面 $c_i$，即
\[
c_i=\sum_kc_i^k\frac{\partial}{\partial y_k}.
\]
差 $\beta_{ij}=c_j-c_i$ 为 $TB$ 的平坦截面，即 $U_i\cap U_j$ 上的 $\Lambda_{\mathbb R}$ 截面；在所采用的对偶符号约定下，$(\beta_{ij})$ 为代表 $\mathbf B$ 的 Čech 1上闭链。然后，考虑 $f^{-1}(U_i)$。标准复结构由全纯体积形式
\[
\Omega=\bigwedge_{j=1}^n(\mathrm dx_j+\mathrm i\,\mathrm dy_j)
\]
诱导。

\SourceParagraph{ch06-6.2.3-024}{451}
若 $T_{c_i}\colon f^{-1}(U_i)\to f^{-1}(U_i)$ 表示沿纤维平移 $c_i$，即
\[
T_{c_i}(y_1,\ldots,y_n,x_1,\ldots,x_n)
=(y_1,\ldots,y_n,x_1+c_i^1,\ldots,x_n+c_i^n),
\]
则
\[
\begin{aligned}
T_{c_i}^*\Omega
&=\bigwedge_{j=1}^n(\mathrm dx_j+\mathrm dc_i^j+\mathrm i\,\mathrm dy_j)\\
&=\bigwedge_{j=1}^n\left(\mathrm dx_j
+\sum_{k=1}^n(b_{kj}+\mathrm i\delta_{kj})\,\mathrm dy_k\right).
\end{aligned}
\]
因此，若用 $f\colon X(B)\to B$ 表示带有标准复结构的环面丛，用 $f'\colon X(B)'\to B$ 表示带有 $\mu^{-1}(e^{\mathbf B+\mathrm i\omega})$ 所诱导复结构的环面丛，则存在局部同构 $T_{c_i}\colon(f')^{-1}(U_i)\to f^{-1}(U_i)$。在 $U_i\cap U_j$ 上，得到交换图
\[
\BookImage{02_Unnumbered_Diagrams/Chapter_06/Ch06_Diagram_PDFp451_01_Bookp441.pdf}
\]
所以，得到同构 $X(B)'\cong X(B,(\beta_{ij}))$，正如所要证明的。
\end{proof}

\SourceParagraph{ch06-6.2.3-025}{451}
\BookStatementNumber{27}
\begin{example}\label{ex:6.27}
回到椭圆曲线的例子，令 $B=\mathbb R/\tau_2\mathbb Z$。在 $X(B)$ 上取
\[
\Omega=\mathrm dx+(\mathrm i+\tau_1/\tau_2)\,\mathrm dy,
\]
在 $\check X(B)$ 上取
\[
\mathbf B+\mathrm i\omega
=(\mathrm i+\tau_1/\tau_2)\,\mathrm d\check x\wedge\mathrm dy.
\]
若改用坐标 $y'=y/\tau_2$，使得 $y'$ 的周期为1，则在 $(\mathbb R/\mathbb Z)^2$ 上，它们的形式为
\[
\begin{aligned}
\Omega&=\mathrm dx+(\tau_1+\mathrm i\tau_2)\,\mathrm dy',\\
\mathbf B+\mathrm i\omega&=(\tau_1+\mathrm i\tau_2)\,
\mathrm d\check x\wedge\mathrm dy',
\end{aligned}
\]
并且 $\mu^{-1}(e^{\mathbf B+\mathrm i\omega})=\Omega$。
\end{example}

\subsection{复流形族}\label{sec:6.2.4}

\SourceParagraph{ch06-6.2.4-001}{451}
现在说明，在一些附加假设下，如何从仿射流形 $B$ 得到一整个非平凡的复结构族。这很重要，因为它有助于明确大复结构极限的概念。

\SourceParagraph{ch06-6.2.4-002}{451}
假定 $B$ 为整仿射流形。那么，实际上可以从这些构造中得到更多内容。首先，令 $B'=B\times\mathbb R_{>0}$，并通过指定展开映射
\[
\delta'\colon\widetilde{B'}=\widetilde B\times\mathbb R_{>0}
\longrightarrow M_{\mathbb R}\oplus\mathbb R,
\qquad
\delta'(b,r)=(r\delta(b),r),
\]
赋予 $B'$ 整仿射结构。注意，若 $\gamma\in\pi_1(B)=\pi_1(B')$，且 $\rho(\gamma)\in\operatorname{Aff}(M)$ 满足 $\rho(\gamma)\circ\delta\circ\Psi_\gamma=\delta$，则可以写成 $\rho(\gamma)(m)=Am+b$，其中，$A\in\operatorname{GL}(M)$、$b\in M$。令 $\Psi_\gamma'$ 为由 $\gamma$ 诱导的 $\widetilde{B'}$ 的覆叠变换，令 $\rho'(\gamma)(m,r)=(Am+rb,r)$。那么，
\[
\begin{aligned}
\rho'(\gamma)\circ\delta'\circ\Psi_\gamma'(b,r)
&=\rho'(\gamma)\circ\delta'(\Psi_\gamma(b),r)\\
&=\rho'(\gamma)(r\delta\Psi_\gamma(b),r)\\
&=(rA\delta\Psi_\gamma(b)+rb,r)\\
&=(r\rho(\gamma)\delta\Psi_\gamma(b),r)\\
&=(r\delta(b),r)\\
&=\delta'(b,r).
\end{aligned}
\]
因此，$\delta'$ 的完整表示由 $\rho'\colon\pi_1(B)\to\operatorname{GL}(M\oplus\mathbb Z)$ 给出，因为 $\rho'(\gamma)$ 始终为线性变换。若 $b$ 不属于 $M$，而只属于 $M_{\mathbb R}$，则这个构造不再保证 $B'$ 具有所需整结构，因而不能直接讨论这里的 $X(B')$。不过，在现在的整性假设下，可以得到复 $(n+1)$ 维流形 $X(B')$，其中，$n=\dim M_{\mathbb R}$。还有 $X(\mathbb R_{>0})$，其中，包含 $\mathbb R_{>0}\hookrightarrow\mathbb R$ 给出底空间的仿射结构。坐标
\[
q=\exp\bigl(2\pi\mathrm i(x+\mathrm i y)\bigr)
\]
将 $X(\mathbb R_{>0})$ 与穿孔单位圆盘 $D^*\subseteq\mathbb C$ 识别。同时，投影 $B'\to\mathbb R_{>0}$ 在切丛之间诱导映射。由于这一投影为仿射线性映射，它诱导映射 $f\colon X(B')\to D^*$。

\SourceParagraph{ch06-6.2.4-003}{452}
可以把这个穿孔圆盘上的族，看作一个大复结构极限退化。当然，缺少的是0点上的纤维；理想情况下，希望在0点上有一个退化的 Calabi--Yau 流形。不过，这是另一回事，见文献\BookCite{ref205,ref206,ref207,ref208}。

\SourceParagraph{ch06-6.2.4-004}{452}
\BookStatementNumber{28}
\begin{exercise}\label{exer:6.28}
证明：若 $B=\mathbb R/n\mathbb Z$，其中，$n$ 为正整数，则 $f\colon X(B')\to D^*$ 是穿孔圆盘 $D^*$ 上的椭圆曲线族，并且
\[
f^{-1}(q)\cong\mathbb C/\left\langle1,\frac{n}{2\pi\mathrm i}\log q\right\rangle.
\]
\end{exercise}

\SourceParagraph{ch06-6.2.4-005}{452}
可以如下推广这一构造。暂时不假定 $B$ 已有仿射结构。设给定由连续函数 $\widetilde B\to M_{\mathbb R}$ 构成的实向量空间 $V_{\mathbb R}$，满足如下性质：存在表示 $\widetilde\rho\colon\pi_1(B)\to\operatorname{GL}(M)$，使得对每个 $\delta\in V_{\mathbb R}$，都存在表示 $\rho_\delta\colon\pi_1(B)\to\operatorname{Aff}(M_{\mathbb R})$，其与投影 $\operatorname{Aff}(M_{\mathbb R})\to\operatorname{GL}(M_{\mathbb R})$ 的复合为 $\widetilde\rho$，并且
\[
\rho_\delta(\gamma)\circ\delta\circ\Psi_\gamma=\delta.
\]
进一步假定，$V_{\mathbb R}$ 包含格 $V$，使得只要 $\delta\in V$，$\rho_\delta$ 的像就属于 $\operatorname{Aff}(M)$。最后，令 $\mathcal V\subseteq V_{\mathbb R}$ 为子集
\[
\mathcal V:=\{\delta\in V_{\mathbb R}\mid\delta\text{ 为浸入}\},
\]
并假定它为非空开集。每个 $\delta\in\mathcal V$，都决定 $B$ 上的一个仿射结构。此外，$\mathcal V$ 从 $V_{\mathbb R}$ 继承仿射结构。如果把 $V$ 看作给出 $V_{\mathbb R}$ 的整结构，就可以讨论 $X(\mathcal V)$。

\SourceParagraph{ch06-6.2.4-006}{453}
现在，利用覆盖上的展开映射
\[
\delta'\colon\widetilde B\times\mathcal V
\longrightarrow M_{\mathbb R}\times V_{\mathbb R},
\qquad
\delta'(b,\delta)=(\delta(b),\delta),
\]
赋予 $B'=B\times\mathcal V$ 仿射结构。可以像前面一样计算 $\delta'$ 的完整表示。对 $\gamma\in\pi_1(B)$，写成
\[
\rho_\delta(\gamma)(m)
=\widetilde\rho(\gamma)(m)+\operatorname{Trans}_\gamma(\delta),
\]
其中，$\operatorname{Trans}_\gamma(\delta)$ 为 $\rho_\delta(\gamma)$ 的平移部分。特别地，映射 $\operatorname{Trans}_\gamma\colon V_{\mathbb R}\to M_{\mathbb R}$ 为线性映射。令
\[
\rho'(\gamma)(m,\delta)
=\bigl(\widetilde\rho(\gamma)(m)+\operatorname{Trans}_\gamma(\delta),\;\delta\bigr),
\]
就容易像前面一样证明
\[
\rho'(\gamma)\circ\delta'\circ\Psi_\gamma'=\delta'.
\]
注意，只要 $\delta\in V$，就有 $\operatorname{Trans}_\gamma(\delta)\in M$，所以 $\rho'(\gamma)\in\operatorname{GL}(M\oplus V)$。因此，在 $B'$ 上诱导的仿射流形结构仍为整仿射结构，得到复流形 $X(B')$、$X(\mathcal V)$，以及全纯映射 $f\colon X(B')\to X(\mathcal V)$。

\SourceParagraph{ch06-6.2.4-007}{453}
注意，$X(\mathcal V)=(V_{\mathbb R}+\mathrm i\mathcal V)/V$。若 $\mathcal V$ 为锥，则 $V_{\mathbb R}+\mathrm i\mathcal V$ 为管状域，而 $X(\mathcal V)$ 是它对格的商。

\SourceParagraph{ch06-6.2.4-008}{453}
希望明确描述映射 $f\colon X(B')\to X(\mathcal V)$ 的纤维。稍后会看到，这些纤维形如 $X(B,\mathbf B)$，对应于 $B$ 上的不同仿射结构选择与 B 场选择。要理解如何确定 B 场，需要如下构造。

\SourceParagraph{ch06-6.2.4-009}{453}
\BookStatementNumber{29}
\begin{construction}\label{constr:6.29}
考虑 $\delta_1+\mathrm i\delta_2\in V_{\mathbb R}+\mathrm i\mathcal V$。按假设，$\delta_2\colon\widetilde B\to M_{\mathbb R}$ 为浸入，因而在 $B$ 上诱导仿射结构，得到 $B$ 上的层 $\Lambda$ 与 $\Lambda_{\mathbb R}$。给定这一结构，断言 $\delta_1$ 诱导自然的类 $[\delta_1]\in H^1(B,\Lambda_{\mathbb R})$。为说明这一点，令 $\pi\colon\widetilde B\to B$ 为万有覆盖，取任意充分小的开集 $U\subseteq B$，使得 $\pi^{-1}(U)$ 是形如 $\Psi_\gamma(W)$ 的开集的不交并，其中，$\gamma\in\pi_1(B)$，$W\subseteq\widetilde B$ 为某个开集。把 $U$ 与 $W$ 识别，浸入 $\delta_2|_W\colon U\to M_{\mathbb R}$ 的微分 $\delta_{2*}$，给出 $TU$ 与 $\delta_2(U)\times M_{\mathbb R}$ 的识别。因此，$\delta_1|_W\colon U\to M_{\mathbb R}$ 的图像，给出 $TU$ 的截面。

\SourceParagraph{ch06-6.2.4-010}{453}
若将 $W$ 换成 $W'=\Psi_\gamma(W)$，则对于 $i=1,2$，得到 $\delta_i'\colon U\to M_{\mathbb R}$，满足 $\delta_i'=\delta_i\circ\Psi_\gamma=\rho_i(\gamma)^{-1}\circ\delta_i$，其中，$\rho_i=\rho_{\delta_i}$。此外，$\delta_{2*}'\colon TU\to\delta_2'(U)\times M_{\mathbb R}$ 为
\[
\delta_{2*}'=
\bigl(\rho_2(\gamma)^{-1}\times\widetilde\rho_2(\gamma)^{-1}\bigr)\circ\delta_{2*},
\]
因为 $\widetilde\rho_2(\gamma)^{-1}$ 是 $\rho_2(\gamma)^{-1}$ 的微分。于是，有交换图
\[
\BookImage{02_Unnumbered_Diagrams/Chapter_06/Ch06_Diagram_PDFp453_01_Bookp443.pdf}
\]
以及图像
\[
\begin{aligned}
\{(\delta_2(b),\delta_1(b))\mid b\in U\}
&\subseteq\delta_2(U)\times M_{\mathbb R},\\
\{(\delta_2'(b),\delta_1'(b))\mid b\in U\}
&=\{(\rho_2(\gamma)^{-1}\delta_2(b),
       \rho_1(\gamma)^{-1}\delta_1(b))\mid b\in U\}\\
&\subseteq\delta_2'(U)\times M_{\mathbb R}.
\end{aligned}
\]
后一个图像在 $\delta_2(U)\times M_{\mathbb R}$ 中的像为
\[
\begin{aligned}
&\{(\delta_2(b),\widetilde\rho_2(\gamma)\rho_1(\gamma)^{-1}\delta_1(b))\mid b\in U\}\\
&\qquad=\{(\delta_2(b),\delta_1(b)-\operatorname{Trans}_\gamma(\delta_1))\mid b\in U\}.
\end{aligned}
\]
因此，两个截面只相差 $\operatorname{Trans}_\gamma(\delta_1)$，可以将其看作 $\Gamma(U,\Lambda_{\mathbb R})$ 的元素。

\SourceParagraph{ch06-6.2.4-011}{454}
现在，选择 $B$ 的开覆盖 $\{U_i\}$，每个开集都具有上述性质，并选择位于 $U_i$ 上方的开集 $W_i\subseteq\widetilde B$。这样，在 $\delta_2$ 于 $W_i$ 上诱导的识别下，得到作为 $\delta_1$ 图像的截面 $\Gamma_i\in\Gamma(U_i,TB)$。此外，$\Gamma_j-\Gamma_i$ 是 $\Gamma(U_i\cap U_j,\Lambda_{\mathbb R})$ 的元素。这给出类 $[\delta_1]\in H^1(B,\Lambda_{\mathbb R})$ 的 Čech 代表。

\SourceParagraph{ch06-6.2.4-012}{454}
注意，若对所有 $\gamma$，都有 $\operatorname{Trans}_\gamma(\delta_1)\in M$，则 $[\delta_1]$ 来自 $H^1(B,\Lambda)$。当 $\delta_1$ 属于 $V$，而不仅仅属于 $V_{\mathbb R}$ 时，就是这种情形。因此，若 $\delta_1$ 只在模 $V$ 的意义下定义，则 $[\delta_1]$ 定义于
\[
H^1(B,\Lambda_{\mathbb R})/\operatorname{im}H^1(B,\Lambda),
\]
或者定义于 $H^1(B,\Lambda_{\mathbb R}/\Lambda)$。\BookQED
\end{construction}

\SourceParagraph{ch06-6.2.4-013}{454}
\BookStatementNumber{30}
\begin{note}[numbered]\label{note:6.30}
若 $\delta_1=\delta_2=\delta$，则类 $[\delta]\in H^1(B,\Lambda_{\mathbb R})$，称为由 $\delta$ 诱导的 $B$ 上仿射结构的辐射障碍。细节见文献\BookCite{ref206}第1节，或文献\BookCite{ref184}。由于 $\check X(B)$ 上的辛形式 $\omega$ 是 $(1,1)$ 型半平坦2形式，它代表 $H^1(B,\Lambda_{\mathbb R})$ 中的一个类。在前述对偶符号约定下，这个类与 $[\delta]$ 相同。细节留给读者。
\end{note}

\SourceParagraph{ch06-6.2.4-014}{454}
\BookStatementNumber{31}
\begin{proposition}\label{prop:6.31}
令 $\delta_1+\mathrm i\delta_2$ 代表 $X(\mathcal V)=(V_{\mathbb R}+\mathrm i\mathcal V)/V$ 中的点，因而 $\delta_2$ 在 $B$ 上诱导仿射结构。那么，取 $\mathbf B=[\delta_1]\in H^1(B,\Lambda_{\mathbb R}/\Lambda)$，有同构
\[
f^{-1}(\delta_1+\mathrm i\delta_2)\cong X(B,\mathbf B).
\]
\end{proposition}
\begin{proof}
\SourceParagraph{ch06-6.2.4-015}{454}
有 $B'=B\times\mathcal V$，以及 $\delta'$ 在 $B'$ 上诱导的整仿射结构。这给出局部系统 $\Lambda'\subseteq TB'$。另一方面，$B\times\{\delta_2\}\subseteq B'$ 本身为仿射流形，仿射结构由 $\delta_2$ 给出，从而得到局部系统 $\Lambda\subseteq T(B\times\{\delta_2\})$。显然，对 $b\in B\times\{\delta_2\}$，有
\[
\Lambda_b=\{v\in\Lambda_b'\mid v\text{ 与 }B\times\{\delta_2\}\text{ 相切}\}.
\]
现在，映射 $f\colon X(B')\to X(\mathcal V)$，由投影 $B'\to\mathcal V$ 的微分 $p_*\colon TB'\to T\mathcal V=\mathrm i\mathcal V\times V_{\mathbb R}$ 诱导。因此，例如有
\[
p_*^{-1}(\mathrm i\delta_2,0)=T(B\times\{\delta_2\}),
\]
从而
\[
f^{-1}(\mathrm i\delta_2)=T(B\times\{\delta_2\})/\Lambda=X(B).
\]
更一般地，$TB'|_{B\times\{\delta_2\}}=T(B\times\{\delta_2\})\times V_{\mathbb R}$，并且
\[
p_*^{-1}(\mathrm i\delta_2,\delta_1)
=T(B\times\{\delta_2\})\times\{\delta_1\}.
\]
显然，后一个空间，只是 $TB'$ 中子丛 $T(B\times\{\delta_2\})$ 的一个平移，但这一说法并不识别它们的复流形结构。要确定 $T(B\times\{\delta_2\})\times\{\delta_1\}$ 上的复结构，需要更仔细地研究 $\Lambda'|_{B\times\{\delta_2\}}$。

\SourceParagraph{ch06-6.2.4-016}{455}
为此，首先考虑足够小的开集 $U\subseteq B$，使得如\BookRef{constr:6.29}{构造6.29}中一样，$\pi^{-1}(U)=\coprod_\gamma\Psi_\gamma(W)$，得到浸入 $\delta_2\colon U\to M_{\mathbb R}$。先用 $\delta_2$ 把 $U$ 与 $\delta_2(U)\subseteq M_{\mathbb R}$ 识别，使得 $\delta_2$ 实际上为恒等映射，而 $U\times\mathcal V$ 带有由 $\delta'$ 诱导的仿射结构。令 $b\in U$。那么，$T(U\times\mathcal V)_{(b,\delta_2)}=M_{\mathbb R}\oplus V_{\mathbb R}$，并有微分
\[
\delta_*'\colon T(U\times\mathcal V)_{(b,\delta_2)}
\longrightarrow M_{\mathbb R}\oplus V_{\mathbb R},
\]
按如下方式计算。若 $(m,\delta)\in M_{\mathbb R}\oplus V_{\mathbb R}$，令 $\gamma\colon(-\epsilon,\epsilon)\to U\times\mathcal V$ 为 $\gamma(t)=(b,\delta_2)+t(m,\delta)$。那么，
\[
\begin{aligned}
\delta_*'(m,\delta)
&=\left.\frac{\mathrm d}{\mathrm dt}(\delta'\circ\gamma(t))\right|_{t=0}\\
&=\left.\frac{\mathrm d}{\mathrm dt}\delta'(b+tm,\delta_2+t\delta)\right|_{t=0}\\
&=\left.\frac{\mathrm d}{\mathrm dt}
\bigl(\delta_2(b+tm)+t\delta(b+tm),\;\delta_2+t\delta\bigr)\right|_{t=0}.
\end{aligned}
\]
回顾一下，已经取 $\delta_2$ 为恒等映射，因此得到
\[
\delta_*'(m,\delta)=(m+\delta(b),\delta).
\]
特别地，$\delta_*'(-\delta(b),\delta)=(0,\delta)$，而
\[
b\longmapsto(-\delta(b),\delta)
\]
是 $T(U\times\mathcal V)|_{U\times\{\delta_2\}}$ 的截面，并且属于 $\Lambda_{\mathbb R}'|_{U\times\{\delta_2\}}$。所以，通过截面 $b\mapsto(-\delta_1(b),\delta_1)$ 作平移，就在
\[
T(U\times\{\delta_2\})\times\{0\}
\subseteq T(U\times\mathcal V)|_{U\times\{\delta_2\}}
\]
与 $T(U\times\{\delta_2\})\times\{\delta_1\}$ 之间诱导双全纯映射。

\SourceParagraph{ch06-6.2.4-017}{455}
现在，如\BookRef{constr:6.29}{构造6.29}中一样，用开集 $\{U_i\}$ 覆盖 $B$。作与前面一样的选择，得到 $\delta_1$ 的图像 $\Gamma_i\subseteq T(U_i\times\{\delta_2\})$。那么，已经说明，映射
\[
\psi_i\colon T(U_i\times\{\delta_2\})\times\{0\}
\longrightarrow T(U_i\times\{\delta_2\})\times\{\delta_1\},
\]
由 $T(U_i\times\mathcal V)$ 的截面 $(-\Gamma_i,\delta_1)$ 作平移给出，并且为双全纯映射。因此，有交换图
\[
\BookImage{02_Unnumbered_Diagrams/Chapter_06/Ch06_Diagram_PDFp456_01_Bookp446.pdf}
\]
并且看到，$p_*^{-1}(\mathrm i\delta_2,\delta_1)$ 是用上闭链
\[
(U_{ij},\Gamma_j-\Gamma_i)=[\delta_1]
\]
重新粘合 $TB$ 得到的。因此，$f^{-1}(\delta_1+\mathrm i\delta_2)$ 是用
\[
\mathbf B=[\delta_1]\in H^1(B,\Lambda_{\mathbb R}/\Lambda)
\]
重新粘合 $X(B)$ 得到的。
\end{proof}

\SourceParagraph{ch06-6.2.4-018}{456}
族 $f\colon X(B')\to X(\mathcal V)$，在拓扑上并不平凡。要确切地理解这一点，选取 $\delta_2\in\mathcal V$ 与非零 $\delta_1\in V$，考虑 $X(\mathcal V)$ 中由
\[
S^1=(\mathbb R\delta_1/\mathbb Z\delta_1)+\mathrm i\delta_2
\]
参数化的环路。固定 $0+\mathrm i\delta_2$ 为基点。对 $0\leq r\leq1$，将定义一个连续的微分同胚族
\[
\varphi_r\colon f^{-1}(\mathrm i\delta_2)
\longrightarrow f^{-1}(r\delta_1+\mathrm i\delta_2).
\]
由于 $f^{-1}(\mathrm i\delta_2)=f^{-1}(\delta_1+\mathrm i\delta_2)$，$\varphi_1$ 给出 $f^{-1}(\mathrm i\delta_2)$ 的一个自同构，称为单值变换微分同胚。若它与恒等映射不同痕，则沿这一环路拉回的族 $f^{-1}(S^1)\to S^1$ 非平凡。更准确地说，在拓扑上，可以从 $[0,1]\times f^{-1}(\mathrm i\delta_2)$ 出发，通过 $\varphi_1$ 将 $\{0\}\times f^{-1}(\mathrm i\delta_2)$ 与 $\{1\}\times f^{-1}(\mathrm i\delta_2)$ 识别，得到这个族。

\SourceParagraph{ch06-6.2.4-019}{456}
要定义 $\varphi_r$，只需在切丛层面考察
\[
p_*\colon T(B\times\{\delta_2\})\times V_{\mathbb R}
\longrightarrow V_{\mathbb R}.
\]
通过将 $p_*^{-1}(0)$ 中的每个点 $(x,c)$ 加上 $(0,r\delta_1)$，得到
\[
\widetilde\varphi_r\colon p_*^{-1}(0)\longrightarrow p_*^{-1}(r\delta_1).
\]
不过，在\BookRef{prop:6.31}{命题6.31}的证明中已经局部地看到，模去 $\Lambda'$ 后，可以通过加上 $TB'|_{U\times\{\delta_2\}}$ 的局部截面 $b\mapsto(-\delta_1(b),\delta_1)$，将 $p_*^{-1}(0)$ 与 $p_*^{-1}(\delta_1)$ 识别。因此，在 $f^{-1}(\mathrm i\delta_2)=TB/\Lambda$ 上，$\varphi_1$ 局部地由沿纤维平移 $T_bB\ni y\mapsto y+\delta_1(b)$ 描述。对于不同的开集选择，截面 $b\mapsto\delta_1(b)$ 只相差 $\Lambda$ 中的元素，因为 $\delta_1$ 属于 $V$，而不仅仅属于 $V_{\mathbb R}$。所以，$b\mapsto\delta_1(b)\bmod\Lambda$ 定义 $X(B,0)\to B$ 的截面，而 $\varphi_1$ 由这一截面给出的沿纤维平移构成。

\SourceParagraph{ch06-6.2.4-020}{457}
因此，已经证明如下命题。

\SourceParagraph{ch06-6.2.4-021}{457}
\BookStatementNumber{32}
\begin{proposition}\label{prop:6.32}
由 $X(\mathcal V)$ 中形如 $(\mathbb R\delta_1/\mathbb Z\delta_1)+\mathrm i\delta_2$ 的环路诱导的单值变换 $\varphi_1\colon X(B)\to X(B)$，由 $X(B)$ 的一个截面对 $X(B)$ 作平移给出，而这一截面由 $\delta_1$ 决定。
\end{proposition}

\subsection{广义复结构}\label{sec:6.2.5}

\SourceParagraph{ch06-6.2.5-001}{457}
现在，来处理讨论 B 场时出现的一个疑问。传统上，对于 Calabi--Yau 流形 $X$，B 场取值于 $H^2(X,\mathbb R)/\operatorname{im}H^2(X,\mathbb Z)$，或 $H^2(X,\mathbb R/\mathbb Z)$。不过，在\BookSectionRef{sec:6.2.3}{6.2.3}中，已经把 B 场放在 $H^1(B,\Lambda_{\mathbb R}/\Lambda)$ 中，或者用 $(1,1)$ 型半平坦2形式代表 B 场。一般而言，这不能给出 $H^2(X,\mathbb R/\mathbb Z)$ 的全部元素。例如，若 $B$ 是带有标准仿射结构的二维环面，则 $H^1(B,\Lambda_{\mathbb R}/\Lambda)=(\mathbb R/\mathbb Z)^4$，而 $H^2(X(B),\mathbb R/\mathbb Z)=(\mathbb R/\mathbb Z)^6$。再考虑到复二维环面的复结构模空间具有复维数4，这种维数上的不一致就更加明显。因此，存在一些额外的 B 场选择，它们的意义并不立即清楚。

\SourceParagraph{ch06-6.2.5-002}{457}
Kapustin\BookCite{ref284}首先提出，解决这一疑问的方法，是考虑 Hitchin\BookCite{ref236}的广义复结构概念，从而扩大所涉及的模空间。Ben-Bassat\BookCite{ref038,ref039}与 Chiantese、Gmeiner、Jeschek\BookCite{ref096}进一步发展了这一想法。这里来解释这个概念。

\SourceParagraph{ch06-6.2.5-003}{457}
从 Lie 括号的一个推广开始。令 $M$ 为实流形，令 $X+\xi,Y+\eta$ 为 $TM\oplus\bigwedge^pT^*M$ 的 $C^\infty$ 截面。如下定义 Courant 括号：
\[
[X+\xi,Y+\eta]
=[X,Y]+\mathcal L_X\eta-\mathcal L_Y\xi
-\tfrac12\mathrm d\bigl(\iota(X)\eta-\iota(Y)\xi\bigr).
\]
这里，$\mathcal L$ 表示 Lie 导数。特别关注 $p=1$ 的情形。将利用 $TM\oplus T^*M$ 上自然的、符号型为 $(n,n)$ 的内积，其中，$n=\dim_{\mathbb R}M$：
\[
\langle X+\xi,Y+\eta\rangle
=\tfrac12\bigl(\xi(Y)+\eta(X)\bigr).
\]

\SourceParagraph{ch06-6.2.5-004}{457}
于是，有如下定义。

\SourceParagraph{ch06-6.2.5-005}{457}
\BookStatementNumber{33}
\begin{definition}\label{def:6.33}
偶数维流形 $M$ 上的广义复结构，是子丛 $E\subseteq(TM\oplus T^*M)\otimes\mathbb C$，满足：
\begin{enumerate}[label=(\arabic*)]
\item $(TM\oplus T^*M)\otimes\mathbb C=E\oplus\overline E$；
\item $E$ 的截面空间在 Courant 括号下封闭；
\item $E$ 相对于 $TM\oplus T^*M$ 上的内积为各向同性子丛。
\end{enumerate}
\end{definition}

\SourceParagraph{ch06-6.2.5-006}{457}
第三个条件可以如下重新表述。把 $E$ 与 $\overline E$ 看作算子 $\mathcal J\colon TM\oplus T^*M\to TM\oplus T^*M$ 的 $\pm\mathrm i$ 特征空间，其中，$\mathcal J^2=-1$。那么，第三个条件等价于 $\mathcal J$ 相对于不定内积为正交算子。事实上，若 $\mathcal J$ 正交且 $X+\xi\in E$，则
\[
\langle X+\xi,X+\xi\rangle
=\langle\mathcal J(X+\xi),\mathcal J(X+\xi)\rangle
=-\langle X+\xi,X+\xi\rangle.
\]
反向结论，也同样容易证明。

\SourceParagraph{ch06-6.2.5-007}{458}
\BookStatementNumber{34}
\begin{example}\label{ex:6.34}
\begin{enumerate}[label=(\arabic*)]
\item 若 $M$ 有概复结构 $J$，则可以取 $E=T^{0,1}M\oplus\Omega^{1,0}M$，即取 $(0,1)$ 型向量场与 $(1,0)$ 型1形式。广义复结构定义中的条件(1)和(3)显然成立。使用 Lie 导数的 Cartan 公式，以及 $(1,0)$ 形式在 $(0,1)$ 向量场上取值为零这一事实，条件(2)化为：只要 $X+\xi,Y+\eta\in E$，就有
\[
[X,Y]+\iota(X)\mathrm d\eta-\iota(Y)\mathrm d\xi\in E.
\]
这显然等价于两个条件：$T^{0,1}M$ 在 Lie 括号下封闭，以及 $(1,0)$ 形式的微分具有 $(2,0)+(1,1)$ 型。不过，由 Newlander--Nirenberg 定理，这两个条件都等价于概复结构 $J$ 的可积性。因此，广义复结构定义中的条件(2)，等价于 $J$ 的可积性。注意，
\[
\mathcal J=\begin{pmatrix}-J&0\\0&J^*\end{pmatrix}.
\]
\end{enumerate}

\SourceParagraph{ch06-6.2.5-008}{458}
\begin{enumerate}[label=(\arabic*),start=2]
\item 令 $M$ 为偶数维流形，带非退化2形式 $\omega$。那么，$\omega$ 通过 $\omega(X)=\iota(X)\omega$，诱导识别 $\omega\colon TM\to T^*M$。取
\[
\mathcal J=\begin{pmatrix}0&\omega^{-1}\\-\omega&0\end{pmatrix}.
\]
$+\mathrm i$ 特征空间 $E$ 为
\[
E=\{X+\mathrm i\omega(X)\mid X\in TM\otimes\mathbb C\}.
\]
同样，\BookRef{def:6.33}{定义6.33}的条件(1)和(3)显然成立。对于条件(2)，前一个例子中的论证说明，$E$ 在 Courant 括号下封闭，等价于 $\omega$ 为闭形式，也就是 $M$ 为辛流形。
\end{enumerate}

\SourceParagraph{ch06-6.2.5-009}{458}
\begin{enumerate}[label=(\arabic*),start=3]
\item 若 $B$ 为闭2形式，则得到 $TM\oplus T^*M$ 的自同构 $\mathbf B$：
\[
\mathbf B\colon X+\xi\longmapsto X+\xi+\iota(X)B.
\]
它满足 $\mathbf B[X+\xi,Y+\eta]=[\mathbf B(X+\xi),\mathbf B(Y+\eta)]$。事实上，反复使用 Cartan 公式，得到
\[
\begin{aligned}
[\mathbf B(X+\xi),\mathbf B(Y+\eta)]
&=[X+\xi,Y+\eta]+[X,\iota(Y)B]+[\iota(X)B,Y]\\
&=[X+\xi,Y+\eta]+\mathcal L_X\iota(Y)B-\mathcal L_Y\iota(X)B\\
&\qquad+\mathrm d(\iota(Y)\iota(X)B)\\
&=[X+\xi,Y+\eta]+\mathcal L_X\iota(Y)B-\mathrm d(\iota(Y)\iota(X)B)\\
&\qquad-\iota(Y)\mathrm d(\iota(X)B)+\mathrm d(\iota(Y)\iota(X)B)\\
&=[X+\xi,Y+\eta]+\mathcal L_X\iota(Y)B-\iota(Y)\mathcal L_XB
+\iota(Y)\iota(X)\mathrm dB\\
&=[X+\xi,Y+\eta]+\iota([X,Y])B+\iota(Y)\iota(X)\mathrm dB\\
&=\mathbf B[X+\xi,Y+\eta].
\end{aligned}
\]
倒数第二个等式是 Lie 导数的一个性质，参见例如文献\BookCite{ref319}第140页的命题5.3。最后一行，由 $B$ 为闭形式得到。

\SourceParagraph{ch06-6.2.5-010}{459}
此外，$TM\oplus T^*M$ 上的内积，显然在 $\mathbf B$ 的作用下保持不变。因此，给定广义复结构 $E$，$\mathbf B(E)$ 也是广义复结构。这就是用 B 场作扭曲。
\end{enumerate}
\end{example}

\SourceParagraph{ch06-6.2.5-011}{459}
有了广义复结构的概念，也应当有广义 Calabi--Yau 流形的概念。\BookRef{prop:6.24}{命题6.24}为此提供启发：Calabi--Yau 流形上的复结构，完全由全纯 $n$ 形式确定。将用一种更一般的形式，称为纯旋量，替代这一形式。考虑 $M$ 上的正、负旋量丛：若 $\dim M=n$，则
\[
\begin{aligned}
S^+&=\left(\bigwedge\nolimits^{\mathrm{even}}T^*M\right)
\otimes\left(\bigwedge^nTM\right)^{1/2},\\
S^-&=\left(\bigwedge\nolimits^{\mathrm{odd}}T^*M\right)
\otimes\left(\bigwedge^nTM\right)^{1/2}.
\end{aligned}
\]
用 $S^\pm$ 表示 $S^+$ 或 $S^-$ 中的任意一个。对于纯次数为 $p$ 的形式 $\varphi$，定义 $\sigma\colon S^\pm\to S^\pm$ 为
\[
\sigma(\varphi)=(-1)^{p(p-1)/2}\varphi.
\]
于是，有自然的双线性形式
\[
\langle\varphi,\psi\rangle=(\sigma(\varphi)\wedge\psi)_{\mathrm{top}}.
\]
这里，取最高次数部分，得到 $\bigwedge^nT^*M\otimes\bigwedge^nTM$ 的截面，再取迹，得到一个数。等价地，也可以从 $S^\pm$ 的定义中去掉 $\bigwedge^nTM$ 的平方根因子，使配对取值于 $\bigwedge^nT^*M$；以下将采用这一做法，把旋量直接看作微分形式。$TM\oplus T^*M$ 通过 Clifford 乘法作用于旋量，交换奇偶分量：
\[
(X+\xi)\cdot\varphi=\iota(X)\varphi+\xi\wedge\varphi.
\]

\SourceParagraph{ch06-6.2.5-012}{460}
\begin{definition*}[纯旋量]
旋量 $\varphi\in S^\pm\otimes\mathbb C$ 为纯旋量，是指子丛
\[
E_\varphi=\{X+\xi\in(TM\oplus T^*M)\otimes\mathbb C
\mid(X+\xi)\cdot\varphi=0\}
\]
是 $(TM\oplus T^*M)\otimes\mathbb C$ 的极大各向同性子丛。
\end{definition*}

\SourceParagraph{ch06-6.2.5-013}{460}
将需要以下引理。

\SourceParagraph{ch06-6.2.5-014}{460}
\BookStatementNumber{35}
\begin{lemma}\label{lem:6.35}
对于纯旋量 $\varphi,\psi$，$E_\varphi\cap E_\psi=0$，当且仅当 $\langle\varphi,\psi\rangle\ne0$。
\end{lemma}
\begin{proof}
\SourceParagraph{ch06-6.2.5-015}{460}
见文献\BookCite{ref088}，III 2.4。
\end{proof}

\SourceParagraph{ch06-6.2.5-016}{460}
\BookStatementNumber{36}
\begin{definition}\label{def:6.36}
广义 Calabi--Yau 流形，是偶数维流形 $M$，带有闭形式 $\varphi\in S^\pm\otimes\mathbb C$，其中，$\varphi$ 为纯旋量，且 $\langle\varphi,\overline\varphi\rangle$ 处处非零。
\end{definition}

\SourceParagraph{ch06-6.2.5-017}{460}
\BookStatementNumber{37}
\begin{example}\label{ex:6.37}
\begin{enumerate}[label=(\arabic*)]
\item 若 $M$ 为 Calabi--Yau 流形，带有处处非零的全纯 $n$ 形式 $\Omega$，则 $\Omega$ 为纯旋量：$(X+\xi)\cdot\Omega=0$，当且仅当 $X$ 为 $(0,1)$ 型向量场，$\xi$ 为 $(1,0)$ 型形式。因此，$E_\Omega$ 是与 $M$ 的复结构关联的广义复结构。可以看到，广义 Calabi--Yau 流形定义中的条件，与\BookRef{prop:6.24}{命题6.24}中的条件对应。
\end{enumerate}

\SourceParagraph{ch06-6.2.5-018}{460}
\begin{enumerate}[label=(\arabic*),start=2]
\item 若 $\varphi=\exp(\mathrm i\omega)$，其中，$\omega$ 为 $M$ 上的辛形式，则 $(X+\xi)\cdot\varphi=0$，当且仅当 $\xi=-\mathrm i\iota(X)\omega$。因此，$E_\varphi$ 是\BookRef{ex:6.34}{例题6.34}(2)中的广义复结构的共轭；该例中的 $+\mathrm i$ 特征丛本身，由 $\exp(-\mathrm i\omega)$ 给出。
\end{enumerate}

\SourceParagraph{ch06-6.2.5-019}{460}
\begin{enumerate}[label=(\arabic*),start=3]
\item B 场 $B$ 通过 $\varphi\mapsto\exp(B)\wedge\varphi$ 作用于旋量。现在，
\[
\begin{aligned}
\bigl(\mathbf B^{-1}(X+\xi)\bigr)\cdot(\exp(B)\wedge\varphi)
&=(X+\xi-\iota(X)B)\cdot(\exp(B)\wedge\varphi)\\
&=\iota(X)(\exp(B)\wedge\varphi)
+\xi\wedge\exp(B)\wedge\varphi\\
&\qquad-(\iota(X)B)\wedge\exp(B)\wedge\varphi\\
&=(\iota(X)B)\wedge\exp(B)\wedge\varphi
+\exp(B)\wedge\iota(X)\varphi\\
&\qquad+\xi\wedge\exp(B)\wedge\varphi
-(\iota(X)B)\wedge\exp(B)\wedge\varphi\\
&=\exp(B)\wedge\bigl((X+\xi)\cdot\varphi\bigr).
\end{aligned}
\]
因此，若 $\varphi$ 为纯旋量，则 $\exp(B)\wedge\varphi$ 也为纯旋量。此外，若 $\varphi$ 定义广义 Calabi--Yau 结构，且 B 场为闭实2形式，则
\[
\langle\exp(B)\wedge\varphi,\exp(B)\wedge\overline\varphi\rangle
=\langle\varphi,\overline\varphi\rangle
\]
处处非零。因此，$\exp(B)\wedge\varphi$ 定义广义 Calabi--Yau 结构，对应于用 $-B$ 扭曲相关联的广义复结构。
\end{enumerate}
\end{example}

\SourceParagraph{ch06-6.2.5-020}{460}
这里的关键是，广义 Calabi--Yau 结构给出广义复结构。

\SourceParagraph{ch06-6.2.5-021}{460}
\BookStatementNumber{38}
\begin{theorem}\label{thm:6.38}
若 $\varphi$ 给出广义 Calabi--Yau 结构，则 $E_\varphi$ 为广义复结构。
\end{theorem}
\begin{proof}
\SourceParagraph{ch06-6.2.5-022}{461}
已经知道，$E_\varphi$ 为极大各向同性子丛，因此满足广义复结构的第三个条件。另一方面，由于 $\langle\varphi,\overline\varphi\rangle$ 处处非零，\BookRef{lem:6.35}{引理6.35}给出
\[
0=E_\varphi\cap E_{\overline\varphi}
=E_\varphi\cap\overline{E_\varphi}.
\]
这给出第一个条件。所以，需要证明 $E_\varphi$ 在 Courant 括号下封闭。设 $X+\xi,Y+\eta\in E_\varphi$。那么，
\[
\iota(X)\varphi+\xi\wedge\varphi=0
=\iota(Y)\varphi+\eta\wedge\varphi.
\]
利用 $\mathrm d\varphi=0$ 与 Cartan 公式，得到
\[
\begin{aligned}
\iota([X,Y])\varphi
&=\mathcal L_X(\iota(Y)\varphi)-\iota(Y)\mathcal L_X\varphi\\
&=-\mathcal L_X(\eta\wedge\varphi)-\iota(Y)\mathrm d(\iota(X)\varphi)\\
&=-(\mathcal L_X\eta)\wedge\varphi-\eta\wedge\mathcal L_X\varphi
+\iota(Y)\mathrm d(\xi\wedge\varphi)\\
&=-(\mathcal L_X\eta)\wedge\varphi-\eta\wedge\mathrm d(\iota(X)\varphi)
+\iota(Y)((\mathrm d\xi)\wedge\varphi)\\
&=-(\mathcal L_X\eta)\wedge\varphi+\eta\wedge\mathrm d(\xi\wedge\varphi)
+\iota(Y)((\mathrm d\xi)\wedge\varphi)\\
&=-(\mathcal L_X\eta)\wedge\varphi+\eta\wedge\mathrm d\xi\wedge\varphi
+(\iota(Y)\mathrm d\xi)\wedge\varphi+(\mathrm d\xi)\wedge\iota(Y)\varphi\\
&=-(\mathcal L_X\eta)\wedge\varphi+\eta\wedge\mathrm d\xi\wedge\varphi
+(\iota(Y)\mathrm d\xi)\wedge\varphi-(\mathrm d\xi)\wedge\eta\wedge\varphi\\
&=-(\mathcal L_X\eta)\wedge\varphi+(\iota(Y)\mathrm d\xi)\wedge\varphi.
\end{aligned}
\]
因此，由反对称性，
\[
\begin{aligned}
\iota([X,Y])\varphi
&=\tfrac12\bigl(\iota([X,Y])\varphi-\iota([Y,X])\varphi\bigr)\\
&=\tfrac12\bigl(-(\mathcal L_X\eta)\wedge\varphi
+(\iota(Y)\mathrm d\xi)\wedge\varphi\\
&\hspace{17mm}+(\mathcal L_Y\xi)\wedge\varphi
-(\iota(X)\mathrm d\eta)\wedge\varphi\bigr)\\
&=\bigl[\iota(Y)\mathrm d\xi+\tfrac12\mathrm d(\iota(Y)\xi)
-\iota(X)\mathrm d\eta-\tfrac12\mathrm d(\iota(X)\eta)\bigr]\wedge\varphi\\
&=\bigl[\mathcal L_Y\xi-\mathcal L_X\eta
-\tfrac12\mathrm d(\iota(Y)\xi-\iota(X)\eta)\bigr]\wedge\varphi.
\end{aligned}
\]
由 Courant 括号的定义，这蕴含 $[X+\xi,Y+\eta]\cdot\varphi=0$，正如所要证明的。
\end{proof}

\SourceParagraph{ch06-6.2.5-023}{461}
现在，回到镜像对称。重新采用基本的设定：$B$ 为仿射流形，转移映射属于 $\mathbb R^n\rtimes\operatorname{SL}_n(\mathbb Z)$。只要旋量 $\varphi$ 是各种类型的半平坦形式之和，就可以称它为半平坦旋量。于是，得到 $X(B)$ 与 $\check X(B)$ 上的半平坦广义 Calabi--Yau 结构。特别地，利用\BookSectionRef{sec:6.2.2}{6.2.2}中的 $\mu$ 的定义，有以下定理。

\SourceParagraph{ch06-6.2.5-024}{461}
\BookStatementNumber{39}
\begin{theorem}\label{thm:6.39}
若 $\varphi$ 为半平坦纯旋量，定义 $X(B)$ 上的广义 Calabi--Yau 结构，则 $\mu(\varphi)$ 是半平坦纯旋量，定义 $\check X(B)$ 上的广义 Calabi--Yau 结构。
\end{theorem}
\begin{proof}
\SourceParagraph{ch06-6.2.5-025}{462}
需要理解，镜像变换如何作用于 $T\oplus T^*$ 的元素。令 $S_{1,0}$ 与 $S_{0,1}$ 分别表示 $X(B)$ 上 $(1,0)$ 型与 $(0,1)$ 型半平坦向量场的层。这里，$S_{1,0}$ 由形如 $\sum_{i=1}^na_i(y)\partial/\partial y_i$ 的向量场给出，而 $S_{0,1}$ 由形如 $\sum_{i=1}^na_i(y)\partial/\partial x_i$ 的向量场给出。同样，这些层不依赖于局部仿射坐标的选择。类似地，定义 $\check S_{1,0}$ 与 $\check S_{0,1}$ 为 $\check X(B)$ 上这两种类型的半平坦向量场层，分别由局部形如 $\sum_{i=1}^na_i(y)\partial/\partial y_i$ 与 $\sum_{i=1}^na_i(y)\partial/\partial\check x_i$ 的向量场构成。继续沿用\BookSectionRef{sec:6.2.2}{6.2.2}中的记号 $S^{p,q}$ 表示 $(p,q)$ 型半平坦形式，这里也用同一记号表示相应的层。注意，有
\[
\begin{aligned}
S^{1,0}&\cong T^*B,&
S^{0,1}&\cong\check\Lambda_{\mathbb R}\otimes_{\mathbb R}C^\infty(B),\\
\check S^{1,0}&\cong T^*B,&
\check S^{0,1}&\cong\Lambda_{\mathbb R}\otimes_{\mathbb R}C^\infty(B),\\
S_{1,0}&\cong TB,&
S_{0,1}&\cong\Lambda_{\mathbb R}\otimes_{\mathbb R}C^\infty(B),\\
\check S_{1,0}&\cong TB,&
\check S_{0,1}&\cong\check\Lambda_{\mathbb R}\otimes_{\mathbb R}C^\infty(B).
\end{aligned}
\]
因此，存在自然同构
\[
\begin{aligned}
\widetilde\mu\colon&S_{1,0}\oplus S_{0,1}\oplus S^{1,0}\oplus S^{0,1}\\
&\longrightarrow\check S_{1,0}\oplus\check S_{0,1}
\oplus\check S^{1,0}\oplus\check S^{0,1},
\end{aligned}
\]
它将 $S_{0,1}$ 与 $\check S^{0,1}$ 识别，将 $S^{0,1}$ 与 $\check S_{0,1}$ 识别。具体地，在局部坐标中，采用与 $\mu$ 相容的符号约定，
\[
\begin{aligned}
\widetilde\mu(\mathrm dy_i)&=-\mathrm dy_i,&
\widetilde\mu\!\left(\frac{\partial}{\partial y_i}\right)&=-\frac{\partial}{\partial y_i},\\
\widetilde\mu(\mathrm dx_i)&=\frac{\partial}{\partial\check x_i},&
\widetilde\mu\!\left(\frac{\partial}{\partial x_i}\right)&=\mathrm d\check x_i.
\end{aligned}
\]

\SourceParagraph{ch06-6.2.5-026}{462}
若 $\varphi\in S^{p,q}$，则对于 $X+\xi\in S_{1,0}\oplus S_{0,1}\oplus S^{1,0}\oplus S^{0,1}$，注意有
\begin{equation}\label{eq:6.1}
\widetilde\mu(X+\xi)\cdot\mu(\varphi)
=\mu\bigl((X+\xi)\cdot\varphi\bigr).
\end{equation}
事实上，只需对 $X+\xi$ 分别为 $\partial/\partial y_i$、$\partial/\partial x_i$、$\mathrm dy_i$ 与 $\mathrm dx_i$ 的情形证明这一点。对于 $X+\xi=\partial/\partial y_i$，结论直接成立，因为 $\mu$ 对纤维部分作内乘，而底空间次数减少1时，其符号因子 $(-1)^p$ 改变符号。对于 $X+\xi=\mathrm dy_i$，同理。如果 $J\subseteq\{1,\ldots,n\}$ 为指标集，$j\notin J$ 且 $\#J=q$，那么
\[
\begin{aligned}
\widetilde\mu\!\left(\frac{\partial}{\partial x_j}\right)
\cdot\mu(\mathrm dy_I\wedge\mathrm dx_J)
&=\pm\mathrm d\check x_j\wedge\mathrm dy_I\wedge
\iota\!\left(\frac{\partial}{\partial\check x_J}\right)
(\mathrm d\check x_1\wedge\cdots\wedge\mathrm d\check x_n)\\
&=0,
\end{aligned}
\]
并且 $(\partial/\partial x_j)\cdot(\mathrm dy_I\wedge\mathrm dx_J)=0$。另一方面，若 $j\notin J$ 且 $\#J=q-1$，则
\[
\begin{aligned}
&\widetilde\mu\!\left(\frac{\partial}{\partial x_j}\right)
\cdot\mu(\mathrm dy_I\wedge\mathrm dx_j\wedge\mathrm dx_J)\\
&\quad=\mathrm d\check x_j\wedge(-1)^{p+q(q-1)/2+n(n-1)/2}\mathrm dy_I\wedge
\iota\!\left(\frac{\partial}{\partial\check x_j}\wedge
\frac{\partial}{\partial\check x_J}\right)
(\mathrm d\check x_1\wedge\cdots\wedge\mathrm d\check x_n)\\
&\quad=(-1)^{p+q(q-1)/2+n(n-1)/2}(-1)^p(-1)^{q-1}\mathrm dy_I\wedge
\iota\!\left(\frac{\partial}{\partial\check x_J}\right)
(\mathrm d\check x_1\wedge\cdots\wedge\mathrm d\check x_n)\\
&\quad=(-1)^{(q-1)(q-2)/2+n(n-1)/2}\mathrm dy_I\wedge
\iota\!\left(\frac{\partial}{\partial\check x_J}\right)
(\mathrm d\check x_1\wedge\cdots\wedge\mathrm d\check x_n),
\end{aligned}
\]
而
\[
\begin{aligned}
&\mu\!\left(\frac{\partial}{\partial x_j}\cdot
(\mathrm dy_I\wedge\mathrm dx_j\wedge\mathrm dx_J)\right)\\
&\quad=\mu\bigl((-1)^p\mathrm dy_I\wedge\mathrm dx_J\bigr)\\
&\quad=(-1)^{(q-1)(q-2)/2+n(n-1)/2}\mathrm dy_I\wedge
\iota\!\left(\frac{\partial}{\partial\check x_J}\right)
(\mathrm d\check x_1\wedge\cdots\wedge\mathrm d\check x_n).
\end{aligned}
\]

\SourceParagraph{ch06-6.2.5-027}{463}
同样，若 $j\in J$，则
\[
\begin{aligned}
\widetilde\mu(\mathrm dx_j)\cdot\mu(\mathrm dy_I\wedge\mathrm dx_J)
&=\frac{\partial}{\partial\check x_j}\cdot
\left[\pm\mathrm dy_I\wedge
\iota\!\left(\frac{\partial}{\partial\check x_J}\right)
(\mathrm d\check x_1\wedge\cdots\wedge\mathrm d\check x_n)\right]\\
&=0=\mu(\mathrm dx_j\cdot\mathrm dy_I\wedge\mathrm dx_J).
\end{aligned}
\]
另一方面，若 $j\notin J$，则
\[
\begin{aligned}
&\widetilde\mu(\mathrm dx_j)\cdot\mu(\mathrm dy_I\wedge\mathrm dx_J)\\
&\quad=\frac{\partial}{\partial\check x_j}\cdot
\left[(-1)^{p+q(q-1)/2+n(n-1)/2}\mathrm dy_I\wedge
\iota\!\left(\frac{\partial}{\partial\check x_J}\right)
(\mathrm d\check x_1\wedge\cdots\wedge\mathrm d\check x_n)\right]\\
&\quad=(-1)^{q(q-1)/2+n(n-1)/2}\mathrm dy_I\wedge
\iota\!\left(\frac{\partial}{\partial\check x_J}\wedge
\frac{\partial}{\partial\check x_j}\right)
(\mathrm d\check x_1\wedge\cdots\wedge\mathrm d\check x_n),
\end{aligned}
\]
而
\[
\begin{aligned}
&\mu\bigl((\mathrm dx_j)\cdot(\mathrm dy_I\wedge\mathrm dx_J)\bigr)\\
&\quad=\mu\bigl((-1)^{p+q}\mathrm dy_I\wedge\mathrm dx_J\wedge\mathrm dx_j\bigr)\\
&\quad=(-1)^{p+q(q+1)/2+n(n-1)/2}(-1)^{p+q}\mathrm dy_I\wedge
\iota\!\left(\frac{\partial}{\partial\check x_J}\wedge
\frac{\partial}{\partial\check x_j}\right)
(\mathrm d\check x_1\wedge\cdots\wedge\mathrm d\check x_n)\\
&\quad=(-1)^{q(q-1)/2+n(n-1)/2}\mathrm dy_I\wedge
\iota\!\left(\frac{\partial}{\partial\check x_J}\wedge
\frac{\partial}{\partial\check x_j}\right)
(\mathrm d\check x_1\wedge\cdots\wedge\mathrm d\check x_n).
\end{aligned}
\]
这说明，$\widetilde\mu$ 与 $\mu$ 保持 Clifford 乘法。另一方面，$\widetilde\mu$ 显然保持 $T\oplus T^*$ 的半平坦截面上的不定内积。

\SourceParagraph{ch06-6.2.5-028}{463}
由于 $\varphi$ 半平坦，$E_\varphi\subseteq(TX(B)\oplus T^*X(B))\otimes\mathbb C$ 由半平坦截面生成。因此，$E_\varphi$ 为极大各向同性子丛，当且仅当 $E_{\mu(\varphi)}$ 为极大各向同性子丛。由\BookRef{exer:6.19}{练习6.19}，$\varphi$ 为闭形式，当且仅当 $\mu(\varphi)$ 为闭形式。

\SourceParagraph{ch06-6.2.5-029}{463}
最后，计算半平坦旋量空间上的配对在 $\mu$ 下的行为。令 $\mathrm dy_I\wedge\mathrm dx_J$ 具有 $(p,q)$ 型，$\mathrm dy_{I'}\wedge\mathrm dx_{J'}$ 具有 $(n-p,n-q)$ 型。那么，$\langle\mathrm dy_I\wedge\mathrm dx_J,\mathrm dy_{I'}\wedge\mathrm dx_{J'}\rangle=0$，除非 $I'=\{1,\ldots,n\}\setminus I$ 且 $J'=\{1,\ldots,n\}\setminus J$。若是这种情形，则
\[
\begin{aligned}
&\langle\mathrm dy_I\wedge\mathrm dx_J,\mathrm dy_{I'}\wedge\mathrm dx_{J'}\rangle\\
&\quad=(-1)^{(p+q)(p+q-1)/2}\mathrm dy_I\wedge\mathrm dx_J
\wedge\mathrm dy_{I'}\wedge\mathrm dx_{J'}\\
&\quad=(-1)^{(p+q)(p+q-1)/2+q(n-p)}\mathrm dy_I\wedge\mathrm dy_{I'}
\wedge\mathrm dx_J\wedge\mathrm dx_{J'},
\end{aligned}
\]
而
\[
\begin{aligned}
&\langle\mu(\mathrm dy_I\wedge\mathrm dx_J),
\mu(\mathrm dy_{I'}\wedge\mathrm dx_{J'})\rangle\\
&\quad=(-1)^{n+q(q-1)/2+(n-q)(n-q-1)/2}\\
&\qquad\cdot\left\langle
\mathrm dy_I\wedge(-1)^{q(q+1)/2+\sum_{j\in J}j}\mathrm d\check x_{J'},\right.\\
&\hspace{35mm}\left.
\mathrm dy_{I'}\wedge(-1)^{(n-q)(n-q+1)/2+\sum_{j\in J'}j}\mathrm d\check x_J
\right\rangle\\
&\quad=(-1)^{n+n(n-1)/2+(p+n-q)(p+n-q-1)/2}
\mathrm dy_I\wedge\mathrm d\check x_{J'}\wedge\mathrm dy_{I'}\wedge\mathrm d\check x_J\\
&\quad=(-1)^{n+n(n-1)/2+(p+n-q)(p+n-q-1)/2+(n-q)(n-p+q)}\\
&\hspace{25mm}\cdot\mathrm dy_I\wedge\mathrm dy_{I'}
\wedge\mathrm d\check x_J\wedge\mathrm d\check x_{J'}.
\end{aligned}
\]
可以证明，在所采用的符号约定下，两个表达式前面的符号相同。当然，得到的是两个不同流形上的 $2n$ 形式。要比较它们，需要分别选择 $X(B)$ 与 $\check X(B)$ 上处处非零的 $2n$ 形式，将配对化为数值。不过，有典范选择：分别取
\[
\mathrm dy_1\wedge\cdots\wedge\mathrm dy_n\wedge\mathrm dx_1\wedge\cdots\wedge\mathrm dx_n
\]
与
\[
\mathrm dy_1\wedge\cdots\wedge\mathrm dy_n\wedge\mathrm d\check x_1\wedge\cdots\wedge\mathrm d\check x_n.
\]
这些形式分别与 $\Omega\wedge\overline\Omega$ 和 $\omega^n$ 成比例，其中，$\Omega$ 为 $X(B)$ 上的典范全纯 $n$ 形式，$\omega$ 为 $\check X(B)$ 上的典范辛形式。无论如何，对于半平坦旋量 $\varphi,\psi$，得到
\begin{equation}\label{eq:6.2}
\langle\varphi,\psi\rangle=\langle\mu(\varphi),\mu(\psi)\rangle.
\end{equation}
因此，若 $\langle\varphi,\overline\varphi\rangle\ne0$，则 $\langle\mu(\varphi),\mu(\overline\varphi)\rangle\ne0$。
\end{proof}

\SourceParagraph{ch06-6.2.5-030}{464}
当然，常常关心的是带有度量的 Calabi--Yau 流形。因此，一个自然的问题是：Kähler 流形，或者广义 Calabi--Yau 流形上的 Ricci 平坦度量，应当如何推广？Gualtieri 的学位论文\BookCite{ref212}给出了一个解答。

\SourceParagraph{ch06-6.2.5-031}{464}
\BookStatementNumber{40}
\begin{definition}\label{def:6.40}
实流形 $M$ 上的广义 Kähler 结构，是一对彼此对易的广义复结构 $(\mathcal J_1,\mathcal J_2)$，使得 $G=-\mathcal J_1\mathcal J_2$ 通过
\[
G(X+\xi,Y+\eta):=\langle G(X+\xi),Y+\eta\rangle
\]
定义 $TM\oplus T^*M$ 上的正定度量。
\end{definition}

\SourceParagraph{ch06-6.2.5-032}{464}
\BookStatementNumber{41}
\begin{example}\label{ex:6.41}
令 $(g,J,\omega)$ 为流形 $M$ 上通常的 Kähler 结构，其中，$g$ 为 Riemann 度量，$J$ 为可积概复结构，$\omega$ 为 Kähler 形式，满足 $\omega(\cdot,J\cdot)=g(\cdot,\cdot)$。那么，$J$ 与 $\omega$ 都定义广义复结构：
\[
\mathcal J_1=\begin{pmatrix}-J&0\\0&J^*\end{pmatrix},
\qquad
\mathcal J_2=\begin{pmatrix}0&\omega^{-1}\\-\omega&0\end{pmatrix}.
\]
注意，
\[
G=-\mathcal J_1\mathcal J_2=-\mathcal J_2\mathcal J_1
=\begin{pmatrix}0&g^{-1}\\g&0\end{pmatrix},
\]
其中，$g\colon TM\to T^*M$ 为由度量 $g$ 诱导的识别。所以，
\[
\begin{aligned}
G(X+\xi,Y+\eta)
&=\langle g^{-1}(\xi)+g(X),Y+\eta\rangle\\
&=\tfrac12\bigl(g^{-1}(\xi,\eta)+g(X,Y)\bigr).
\end{aligned}
\]
\end{example}

\SourceParagraph{ch06-6.2.5-033}{464}
\BookStatementNumber{42}
\begin{exercise}\label{exer:6.42}
通过同时扭曲 $\mathcal J_1$ 与 $\mathcal J_2$，可以用 B 场扭曲任意广义 Kähler 结构。对于上面讨论的通常 Kähler 结构，写出用 B 场作扭曲的公式。
\end{exercise}

\SourceParagraph{ch06-6.2.5-034}{464}
现在，可以定义广义 Calabi--Yau 度量。

\SourceParagraph{ch06-6.2.5-035}{464}
\BookStatementNumber{43}
\begin{definition}\label{def:6.43}
广义 Calabi--Yau 度量，由广义 Kähler 结构 $(\mathcal J_1,\mathcal J_2)$ 定义，其中，这两个结构分别来自定义广义 Calabi--Yau 结构的闭纯旋量 $\varphi_1,\varphi_2$，并满足
\[
\langle\varphi_1,\overline\varphi_1\rangle
=C\langle\varphi_2,\overline\varphi_2\rangle,
\]
$C$ 为某个非零常数。
\end{definition}

\SourceParagraph{ch06-6.2.5-036}{464}
\BookStatementNumber{44}
\begin{example}\label{ex:6.44}
在\BookRef{ex:6.41}{例题6.41}中，若流形 $M$ 为 Calabi--Yau 流形，并选择相容的全纯体积形式 $\Omega$，使其相对于 Ricci 平坦度量平行，则体积归一化要求 $\omega^n$ 为 $\Omega\wedge\overline\Omega$ 的常数倍。这正是\BookRef{def:6.43}{定义6.43}的条件。按照所用的 Clifford 约定，两个纯旋量可以取为 $\Omega$ 与 $\exp(-\mathrm i\omega)$。
\end{example}

\SourceParagraph{ch06-6.2.5-037}{464}
现在，镜像对称作用于广义 Calabi--Yau 度量。

\SourceParagraph{ch06-6.2.5-038}{464}
\BookStatementNumber{45}
\begin{theorem}\label{thm:6.45}
令 $\varphi_1,\varphi_2$ 为纯半平坦旋量，确定 $X(B)$ 上的广义 Calabi--Yau 度量。则 $\mu(\varphi_1),\mu(\varphi_2)$ 确定 $\check X(B)$ 上的广义 Calabi--Yau 度量。
\end{theorem}
\begin{proof}
\SourceParagraph{ch06-6.2.5-039}{465}
由\BookRef{thm:6.39}{定理6.39}，$\mu(\varphi_1)$ 与 $\mu(\varphi_2)$ 都定义 $\check X(B)$ 上的广义 Calabi--Yau 结构。此外，由\BookEqRef{eq:6.2}{6.2}，条件
\[
\langle\mu(\varphi_1),\mu(\overline\varphi_1)\rangle
=C\langle\mu(\varphi_2),\mu(\overline\varphi_2)\rangle
\]
成立，因此，只需证明 $\mu(\varphi_1)$ 与 $\mu(\varphi_2)$ 定义广义 Kähler 结构。用 $\check{\mathcal J}_1,\check{\mathcal J}_2$ 表示 $\mu(\varphi_1),\mu(\varphi_2)$ 在 $T\check X(B)\oplus T^*\check X(B)$ 上定义的概复结构。那么，在 $\check S_{1,0}\oplus\check S_{0,1}\oplus\check S^{1,0}\oplus\check S^{0,1}$ 上，
\[
\check{\mathcal J}_i=\widetilde\mu\mathcal J_i\widetilde\mu^{-1},
\]
因为由\BookEqRef{eq:6.1}{6.1}，$E_{\mu(\varphi_i)}$ 的半平坦截面空间，就是 $E_{\varphi_i}$ 的半平坦截面空间在 $\widetilde\mu$ 下的像。于是，它们彼此对易，并且
\[
-\check{\mathcal J}_1\check{\mathcal J}_2
=-\widetilde\mu\mathcal J_1\mathcal J_2\widetilde\mu^{-1}.
\]
所以，若 $-\mathcal J_1\mathcal J_2$ 定义正定度量，则 $-\check{\mathcal J}_1\check{\mathcal J}_2$ 也定义正定度量。
\end{proof}

\SourceParagraph{ch06-6.2.5-040}{465}
\BookStatementNumber{46}
\begin{example}\label{ex:6.46}
设 $\dim_{\mathbb R}M=2$。若 $\varphi$ 为 $M$ 上的奇纯旋量，则广义 Calabi--Yau 条件为 $\mathrm d\varphi=0$，以及 $\varphi\wedge\overline\varphi$ 处处非零。在这种情形中，$\varphi$ 定义通常的复结构。

\SourceParagraph{ch06-6.2.5-041}{465}
若 $\varphi$ 为定义广义 Calabi--Yau 结构的偶旋量，则 $\varphi=c+\beta$，其中，$c$ 为常数，$\beta$ 为2形式。那么，
\[
\langle c+\beta,\bar c+\overline\beta\rangle
=c\overline\beta-\bar c\beta\ne0
\]
蕴含 $c\ne0$、$\operatorname{Im}(\beta/c)\ne0$，所以 $\beta/c=B+\mathrm i\omega$，其中，$B$ 为任意闭实2形式，$\omega$ 为辛形式。因此，
\[
\varphi=c\exp(B+\mathrm i\omega).
\]
\end{example}

\SourceParagraph{ch06-6.2.5-042}{465}
\BookStatementNumber{47}
\begin{exercise}\label{exer:6.47}
证明：若 $\dim_{\mathbb R}M=2$，则 $(\varphi_1,\varphi_2)$ 定义广义 Calabi--Yau 度量，当且仅当一个旋量为偶旋量，另一个为奇旋量，偶旋量形如 $c\exp(B-\mathrm i\omega)$，其中，$\omega$ 为相对于另一个旋量确定的复结构的 Kähler 形式，并且两个旋量的配对 $\langle\varphi_i,\overline\varphi_i\rangle$ 之比为非零常数。
\end{exercise}

\SourceParagraph{ch06-6.2.5-043}{465}
\BookStatementNumber{48}
\begin{example}\label{ex:6.48}
设 $\dim_{\mathbb R}M=4$。首先，设 $\varphi$ 为定义广义 Calabi--Yau 结构的奇纯旋量，$\varphi=\beta+\gamma$，其中，$\beta$ 为1形式，$\gamma$ 为3形式。先考虑 $E_\varphi$ 为极大各向同性子丛这一条件。若 $X+\xi\in E_\varphi$，则
\[
0=(X+\xi)\cdot(\beta+\gamma)
=\iota(X)\beta+(\xi\wedge\beta+\iota(X)\gamma)+\xi\wedge\gamma.
\]
这说明，$X\in\ker\beta$，因此在 $M$ 的每点，$X$ 都限于一个三维复空间。另一方面，$\xi\wedge\beta+\iota(X)\gamma=0$ 蕴含，$\xi$ 由 $\iota(X)\gamma$ 确定，至多相差 $\beta$ 的一个倍数。因此，$E_\varphi$ 要达到四维，就必须只要 $X+\xi\in E_\varphi$，都有 $X+\xi+c\beta\in E_\varphi$。不过，这只有在 $\beta\wedge\gamma=0$ 时才成立，也就是说，对于某个2形式 $\nu$，有 $\gamma=\beta\wedge\nu$。另一方面，若 $\varphi=\beta\wedge(1+\nu)$ 且 $X\in\ker\beta$，则
\[
(X+\xi)\cdot\bigl(\beta\wedge(1+\nu)\bigr)
=\xi\wedge\beta\wedge(1+\nu)-\beta\wedge\iota(X)\nu,
\]
当 $\xi=-\iota(X)\nu+c\beta$ 时，它为零。因此，
\[
E_\varphi=\{X-\iota(X)\nu+c\beta\mid X\in\ker\beta,\ c\in\mathbb C\}
\]
为极大各向同性子丛。现在，
\[
\langle\beta+\gamma,\overline\beta+\overline\gamma\rangle
=\beta\wedge\overline\gamma+\overline\beta\wedge\gamma
=\beta\wedge\overline\beta\wedge(\overline\nu-\nu).
\]
因此，这一配对处处非零与 $\mathrm d\varphi=0$ 的条件说明，局部地存在函数 $f\colon M\to\mathbb C$，满足 $\mathrm df=\beta$，在 $M$ 上局部定义纤维化；$\operatorname{Im}\nu$ 在纤维上定义辛结构，$\operatorname{Re}\nu$ 在纤维上定义 B 场。所以，这些广义复结构，是复结构与辛结构的一种混合。
\end{example}

\SourceParagraph{ch06-6.2.5-044}{466}
\BookStatementNumber{49}
\begin{exercise}\label{exer:6.49}
确定：两个具有上述形式的奇旋量 $\varphi_1,\varphi_2$，何时给出广义 Calabi--Yau 度量。
\end{exercise}

\SourceParagraph{ch06-6.2.5-045}{466}
\begin{example*}
接下来，考虑偶纯旋量，$\varphi=c+\beta+\gamma$，其中，$c$ 为常数，$\beta$ 为2形式，$\gamma$ 为4形式。先假定 $c\ne0$。通过缩放，可以假定 $c=1$。那么，若 $X+\xi\in E_\varphi$，有
\[
0=(\xi+\iota(X)\beta)+(\xi\wedge\beta+\iota(X)\gamma).
\]
所以，$\xi=-\iota(X)\beta$ 由 $X$ 确定，然后有
\[
0=-(\iota(X)\beta)\wedge\beta+\iota(X)\gamma,
\]
或者
\[
0=\iota(X)(-\beta\wedge\beta/2+\gamma).
\]
要使 $E_\varphi$ 极大，这必须对所有 $X$ 成立。因此，必须有
\[
\gamma=\beta\wedge\beta/2,
\qquad
\varphi=\exp(\beta).
\]
此外，
\[
\begin{aligned}
0&\ne\langle1+\beta+\beta\wedge\beta/2,
1+\overline\beta+\overline\beta\wedge\overline\beta/2\rangle\\
&=\frac{\overline\beta\wedge\overline\beta}{2}
-\beta\wedge\overline\beta+\frac{\beta\wedge\beta}{2}
=\tfrac12(\overline\beta-\beta)^2.
\end{aligned}
\]
因此，$\beta=B+\mathrm i\omega$，其中，$\omega$ 为辛形式。所以，$\varphi$ 诱导用 B 场扭曲的辛流形。

\SourceParagraph{ch06-6.2.5-046}{466}
若 $c=0$，则由非零配对可知 $\beta\ne0$，并且
\[
\dim_{\mathbb C}\{X\in T_bM\otimes\mathbb C\mid\iota(X)\beta=0\}\leq2,
\]
等号成立，当且仅当 $\beta$ 可分解，等价地，$\beta\wedge\beta=0$。同样，
\[
\dim_{\mathbb C}\{\xi\in T_b^*M\otimes\mathbb C\mid\xi\wedge\beta=0\}\leq2,
\]
等号成立，当且仅当 $\beta$ 可分解。由于 $(X+\xi)\cdot\varphi=0$ 当且仅当 $\iota(X)\beta=0$ 且 $\iota(X)\gamma=-\xi\wedge\beta$，所以，$E_\varphi$ 为极大各向同性子丛，当且仅当 $\beta\wedge\beta=0$。此外，$\langle\varphi,\overline\varphi\rangle=-\beta\wedge\overline\beta\ne0$，这说明，$\beta$ 是某个复结构的全纯2形式。特别地，可以写成 $\varphi=\beta+\beta\wedge B$，其中，$B$ 为任意2形式，不一定闭。由于相对于 $\beta$ 诱导的复结构，只有 $B$ 的 $(0,2)$ 部分起作用，可以假定对某个函数 $f$，有 $B=f\overline\beta$，从而 $\varphi=\beta+f\beta\wedge\overline\beta$。
\end{example*}

\SourceParagraph{ch06-6.2.5-047}{467}
\BookStatementNumber{50}
\begin{exercise}\label{exer:6.50}
\begin{enumerate}[label=(\arabic*)]
\item 令 $\varphi_1,\varphi_2$ 在 $2n$ 维流形 $M$ 上定义广义 Calabi--Yau 度量。证明
\[
\dim(E_{\varphi_1}\cap E_{\varphi_2})
=\dim(E_{\varphi_1}\cap\overline{E_{\varphi_2}})=n.
\]
因此，特别地，$\langle\varphi_1,\varphi_2\rangle
=\langle\varphi_1,\overline\varphi_2\rangle=0$。提示：证明，$G$ 作用于 $TM\oplus T^*M$ 的 $\pm1$ 特征空间 $G^\pm$ 都为 $2n$ 维。然后，用 $E_{\varphi_i}$ 与 $E_{\overline\varphi_i}$ 表达 $G^\pm\otimes\mathbb C$。
\item 反过来，若 $n=2$，且 $\varphi_1,\varphi_2$ 具有相同的奇偶性，并定义广义 Calabi--Yau 结构，证明：它们共同定义广义 Calabi--Yau 度量，当且仅当
\[
\langle\varphi_1,\varphi_2\rangle
=\langle\varphi_1,\overline\varphi_2\rangle=0,
\qquad
\langle\varphi_1,\overline\varphi_1\rangle
=C\langle\varphi_2,\overline\varphi_2\rangle,
\]
其中，$C>0$ 为常数，并且 $G=-\mathcal J_1\mathcal J_2$ 定义正定度量。提示：利用如下事实，两个纯偶旋量 $\varphi,\psi$ 的 $\dim(E_\varphi\cap E_\psi)$ 为偶数；纯奇旋量的情形同理。
\end{enumerate}
\end{exercise}

\SourceParagraph{ch06-6.2.5-048}{467}
\BookStatementNumber{51}
\begin{example}\label{ex:6.51}
终于可以理解本小节开头问题的答案。若 $B$ 为仿射流形，转移映射属于 $\mathbb R^n\rtimes\operatorname{SL}_n(\mathbb Z)$，现在允许 B 场为任意闭半平坦2形式。那么，$X(B)$ 上由扭曲辛形式 $\exp(B+\mathrm i\omega)$ 给出的广义 Calabi--Yau 结构，在 $\check X(B)$ 上定义镜像的广义 Calabi--Yau 结构 $\mu(\exp(B+\mathrm i\omega))$。\BookRef{prop:6.26}{命题6.26}说明，若 $B$ 与 $\omega$ 为 $(1,1)$ 型半平坦形式，则 $\mu(\exp(B+\mathrm i\omega))$ 在 $\check X(B)$ 上定义通常的复结构。否则，一般只能得到广义复结构。由上述实四维情形的分析，若 $\dim_{\mathbb R}B=2$，则得到的或者是扭曲复结构，或者是扭曲辛结构。事实上，在这个二维底空间的情形中可以看到，若 $B+\mathrm i\omega$ 有非零的 $(0,2)$ 型部分，则 $\mu(\exp(B+\mathrm i\omega))$ 有 $(0,0)$ 型分量，因此，其镜像必须给出扭曲辛结构。否则，镜像为扭曲复结构。当然，在更高维中，情形可能更加复杂。
\end{example}

\section{半平坦情形中的 A 膜与 B 膜}\label{sec:6.3}

\SourceParagraph{ch06-6.3-intro-001}{467}
现在，回到最初启发 SYZ 猜想的 A 膜与 B 膜之间的对偶，如\BookSectionRef{sec:1.1.5}{1.1.5}所讨论的。将注意力限制在一类特别简单的膜上，就会得到半平坦镜像对上 A 膜与 B 膜模空间之间的精确同构。

\subsection{B 膜}\label{sec:6.3.1}

\SourceParagraph{ch06-6.3.1-001}{467}
现在，如\BookSectionRef{sec:6.2.3}{6.2.3}中一样，固定仿射流形 $B$，其转移映射属于 $\mathbb R^n\rtimes\operatorname{GL}(n,\mathbb Z)$，并固定 B 场 $\mathbf B\in H^1(B,\Lambda_{\mathbb R}/\Lambda)$。这产生复流形 $X(B,\mathbf B)$。考虑 $X(B,\mathbf B)$ 上一些相对简单的 B 膜，由二元组 $(Y,\mathcal L)$ 给出。这里，$Y$ 是 $X(B,\mathbf B)$ 的复子流形，使得 $f|_Y\colon Y\to f(Y)\subseteq B$，是 $B$ 的一个 $d$ 维子流形上的 $d$ 维环面纤维化。此外，$\mathcal L$ 是 $Y$ 上某种特定类型的线丛。

\SourceParagraph{ch06-6.3.1-002}{468}
先作局部考察：开集 $U\subseteq B$ 与 $M_{\mathbb R}$ 的一个子集识别，其中，与通常一样，$M=\mathbb Z^n$、$M_{\mathbb R}=M\otimes_{\mathbb Z}\mathbb R$。假定 B 场在这一开集上平凡。令 $X=f^{-1}(U)=U\times(M_{\mathbb R}/M)$，复结构仍由纤维坐标与底坐标组成的 $x+\mathrm i y$ 给出。设 $Y\subseteq X$ 为上述复流形，在 $U$ 的一个子流形上以环面为纤维。考察点 $x\in Y$，$f(x)=b$；切空间 $T_xY$ 必须是 $T_xX$ 的复子空间。因此，可以分裂为 $T_xY=V\oplus JV$，其中，$J$ 是 $X$ 上的概复结构，$V$ 是 $Y$ 的垂直切空间，即纤维 $Y_b=X_b\cap Y$ 在 $x$ 点的切空间。那么，$f_*(JV)$ 是 $f(Y)$ 在 $b$ 点的切空间。另一方面，$f_*\colon J(T_xX_b)\to T_bB$ 为同构，所以，向量空间 $JV$ 完全由 $T_b(f(Y))$ 确定，不依赖于 $x\in Y_b$。因此，$V$ 本身也不依赖于 $x\in Y_b$，所以，$Y_b$ 只是一个线性环面。于是，$Y_b=(V/(V\cap M))+g(b)$，其中，函数 $g\colon f(Y)\to M_{\mathbb R}$ 给出把环面从 $M_{\mathbb R}$ 的原点移开的平移。

\SourceParagraph{ch06-6.3.1-003}{468}
$V$ 还须满足一个约束：它必须是 $M_{\mathbb R}$ 的一个在 $\mathbb Q$ 上定义的子空间，才能使 $V/(V\cap M)$ 为环面。因此，随着 $b\in f(Y)$ 变化，由于 $V\subseteq M_{\mathbb R}$ 必须连续变化，且始终在 $\mathbb Q$ 上定义，$V$ 实际上必须保持不变。所以，$f(Y)$ 局部张成 $M_{\mathbb R}$ 的一个有理仿射子空间。

\SourceParagraph{ch06-6.3.1-004}{468}
注意，还可以限制函数 $g$ 的性质：要使 $Y$ 在点 $x$ 处的切空间形如 $V\oplus JV$，$g$ 必须在模 $V$ 的意义下为常数。这很容易证明。

\SourceParagraph{ch06-6.3.1-005}{468}
\BookStatementNumber{52}
\begin{note}[numbered]\label{note:6.52}
从方程的角度看，这里所考虑的特定类型的复子流形，可以用\BookSectionRef{sec:6.2.1}{6.2.1}中的复坐标，通过单项式方程写出：
\[
a_1\prod_iq_i^{b_{i1}}=\cdots=a_m\prod_iq_i^{b_{im}}=1,
\]
其中，$(b_{ij})$ 为整矩阵。对这些方程施加 $-\frac1{2\pi}\log|\cdot|$，就得到定义这个复子流形在 $U$ 中的像的仿射线性方程。
\end{note}

\SourceParagraph{ch06-6.3.1-006}{468}
现在，把这些观察推广到 $X(B,\mathbf B)$ 的整体情形。

\SourceParagraph{ch06-6.3.1-007}{468}
\BookStatementNumber{53}
\begin{definition}\label{def:6.53}
若 $B$ 为仿射流形，则 $B$ 的仿射子空间，是一个子流形，局部地由某个仿射线性函数 $M_{\mathbb R}\to\mathbb R^m$ 的零点集给出。若 $B$ 的完整表示的像包含于 $M_{\mathbb R}\rtimes\operatorname{GL}(n,\mathbb Z)$，则当这一仿射线性映射的微分在 $\mathbb Q$ 上定义时，称该仿射子空间为有理的。
\end{definition}

\SourceParagraph{ch06-6.3.1-008}{468}
\BookStatementNumber{54}
\begin{proposition}\label{prop:6.54}
令 $B$ 为仿射流形，完整表示的像包含于 $M_{\mathbb R}\rtimes\operatorname{GL}(n,\mathbb Z)$。令 $\mathbf B\in H^1(B,\Lambda_{\mathbb R}/\Lambda)$，由 $B$ 的开覆盖 $\{U_i\}$ 上的 Čech 1上闭链 $(b_{ij})$ 代表。那么，给出全纯子流形 $Y\subseteq X(B,\mathbf B)$，它在 $d$ 维的 $f(Y)\subseteq B$ 上以 $d$ 维环面为纤维，等价于给出以下数据：
\begin{enumerate}[label=(\arabic*)]
\item 一个 $d$ 维有理仿射子流形 $B'\subseteq B$，其中，$\Lambda'\subseteq\Lambda|_{B'}$ 为与 $B'$ 相切的平坦整切向量所构成的局部系统，$\Lambda_{\mathbb R}':=\Lambda'\otimes_{\mathbb Z}\mathbb R$；
\item 元素 $g_i\in\Gamma(U_i\cap B',\Lambda_{\mathbb R}|_{B'})$，满足
\[
g_j-g_i=b_{ij}|_{U_i\cap U_j\cap B'}
\quad\bmod\Gamma(U_i\cap U_j\cap B',\Lambda_{\mathbb R}'+\Lambda|_{B'}).
\]
这里，$g_i$ 只在模去 $\Lambda_{\mathbb R}'+\Lambda|_{B'}$ 的意义下指定同一子流形。
\end{enumerate}
\end{proposition}
\begin{proof}
\SourceParagraph{ch06-6.3.1-009}{469}
设给定这样的数据。记 $g\colon X(B)\to B$ 与 $f\colon X(B,\mathbf B)\to B$。那么，在 $U_i$ 上，自然地有复子流形
\[
X(B'\cap U_i)=T(B'\cap U_i)/\Lambda'|_{B'\cap U_i}
\subseteq g^{-1}(U_i).
\]
然后，在 $g^{-1}(U_i)$ 内用 $g_i$ 作平移，得到复子流形 $g_i+X(B'\cap U_i)$。现在，通过用 $b_{ij}$ 作平移，把 $g^{-1}(U_i)$ 与 $g^{-1}(U_j)$ 粘合，得到 $X(B,\mathbf B)$。考虑点 $x\in g^{-1}(U_i)$，满足 $g(x)=b\in U_i\cap U_j$，且 $x\in g_i(b)+X(B'\cap U_i)$。那么，$x=g_i(b)+x'$，与
\[
x+b_{ij}=g_i(b)+b_{ij}+x'
=g_j(b)+x'\quad\bmod(\Lambda_{\mathbb R,b}'+\Lambda_b)
\]
识别。但是，$x'$ 加上 $\Lambda_{\mathbb R,b}'+\Lambda_b$ 中的任意元素，仍属于 $X(B'\cap U_j)$，所以，$x+b_{ij}\in g_j+X(B'\cap U_j)$。因此，粘合映射在 $U_i\cap U_j$ 上，将 $g_i+X(B'\cap U_i)$ 与 $g_j+X(B'\cap U_j)$ 识别。

\SourceParagraph{ch06-6.3.1-010}{469}
反向结论，由本小节开头的局部讨论得到。
\end{proof}

\SourceParagraph{ch06-6.3.1-011}{469}
\BookStatementNumber{55}
\begin{exercise}\label{exer:6.55}
令 $Y\subseteq X(B,\mathbf B)$ 由数据 $B'\subseteq B$、$(g_i)$ 确定。证明，存在同构 $Y\cong X(B',\mathbf B')$，其中，$\mathbf B'\in H^1(B',\Lambda_{\mathbb R}'/\Lambda')$ 由 Čech 1上闭链 $(b_{ij}+g_i-g_j)_{ij}$ 代表，且
\[
\begin{aligned}
b_{ij}+g_i-g_j
&\in\Gamma(U_i\cap U_j\cap B',\,(
\Lambda|_{B'}+\Lambda_{\mathbb R}')/\Lambda|_{B'})\\
&=\Gamma(U_i\cap U_j\cap B',\,
\Lambda_{\mathbb R}'/(\Lambda|_{B'}\cap\Lambda_{\mathbb R}'))\\
&=\Gamma(U_i\cap U_j\cap B',\,\Lambda_{\mathbb R}'/\Lambda').
\end{aligned}
\]
注意，有正合列
\begin{equation}\label{eq:6.3}
0\longrightarrow\Lambda_{\mathbb R}'/\Lambda'
\longrightarrow(\Lambda_{\mathbb R}/\Lambda)|_{B'}
\longrightarrow\Lambda_{\mathbb R}|_{B'}/(\Lambda|_{B'}+\Lambda_{\mathbb R}')
\longrightarrow0.
\end{equation}
给定 B 场 $\mathbf B=(b_{ij})$，存在截面 $g_i\in\Gamma(U_i\cap B',\Lambda_{\mathbb R}|_{B'})$，使得 $b_{ij}+g_i-g_j\in\Gamma(U_i\cap U_j\cap B',\Lambda|_{B'}+\Lambda_{\mathbb R}')$，等价于
\[
\mathbf B|_{B'}\in H^1(B',(\Lambda_{\mathbb R}/\Lambda)|_{B'})
\]
在 $H^1(B',\Lambda_{\mathbb R}|_{B'}/(\Lambda|_{B'}+\Lambda_{\mathbb R}'))$ 中的像为零。证明：所有在 $B'$ 上以环面为纤维的全纯子流形 $Y\subseteq X(B,\mathbf B)$，与满足 $g_j-g_i=b_{ij}\bmod(\Lambda|_{B'}+\Lambda_{\mathbb R}')$ 的、取值于 $\Lambda_{\mathbb R}|_{B'}/(\Lambda|_{B'}+\Lambda_{\mathbb R}')$ 的 Čech 0上链 $(g_i)$ 一一对应。证明，
\[
H^0(B',\Lambda_{\mathbb R}|_{B'}/(\Lambda|_{B'}+\Lambda_{\mathbb R}'))
\]
自由且传递地作用于这类选择的集合；称这一选择集合为上述群的主齐性空间，即 torsor。\BookQED
\end{exercise}

\SourceParagraph{ch06-6.3.1-012}{469}
\BookStatementNumber{56}
\begin{note}[numbered]\label{note:6.56}
还可以考虑 $Y\to f(Y)$ 的纤维为环面的不交并的情形。为此，允许 $B'\to B$ 是其像上的无分歧覆盖，而非嵌入。
\end{note}

\SourceParagraph{ch06-6.3.1-013}{469}
\BookStatementNumber{57}
\begin{example}\label{ex:6.57}
若 $B=M_{\mathbb R}/\Gamma$，其中，$\Gamma\subseteq M_{\mathbb R}$ 为格，先注意，若 $\Gamma$ 取一般值，则 $B$ 不包含正维的真紧有理仿射子空间。当 $\Gamma$ 没有秩介于0与 $n$ 之间的、张成 $M_{\mathbb R}$ 的有理子空间的子格时，就出现这种情形。不过，即使 $B$ 有这样的紧有理仿射子空间 $B'$，对于一般选择的 $\mathbf B$，它限制到
\[
H^1(B',\Lambda_{\mathbb R}|_{B'}/(\Lambda|_{B'}+\Lambda_{\mathbb R}'))
\]
后也不为零，这很容易证明。这反映了如下事实：一般的复环面不包含非平凡的真复子环面。\BookQED
\end{example}

\SourceParagraph{ch06-6.3.1-014}{470}
作为第一步近似，希望把 B 膜看作不只是一个复子流形，而是一个复子流形连同一个复向量丛。作进一步简化，假定这一向量丛为线丛。由\BookRef{exer:6.55}{练习6.55}，所考虑的复子流形始终形如 $X(B',\mathbf B')$，因此，先描述 $X(B,\mathbf B)$ 上的线丛。

\SourceParagraph{ch06-6.3.1-015}{470}
实际上，只描述 $X(B,\mathbf B)$ 上的某些线丛。设 $B$ 有开覆盖 $\{U_i\}$，带有坐标图 $\psi_i\colon U_i\to M_{\mathbb R}$。可以通过使线丛在 $f^{-1}(U_i)$ 上平凡，并取转移函数 $s_{ij}\colon f^{-1}(U_i\cap U_j)\to\mathbb C^\times$，构造 $X(B,\mathbf B)$ 上的复线丛。换言之，利用映射 $(x,c)\mapsto(x,s_{ij}(x)c)$ 作识别
\[
\begin{aligned}
f^{-1}(U_i)\times\mathbb C
&\supseteq f^{-1}(U_i\cap U_j)\times\mathbb C\\
&\longrightarrow f^{-1}(U_i\cap U_j)\times\mathbb C
\subseteq f^{-1}(U_j)\times\mathbb C.
\end{aligned}
\]
现在，只考虑 $f^{-1}(U_i\cap U_j)$ 上某些特定的可逆函数。若使用坐标图 $\psi_i\colon U_i\to M_{\mathbb R}$，将 $f^{-1}(U_i\cap U_j)$ 与 $\psi_i(U_i\cap U_j)\times(M_{\mathbb R}/M)$ 识别，并通过 $q=\exp(2\pi\mathrm i(x+\mathrm i y))$ 将其嵌入 $M\otimes_{\mathbb Z}\mathbb C^\times$，则对于 $\nu\in N=\operatorname{Hom}_{\mathbb Z}(M,\mathbb Z)$，定义 $q^\nu$ 为由 $\nu$ 确定的特征标 $M\otimes\mathbb C^\times\to\mathbb C^\times$。它是全纯函数。对于 $\nu\in N$，函数 $aq^\nu$，其中，$a\in\mathbb C^\times$，是 $M\otimes\mathbb C^\times$ 上的可逆函数，因而也是 $f^{-1}(U_i\cap U_j)$ 上的可逆函数。当然，还可能有许多其他可逆函数；局部地，$f^{-1}(U_i\cap U_j)$ 上的任意全纯函数，都有由这类项组成的无穷 Laurent 展开。不过，不允许使用这些更复杂的函数。形如 $aq^\nu$ 的函数，称为单项式函数，在 $B$ 上构成层 $\mathcal P_{B,\mathbf B}$，即
\[
\mathcal P_{B,\mathbf B}(U)
=\{f^{-1}(U)\text{ 上的可逆单项式函数}\}.
\]
所关注的线丛群，由上同调群 $H^1(B,\mathcal P_{B,\mathbf B})$ 分类。

\SourceParagraph{ch06-6.3.1-016}{470}
这启发如下定义。

\SourceParagraph{ch06-6.3.1-017}{470}
\BookStatementNumber{58}
\begin{definition}\label{def:6.58}
$X(B,\mathbf B)$ 上的半平坦 B 膜，是 $X(B,\mathbf B)$ 的一个全纯子流形 $Y$，在 $B$ 的有理仿射子流形 $B'$ 上以环面为纤维，并带有 $Y$ 上的线丛 $\mathcal L$，其转移函数为 $Y$ 上的单项式函数。
\end{definition}

\SourceParagraph{ch06-6.3.1-018}{470}
要分类半平坦 B 膜，以更加内在的方式描述 $\mathcal P_{B,\mathbf B}$。有子层 $\mathbb C^\times\subseteq\mathcal P_{B,\mathbf B}$，由常值函数构成。另一方面，若 $b\in B$，则在相差平移的意义下，可以典范地把 $b$ 的某个小邻域 $U$ 与 $\Lambda_{\mathbb R,b}$ 中的邻域 $U'$ 识别。因此，
\[
f^{-1}(U)\cong U'\times(\Lambda_{\mathbb R,b}/\Lambda_b),
\qquad
\mathcal P_{B,\mathbf B}(U)
\cong\mathbb C^\times\times\check\Lambda_b
=\mathbb C^\times\times\check\Lambda(U).
\]
所以，有正合列
\[
1\longrightarrow\mathbb C^\times\longrightarrow\mathcal P_{B,\mathbf B}
\longrightarrow\check\Lambda\longrightarrow0.
\]
这个序列一般不分裂：局部地，$U'$ 的选择影响识别 $\mathcal P_{B,\mathbf B}(U)\cong\mathbb C^\times\times\check\Lambda(U)$，已经清楚地显示了这一点。此外，这一扩张甚至依赖于 B 场。更准确地说，有如下命题。

\SourceParagraph{ch06-6.3.1-019}{471}
\BookStatementNumber{59}
\begin{proposition}\label{prop:6.59}
正合列
\[
1\longrightarrow\mathbb C^\times\longrightarrow\mathcal P_{B,\mathbf B}
\longrightarrow\check\Lambda\longrightarrow0
\]
在 $\operatorname{Ext}_B^1(\check\Lambda,\mathbb C^\times)
\cong H^1(B,\Lambda\otimes_{\mathbb Z}\mathbb C^\times)$ 中的扩张类，由
\[
\exp\bigl(2\pi\mathrm i(\mathbf B+\mathrm i[\delta])\bigr)
\]
代表，其中，$[\delta]$ 为 $B$ 的辐射障碍，见\BookRef{note:6.30}{笔记6.30}。把 $\mathbf B+\mathrm i[\delta]$ 称为 $B$ 的复化辐射障碍。若 B 场有实提升，则它属于 $H^1(B,\Lambda\otimes_{\mathbb Z}\mathbb C)$；一般地，其指数由局部代表的指数定义，因而即使 $\mathbf B$ 只在模 $\Lambda$ 的意义下指定，也仍然良定义。
\end{proposition}
\begin{proof}
\SourceParagraph{ch06-6.3.1-020}{471}
令 $\{U_i\}$ 为开覆盖，每个 $U_i$ 都可缩，令 $\{(U_i\cap U_j,b_{ij})\}$ 为 B 场的 Čech 代表。令 $\psi_i\colon U_i\to M_{\mathbb R}$ 为坐标图。如前所述，这个坐标图给出 $\mathcal P_{B,\mathbf B}|_{U_i}$ 与常值层 $\mathbb C^\times\times N$ 的同构，从而给出局部分裂 $\alpha_i\colon\check\Lambda|_{U_i}\to\mathcal P_{B,\mathbf B}|_{U_i}$。扩张类作为 $H^1(B,\Lambda\otimes\mathbb C^\times)$ 的元素，由这些分裂之差给出的 Čech 1上闭链代表，即 $\{(U_i\cap U_j,\alpha_j/\alpha_i)\}$，其中，$\alpha_j/\alpha_i$ 是映射 $\check\Lambda|_{U_i\cap U_j}\to\mathbb C^\times$，可以看作 $U_i\cap U_j$ 上的 $\Lambda\otimes\mathbb C^\times$ 截面。这里，将群层 $\mathcal P_{B,\mathbf B}$ 写为乘法形式。现在，直接比较 $\alpha_i$ 与 $\alpha_j$。为此，需要理解如何把形如 $X(\psi_i(U_i))$ 的局部片粘合，构成 $X(B,\mathbf B)$。特别地，通过用 $b_{ij}$ 扭曲典范识别，将 $X(\psi_i(U_i\cap U_j))$ 与 $X(\psi_j(U_i\cap U_j))$ 粘合。典范识别如下给出：转移映射 $\psi_j\circ\psi_i^{-1}\in M_{\mathbb R}\rtimes\operatorname{GL}(M)$，在切丛上诱导映射
\[
(\psi_j\circ\psi_i^{-1})_*\colon T(\psi_i(U_i\cap U_j))
\longrightarrow T(\psi_j(U_i\cap U_j)),
\]
进而诱导典范识别 $X(\psi_i(U_i\cap U_j))\to X(\psi_j(U_i\cap U_j))$。再与沿纤维加上 $b_{ij}$ 的映射复合，得到识别
\[
\phi_{ij}\colon X(\psi_i(U_i\cap U_j))
\longrightarrow X(\psi_j(U_i\cap U_j)).
\]
具体地，利用 $X(\psi_i(U_i\cap U_j))$ 与 $X(\psi_j(U_i\cap U_j))$ 在 $M_{\mathbb R}\times(M_{\mathbb R}/M)=M\otimes_{\mathbb Z}\mathbb C^\times$ 中的包含，$\phi_{ij}$ 在 $M_{\mathbb R}$ 上通过 $\psi_j\circ\psi_i^{-1}\in\operatorname{Aff}(M_{\mathbb R})$ 作用，在 $M_{\mathbb R}/M$ 上通过
\[
m\longmapsto\operatorname{Lin}(\psi_j\circ\psi_i^{-1})(m)+b_{ij}
\]
作用；其中，$\operatorname{Lin}\colon\operatorname{Aff}(M_{\mathbb R})\to\operatorname{GL}(M_{\mathbb R})$ 为投影。或者，将元素 $m\in M\otimes_{\mathbb Z}\mathbb C^\times$ 写成 $\exp(2\pi\mathrm i(m_1+\mathrm i m_2))$，其中，$m_1,m_2\in M_{\mathbb R}$，则
\[
\begin{aligned}
\phi_{ij}(m)
&=\phi_{ij}\bigl(\exp(2\pi\mathrm i(m_1+\mathrm i m_2))\bigr)\\
&=\exp\!\left(2\pi\mathrm i\left[
\operatorname{Lin}(\psi_j\circ\psi_i^{-1})(m_1)+b_{ij}
+\mathrm i(\psi_j\circ\psi_i^{-1})(m_2)\right]\right)\\
&=\exp\!\left(2\pi\mathrm i\operatorname{Lin}(\psi_j\circ\psi_i^{-1})
(m_1+\mathrm i m_2)\right)\\
&\qquad\cdot\exp\!\left(2\pi\mathrm i\left[b_{ij}
+\mathrm i\operatorname{Trans}(\psi_j\circ\psi_i^{-1})\right]\right)\\
&=\operatorname{Lin}(\psi_j\circ\psi_i^{-1})(m)
\exp\!\left(2\pi\mathrm i\left[b_{ij}
+\mathrm i\operatorname{Trans}(\psi_j\circ\psi_i^{-1})\right]\right),
\end{aligned}
\]
其中，$\operatorname{Trans}\colon\operatorname{Aff}(M_{\mathbb R})\to M_{\mathbb R}$ 是取平移部分的映射。注意，$\operatorname{Trans}$ 不是同态。

\SourceParagraph{ch06-6.3.1-021}{472}
于是，对于 $\nu\in N$、$m\in M\otimes_{\mathbb Z}\mathbb C^\times$，有
\[
\begin{aligned}
(q^\nu\circ\phi_{ij})(m)
&=q^\nu\!\left(\operatorname{Lin}(\psi_j\circ\psi_i^{-1})(m)
\exp\!\left(2\pi\mathrm i[b_{ij}
+\mathrm i\operatorname{Trans}(\psi_j\circ\psi_i^{-1})]\right)\right)\\
&=q^{{}^t\operatorname{Lin}(\psi_j\circ\psi_i^{-1})(\nu)}(m)
q^\nu\!\left(\exp\!\left(2\pi\mathrm i[b_{ij}
+\mathrm i\operatorname{Trans}(\psi_j\circ\psi_i^{-1})]\right)\right)\\
&=q^{{}^t\operatorname{Lin}(\psi_j\circ\psi_i^{-1})(\nu)}(m)
\exp\!\left(2\pi\mathrm i\left\langle\nu,
b_{ij}+\mathrm i\operatorname{Trans}(\psi_j\circ\psi_i^{-1})\right\rangle\right).
\end{aligned}
\]
因此，
\[
q^\nu\circ\phi_{ij}
=\exp\!\left(2\pi\mathrm i\left\langle\nu,
b_{ij}+\mathrm i\operatorname{Trans}(\psi_j\circ\psi_i^{-1})\right\rangle\right)
q^{{}^t\operatorname{Lin}(\psi_j\circ\psi_i^{-1})(\nu)}.
\]
所以，看到 $\alpha_j/\alpha_i=\exp(2\pi\mathrm i[b_{ij}+\mathrm i\operatorname{Trans}(\psi_j\circ\psi_i^{-1})])$，是 $U_i\cap U_j$ 上的 $\Lambda\otimes\mathbb C^\times$ 截面。结论便由\BookRef{constr:6.29}{构造6.29}与\BookRef{note:6.30}{笔记6.30}中对辐射障碍的描述得到。
\end{proof}

\SourceParagraph{ch06-6.3.1-022}{472}
一个标准事实是：上同调长正合列中的连接同态，由与扩张类的杯积给出。因此，得到如下命题。

\SourceParagraph{ch06-6.3.1-023}{472}
\BookStatementNumber{60}
\begin{proposition}\label{prop:6.60}
在长正合列
\[
H^1(B,\mathbb C^\times)\longrightarrow H^1(B,\mathcal P_{B,\mathbf B})
\longrightarrow H^1(B,\check\Lambda)\longrightarrow H^2(B,\mathbb C^\times)
\]
中，元素 $\overline\alpha\in H^1(B,\check\Lambda)$ 提升为元素 $\alpha\in H^1(B,\mathcal P_{B,\mathbf B})$，当且仅当相对于自然杯积
\[
H^1(B,\check\Lambda)\times H^1(B,\Lambda\otimes\mathbb C^\times)
\longrightarrow H^2(B,\mathbb C^\times),
\]
有
\[
\overline\alpha\cdot\exp\bigl(2\pi\mathrm i(\mathbf B+\mathrm i[\delta])\bigr)=1.
\]
\end{proposition}

\SourceParagraph{ch06-6.3.1-024}{472}
\BookStatementNumber{61}
\begin{example}\label{ex:6.61}
令 $B=M_{\mathbb R}/\Gamma$，其中，$\Gamma\subseteq M_{\mathbb R}$ 为格，令 $\check\Gamma\subseteq N_{\mathbb R}$ 为对偶格。那么，$\Lambda$ 与 $\check\Lambda$ 都是常值局部系统。正合列
\[
1\longrightarrow\mathbb C^\times\longrightarrow\mathcal P_{B,\mathbf B}
\longrightarrow\check\Lambda\longrightarrow0
\]
给出上同调长正合列
\[
\begin{aligned}
1\longrightarrow\mathbb C^\times
&\longrightarrow H^0(B,\mathcal P_{B,\mathbf B})
\longrightarrow N\longrightarrow\check\Gamma\otimes_{\mathbb Z}\mathbb C^\times\\
&\longrightarrow H^1(B,\mathcal P_{B,\mathbf B})
\longrightarrow\check\Gamma\otimes_{\mathbb Z}N
\longrightarrow\bigwedge^2\check\Gamma\otimes_{\mathbb Z}\mathbb C^\times.
\end{aligned}
\]
连接映射 $\check\Gamma\otimes_{\mathbb Z}N\to\bigwedge^2\check\Gamma\otimes_{\mathbb Z}\mathbb C^\times$，由与
\[
\begin{aligned}
\exp\bigl(2\pi\mathrm i(\mathbf B+\mathrm i[\delta])\bigr)
&\in H^1(B,\Lambda\otimes\mathbb C^\times)\\
&=\check\Gamma\otimes_{\mathbb Z}M\otimes_{\mathbb Z}\mathbb C^\times
\end{aligned}
\]
的杯积给出。B 场可以在 $H^1(B,\Lambda_{\mathbb R})=\check\Gamma\otimes_{\mathbb Z}M_{\mathbb R}$ 中取任意值；辐射障碍由典范包含 $\Gamma\to M_{\mathbb R}$ 给出，将其看作 $\operatorname{Hom}(\Gamma,M_{\mathbb R})=\check\Gamma\otimes_{\mathbb Z}M_{\mathbb R}$ 的元素。将这一点作为练习。因此，当 $n\geq2$ 时，看到对于一般选择的 B 场，映射 $H^1(B,\mathcal P_{B,\mathbf B})\to H^1(B,\check\Lambda)$ 实际上为零。

\SourceParagraph{ch06-6.3.1-025}{473}
还注意，$H^0(B,\mathcal P_{B,\mathbf B})=\mathbb C^\times$，因为紧空间 $X(B,\mathbf B)$ 上的整体全纯函数只有常值函数。所以，无论如何，都得到正合列
\[
1\longrightarrow\frac{\check\Gamma\otimes_{\mathbb Z}\mathbb C^\times}{N}
\longrightarrow H^1(B,\mathcal P_{B,\mathbf B})
\longrightarrow\check\Gamma\otimes_{\mathbb Z}N,
\]
当 $n\geq2$ 时，对于一般选择的 $\mathbf B+\mathrm i[\delta]$，其中的第一个映射为同构。
\end{example}

\SourceParagraph{ch06-6.3.1-026}{473}
\BookStatementNumber{62}
\begin{exercise}\label{exer:6.62}
把商 $(\check\Gamma\otimes_{\mathbb Z}\mathbb C^\times)/N$，描述为某个仿射环面 $B'$ 与 B 场 $\mathbf B'$ 所给出的、形如 $X(B',\mathbf B')$ 的复流形。特别注意，$N$ 到 $\check\Gamma\otimes_{\mathbb Z}\mathbb C^\times$ 中的包含，依赖于最初的数据选择。
\end{exercise}

\SourceParagraph{ch06-6.3.1-027}{473}
把上述内容合在一起，得到如下半平坦 B 膜分类。

\SourceParagraph{ch06-6.3.1-028}{473}
\BookStatementNumber{63}
\begin{theorem}\label{thm:6.63}
设给定 $B$ 的有理仿射子流形 $B'$，以及类 $[\mathcal L]\in H^1(B',\check\Lambda')$，其中，$\check\Lambda'$ 为 $\Lambda'$ 的对偶。那么，存在在 $B'$ 上纤维化的半平坦 B 膜 $(Y,\mathcal L)$，其中，$Y$ 同构于某个 $\mathbf B'\in H^1(B',\Lambda_{\mathbb R}'/\Lambda')$ 所给出的 $X(B',\mathbf B')$，线丛 $\mathcal L\in H^1(B',\mathcal P_{B',\mathbf B'})$ 映到 $[\mathcal L]\in H^1(B',\check\Lambda')$，当且仅当可以选择 $\mathbf B'$ 满足：
\begin{enumerate}[label=(\arabic*)]
\item $\mathbf B'$ 是通过正合列\BookEqRef{eq:6.3}{6.3}，把 $\mathbf B|_{B'}$ 提升到 $H^1(B',\Lambda_{\mathbb R}'/\Lambda')$ 所得的提升；
\item 相对于自然杯积
\[
H^1(B',\check\Lambda')\times H^1(B',\Lambda'\otimes\mathbb C^\times)
\longrightarrow H^2(B',\mathbb C^\times),
\]
有
\[
[\mathcal L]\cdot\exp\bigl(2\pi\mathrm i(\mathbf B'+\mathrm i[\delta_{B'}])\bigr)=1.
\]
这里，$[\delta_{B'}]$ 表示 $B'$ 上的辐射障碍类。
\end{enumerate}
此外，$Y$ 由取值于 $\Lambda_{\mathbb R}|_{B'}/(\Lambda|_{B'}+\Lambda_{\mathbb R}')$ 的 Čech 0上链 $(g_i)$ 唯一指定；对任意将 $g_i$ 提升到 $(\Lambda_{\mathbb R}/\Lambda)|_{B'}$ 的选择 $\widetilde g_i$，要求 $(b_{ij}+\widetilde g_i-\widetilde g_j)_{ij}$ 代表 $\mathbf B'$。对于给定的 $\mathbf B'$，$Y$ 的选择集合，是群
\[
\operatorname{im}\!\left(
H^0(B',(\Lambda_{\mathbb R}/\Lambda)|_{B'})
\longrightarrow H^0(B',\Lambda_{\mathbb R}|_{B'}/(\Lambda|_{B'}+\Lambda_{\mathbb R}'))
\right)
\]
的主齐性空间。此外，$\mathcal L$ 在相差与下列群中的线丛作张量积的意义下唯一确定：
\[
\frac{H^1(B',\mathbb C^\times)}
{\operatorname{im}\!\left(H^0(B',\check\Lambda')
\longrightarrow H^1(B',\mathbb C^\times)\right)}
\subseteq H^1(B',\mathcal P_{B',\mathbf B'}).
\]
分母中的映射，是相应长正合列的连接同态。
\end{theorem}
\begin{proof}
\SourceParagraph{ch06-6.3.1-029}{474}
将\BookRef{prop:6.54}{命题6.54}、\BookRef{exer:6.55}{练习6.55}与\BookRef{prop:6.60}{命题6.60}结合即可。
\end{proof}

\SourceParagraph{ch06-6.3.1-030}{474}
还可以如下构造更广泛的一类向量丛。令 $B'\to B$ 为有限覆盖映射。那么，$B$ 上的仿射结构在 $B'$ 上诱导仿射结构；把 B 场 $\mathbf B$ 拉回为 $\mathbf B'\in H^1(B',\Lambda_{\mathbb R}'/\Lambda')$，得到 $X(B',\mathbf B')$ 与自然映射 $\pi\colon X(B',\mathbf B')\to X(B,\mathbf B)$，后者为复流形的映射。给定通过前面构造得到的、$X(B',\mathbf B')$ 上的线丛 $\mathcal L$，则 $\pi_*\mathcal L$ 是 $X(B,\mathbf B)$ 上的向量丛，其秩等于映射 $\pi$ 的次数，即对于任意 $x\in X(B,\mathbf B)$，$\pi^{-1}(x)$ 的点数。此外，可以通过对已经构造的向量丛作扩张，构造新的向量丛。用这种方式，可以构造椭圆曲线上的全部向量丛，参见例如文献\BookCite{ref398}。不过，为使下一小节的讨论足够简单，这里将继续限于线丛。

\subsection{A 膜}\label{sec:6.3.2}

\SourceParagraph{ch06-6.3.2-001}{474}
接下来，考虑 A 膜。同样，只考虑一类受限制的膜：首先，只考虑 Lagrange 膜，而不考虑余各向同性膜；其次，不要求膜为特殊 Lagrange 膜。事实上，在前一小节中，固定了 $X(B,\mathbf B)$ 上的复结构，并考虑带有线丛的全纯子流形；忽略了膜的条件中来自复化 Kähler 形式选择的那一部分。同样，在这一侧，将研究辛流形 $\check X(B)$，连同 B 场 $\mathbf B$ 的选择。对于二元组 $(L,\mathcal L)$，其中，$L\subseteq\check X(B)$ 为 Lagrange 子流形，$\mathcal L$ 为 $L$ 上带有联络 $\nabla_{\mathcal L}$ 的 $\operatorname{U}(1)$ 丛，希望知道，它成为 A 膜所需的条件是什么。若没有 B 场，只需要求联络平坦，即联络的曲率2形式 $F$ 为零。不过，\BookSectionRef{sec:3.5.2.7}{3.5.2.7}的讨论说明，存在 B 场时，B 场与归一化实曲率的相应组合必须为零；用符号相反的实曲率约定，可以写成 $\mathbf B|_L+F_{\mathrm{real}}=0$。本章将联络1形式取值于 $\operatorname{U}(1)$ 的 Lie 代数，并采用 $c_1(\mathcal L)=[\mathrm iF/(2\pi)]$ 的约定；相应的曲率条件为
\[
F=-2\pi\mathrm i\,\mathbf B|_L.
\]
这里涉及曲率的符号与归一化约定的改变，并且已经用2形式代表 B 场。若这一条件成立，则 $\mathbf B|_L$ 代表 $L$ 上的整上同调类。因此，这是 Lagrange 闭链 $L$ 能够承载 A 膜结构的必要条件。若 $\mathbf B|_L$ 为整类，则存在曲率为 $-2\pi\mathrm i\,\mathbf B|_L$ 的 $\operatorname{U}(1)$ 丛。模去规范等价后，固定其拓扑类型的连续选择由 $H^1(L,\mathbb R)/\operatorname{im}H^1(L,\mathbb Z)$ 给出；若允许所有具有这一曲率的丛及联络，包括挠的平坦丛选择，则选择集合是 $H^1(L,\mathbb R/\mathbb Z)$ 的主齐性空间。

\SourceParagraph{ch06-6.3.2-002}{475}
\BookStatementNumber{64}
\begin{definition}\label{def:6.64}
令 $B'\subseteq B$ 为有理仿射子空间。定义 $L(B')\subseteq\check X(B)$ 为
\[
L(B'):=(TB')^\perp/(\check\Lambda|_{B'}\cap(TB')^\perp),
\]
其中，$(TB')^\perp$ 为 $(T^*B)|_{B'}$ 中湮灭 $TB'\subseteq TB|_{B'}$ 的子丛。

\SourceParagraph{ch06-6.3.2-003}{475}
给定 $B'$ 的开覆盖 $\{U_i\}$，以及 $T^*B|_{U_i}$ 的截面 $g_i$，满足 $g_j-g_i\in\Gamma(U_i\cap U_j,\check\Lambda|_{B'}+(TB')^\perp)$，定义 $L(B',\{(U_i,g_i)\})\subseteq\check X(B)$ 为
\[
L(B',\{(U_i,g_i)\}):=\bigcup_i\bigl(L(U_i)+g_i\bigr).
\]
对 $g_j-g_i$ 的条件，保证 $L(U_i)+g_i$ 与 $L(U_j)+g_j$ 在重叠部分相同。

\SourceParagraph{ch06-6.3.2-004}{475}
注意，对 $\{(U_i,g_i)\}$ 的条件说明，这些数据实际上确定截面
\[
\begin{aligned}
g&\in\Gamma(B',T^*B|_{B'}/(\check\Lambda|_{B'}+(TB')^\perp))\\
&\cong\Gamma(B',T^*B'/\check\Lambda'),
\end{aligned}
\]
其中，$\check\Lambda'$ 为 $T^*B'$ 的平坦整截面所构成的子层。因此，可以写成
\[
L(B',g):=L(B',\{(U_i,g_i)\}).
\]
通过正合列
\[
0\longrightarrow\check\Lambda'\longrightarrow T^*B'
\longrightarrow T^*B'/\check\Lambda'\longrightarrow0,
\]
得到 $g$ 的类，即 $g$ 在连接同态下于 $H^1(B',\check\Lambda')$ 中的像 $[g]$。\BookQED
\end{definition}

\SourceParagraph{ch06-6.3.2-005}{475}
一般而言，对于给定的类 $[g]$，使得 $L(B',g)$ 为 Lagrange 子流形的 $g$ 的选择太多。对于这里的目的，Lagrange 子流形不够刚性，一般具有无穷维模空间。\BookSectionRef{sec:1.1.4}{1.1.4}中 A 膜的完整定义，通过要求它为特殊 Lagrange 膜，使 A 膜具有刚性。不过，并未在 $\check X(B)$ 上指定复结构，因此，这里将采用 Floer 理论中的另一种标准方法：在 Hamilton 同痕的意义下考虑 Lagrange 子流形。

\SourceParagraph{ch06-6.3.2-006}{475}
\BookStatementNumber{65}
\begin{exercise}\label{exer:6.65}
给定 $\check f\colon\check X(B)\to B$ 与 $B$ 上的 $C^\infty$ 函数 $\varphi$，证明：由方程 $\iota(H_\varphi)\omega=\mathrm d(\varphi\circ\check f)$ 定义的 Hamilton 向量场 $H_\varphi$，其时间为1的 Hamilton 流，由1形式 $\mathrm d\varphi$ 给出的沿纤维平移构成。
\end{exercise}

\SourceParagraph{ch06-6.3.2-007}{475}
这一练习启发以下等价关系：若 $g-g'$ 是 $B'$ 上某个恰当1形式的像，则将 $L(B',g)$ 与 $L(B',g')$ 看作等价。

\SourceParagraph{ch06-6.3.2-008}{475}
仿照\BookRef{def:6.58}{定义6.58}，作如下定义。

\SourceParagraph{ch06-6.3.2-009}{476}
\BookStatementNumber{66}
\begin{definition}\label{def:6.66}
$\check X(B)$ 上的半平坦 A 膜，是三元组 $(L,\mathcal L,\nabla_{\mathcal L})$：$L$ 为 $\check X(B)$ 的 Lagrange 子流形，在 $B$ 的有理仿射子流形 $B'$ 上以线性环面为纤维；$\mathcal L$ 为 $L$ 上的 $\operatorname{U}(1)$ 丛，带有联络 $\nabla_{\mathcal L}$，其曲率2形式满足 $F=-2\pi\mathrm i\,\mathbf B|_L$。

\SourceParagraph{ch06-6.3.2-010}{476}
称 $(L_1,\mathcal L_1,\nabla_{\mathcal L_1})$ 与 $(L_2,\mathcal L_2,\nabla_{\mathcal L_2})$ 等价，是指存在 $B'$ 上的恰当1形式 $\alpha$，使得 $T_\alpha(L_1)=L_2$，其中，$T_\alpha$ 表示沿纤维平移 $\alpha$；并且 $T_\alpha^*\mathcal L_2\cong\mathcal L_1$，$\nabla_{\mathcal L_2}$ 的拉回与 $\nabla_{\mathcal L_1}$ 规范等价，即二者相差某个函数 $u\colon L_1\to\operatorname{U}(1)$ 给出的规范项 $u^{-1}\mathrm du$。
\end{definition}

\SourceParagraph{ch06-6.3.2-011}{476}
现在，希望在等价意义下分类半平坦 A 膜。

\SourceParagraph{ch06-6.3.2-012}{476}
\BookStatementNumber{67}
\begin{lemma}\label{lem:6.67}
令 $B'\subseteq B$ 为有理仿射子空间，由 $B$ 的坐标图 $\{U_i\}$ 覆盖。设 $y_1,\ldots,y_n$ 为 $U_i$ 上的仿射坐标，选择为使 $B'$ 由 $y_{p+1}=\cdots=y_n=0$ 给出。那么，$L(B',g)$ 上有坐标 $(y_1,\ldots,y_p,\check x_{p+1},\ldots,\check x_n)$，$\check X(B)$ 上有坐标 $(y_1,\ldots,y_n,\check x_1,\ldots,\check x_n)$。令 $g\in\Gamma(B',T^*B'/\check\Lambda')$，局部地由 $T^*B|_{B'}$ 的截面代表，在这些坐标中写成 $g=\sum_{j=1}^ng^j\,\mathrm dy_j$。那么，可以按如下方式参数化 $L(B',g)$：
\[
\begin{aligned}
&(y_1,\ldots,y_p,\check x_{p+1},\ldots,\check x_n)\\
&\quad\longmapsto(y_1,\ldots,y_p,0,\ldots,0,
g^1,\ldots,g^p,\check x_{p+1}+g^{p+1},\ldots,\check x_n+g^n).
\end{aligned}
\]
设 $\mathbf B=\sum_{i,j}b_{ij}\,\mathrm d\check x_j\wedge\mathrm dy_i$。则
\[
\begin{aligned}
(\mathbf B+\mathrm i\omega)|_{L(B',g)}
&=\sum_{i=1}^p\sum_{j=p+1}^nb_{ij}\,\mathrm d\check x_j\wedge\mathrm dy_i\\
&\qquad+\sum_{i=1}^p\sum_{j=1}^n(b_{ij}+\mathrm i\delta_{ij})\,
\mathrm dg^j\wedge\mathrm dy_i.
\end{aligned}
\]
\end{lemma}
\begin{proof}
\SourceParagraph{ch06-6.3.2-013}{476}
沿这一参数化拉回
\[
\mathbf B+\mathrm i\omega
=\sum_{i,j=1}^n(b_{ij}+\mathrm i\delta_{ij})\,
\mathrm d\check x_j\wedge\mathrm dy_i,
\]
得到
\[
\begin{aligned}
&\sum_{i=1}^p\sum_{j=p+1}^n(b_{ij}+\mathrm i\delta_{ij})\,
\mathrm d\check x_j\wedge\mathrm dy_i
+\sum_{i=1}^p\sum_{j=1}^n(b_{ij}+\mathrm i\delta_{ij})\,
\mathrm dg^j\wedge\mathrm dy_i\\
&\quad=\sum_{i=1}^p\sum_{j=p+1}^nb_{ij}\,\mathrm d\check x_j\wedge\mathrm dy_i
+\sum_{i=1}^p\sum_{j=1}^n(b_{ij}+\mathrm i\delta_{ij})\,
\mathrm dg^j\wedge\mathrm dy_i,
\end{aligned}
\]
正如所要证明的。
\end{proof}

\SourceParagraph{ch06-6.3.2-014}{477}
\BookStatementNumber{68}
\begin{note}[numbered]\label{note:6.68}
以下将讨论 $L(B',g)$ 上的半平坦微分形式。与\BookSectionRef{sec:6.2.2}{6.2.2}中一样，这些形式在 $L(B',g)$ 的纤维上具有常系数。不过，与该小节不同，不能在整体上将它们分裂为各个 $(p,q)$ 型的形式。事实上，采用引理中的记号，对于 $I\subseteq\{1,\ldots,p\}$、$J\subseteq\{p+1,\ldots,n\}$，$L(B',g)$ 上局部形如 $\mathrm dy_I\wedge\mathrm d\check x_J$ 的形式，会随着坐标的改变与 $g$ 的代表的选择而改变写法：可能出现一些额外项，具有更少的 $\mathrm d\check x_i$ 与更多的 $\mathrm dy_j$。因此，可以在模去 $(p+1,q-1)+(p+2,q-2)+\cdots$ 型形式的意义下，谈论 $(p,q)$ 型形式。例如，引理中的计算说明，$\mathbf B|_{L(B',g)}$ 在模去 $(2,0)$ 型形式的意义下，为 $(1,1)$ 型形式，而 $(2,0)$ 部分本身并不良定义。

\SourceParagraph{ch06-6.3.2-015}{477}
还有一点可能使后面的讨论产生混淆：将讨论两种不同意义下的 $\mathbf B$ 的限制。可以把 $\mathbf B$ 看作 $\check X(B)$ 上的2形式，此时，可以将其限制到子流形，得到子流形上的2形式。或者，将 $\mathbf B$ 看作 $H^1(B,\Lambda_{\mathbb R})$ 的元素，此时，可以将其限制到 $B'$，得到 $H^1(B',\Lambda_{\mathbb R}|_{B'})$ 中的类。对于子流形 $M\subseteq\check X(B)$，把前一种限制写为 $\mathbf B|_M$，后一种限制则写为 $\mathbf B|_{B'}$。
\end{note}

\SourceParagraph{ch06-6.3.2-016}{477}
\BookStatementNumber{69}
\begin{proposition}\label{prop:6.69}
给定有理仿射子流形 $B'\subseteq B$ 与类 $[g]\in H^1(B',\check\Lambda')$，存在代表 $g\in H^0(B',T^*B'/\check\Lambda')$，使得 $L(B',g)$ 为 Lagrange 子流形，当且仅当 $[g]$ 在以下正合列对应的上同调长正合列的连接同态下，于 $H^2(B',\mathbb R)$ 中的像为零：
\[
0\longrightarrow\mathbb R\longrightarrow\operatorname{Aff}_{\mathbb Z}(B',\mathbb R)
\xrightarrow{\mathrm d}\check\Lambda'\longrightarrow0.
\]
这里，$\operatorname{Aff}_{\mathbb Z}(B',\mathbb R)$ 表示 $B'$ 上具有整斜率的仿射线性函数层，$\mathrm d$ 为外微分。此外，形如 $L(B',g)$ 的 Lagrange 子流形的选择，模去上述等价关系后，与 $[g]$ 到 $H^1(B',\operatorname{Aff}_{\mathbb Z}(B',\mathbb R))$ 的全部提升的集合相同。
\end{proposition}
\begin{proof}
\SourceParagraph{ch06-6.3.2-017}{477}
局部地，写成 $g=\sum_{j=1}^ng^j\,\mathrm dy_j$。由引理看到，$L(B',g)$ 为 Lagrange 子流形，当且仅当 $\sum_{j=1}^p\mathrm dg^j\wedge\mathrm dy_j=0$；等价地，将 $g$ 看作 $T^*B'$ 的截面时，有 $\mathrm dg=0$。现在，用 $(T^*B')_{\mathrm{closed}}$ 表示 $T^*B'$ 中由闭1形式构成的子层。那么，$L(B',g)$ 为 Lagrange 子流形，当且仅当 $g\in H^0(B',T^*B'/\check\Lambda')$ 属于 $H^0(B',(T^*B')_{\mathrm{closed}}/\check\Lambda')$。有正合列
\[
0\longrightarrow\check\Lambda'\longrightarrow(T^*B')_{\mathrm{closed}}
\longrightarrow(T^*B')_{\mathrm{closed}}/\check\Lambda'\longrightarrow0,
\]
给出长正合列
\[
H^0(B',(T^*B')_{\mathrm{closed}}/\check\Lambda')
\longrightarrow H^1(B',\check\Lambda')
\longrightarrow H^1(B',(T^*B')_{\mathrm{closed}}).
\]
所以，$[g]\in H^1(B',\check\Lambda')$ 可以提升到 $g\in H^0(B',(T^*B')_{\mathrm{closed}}/\check\Lambda')$，当且仅当 $[g]$ 在 $H^1(B',(T^*B')_{\mathrm{closed}})$ 中的像为零。不过，还有正合列
\[
0\longrightarrow\mathbb R\longrightarrow C^\infty(B')
\xrightarrow{\mathrm d}(T^*B')_{\mathrm{closed}}\longrightarrow0,
\]
其中，$C^\infty(B')$ 为 $B'$ 上的 $C^\infty$ 函数层，它说明
\[
H^1(B',(T^*B')_{\mathrm{closed}})\cong H^2(B',\mathbb R).
\]

\SourceParagraph{ch06-6.3.2-018}{478}
注意，有如下正合列的交换图：
\[
\BookImage{02_Unnumbered_Diagrams/Chapter_06/Ch06_Diagram_PDFp478_01_Bookp468.pdf}
\]
因此，$[g]$ 提升为给出 Lagrange 子流形 $L(B',g)$ 的截面 $g$，当且仅当 $[g]$ 在第一行的正合列的连接同态下，于 $H^2(B',\mathbb R)$ 中的像为零。此外，由于 $H^1(B',C^\infty(B'))=0$，若 $[g]$ 提升到 $\overline g\in H^1(B',\operatorname{Aff}_{\mathbb Z}(B',\mathbb R))$，则它进而提升到元素
\[
\begin{aligned}
g&\in H^0(B',C^\infty(B')/\operatorname{Aff}_{\mathbb Z}(B',\mathbb R))\\
&\cong H^0(B',(T^*B')_{\mathrm{closed}}/\check\Lambda').
\end{aligned}
\]
任意这样的提升，都使 $L(B',g)$ 为 Lagrange 子流形。两个这样的提升 $g,g'$ 等价，当且仅当 $g-g'$ 来自 $H^0(B',C^\infty(B'))$。所以，看到 $[g]$ 的提升 $\overline g$ 指定 Lagrange 子流形 $L(B',g)$ 的一个等价类。
\end{proof}

\SourceParagraph{ch06-6.3.2-019}{478}
\BookStatementNumber{70}
\begin{corollary}\label{cor:6.70}
给定有理仿射子流形 $B'\subseteq B$ 与类 $[g]\in H^1(B',\check\Lambda')$，存在 $[g]$ 的代表 $g$，使得 $L(B',g)$ 为 Lagrange 子流形，当且仅当相对于自然杯积
\[
H^1(B',\check\Lambda')\times H^1(B',\Lambda_{\mathbb R}')
\longrightarrow H^2(B',\mathbb R),
\]
有 $[g]\cdot[\delta_{B'}]=0$，其中，$[\delta_{B'}]\in H^1(B',\Lambda_{\mathbb R}')$ 为 $B'$ 的辐射障碍。
\end{corollary}
\begin{proof}
\SourceParagraph{ch06-6.3.2-020}{479}
如\BookRef{prop:6.59}{命题6.59}的证明中一样，可以说明，正合列
\[
0\longrightarrow\mathbb R\longrightarrow\operatorname{Aff}_{\mathbb Z}(B',\mathbb R)
\longrightarrow\check\Lambda'\longrightarrow0
\]
的扩张类，就是 $B'$ 的辐射障碍 $[\delta_{B'}]$。由此看到，$[g]\in H^1(B',\check\Lambda')$ 在 $H^2(B',\mathbb R)$ 中的像，正是 $[g]$ 与辐射障碍 $[\delta_{B'}]$ 的杯积，与\BookRef{prop:6.60}{命题6.60}证明中的情形完全相同。参见例如文献\BookCite{ref206}的命题1.12。
\end{proof}

\SourceParagraph{ch06-6.3.2-021}{479}
要理解具有所需性质的 $\operatorname{U}(1)$ 丛何时存在，注意，需要2形式 $\mathbf B$ 在 $L=L(B',g)$ 上的限制，代表一个来自 $H^2(L,\mathbb Z)$ 的整类。一旦满足这一条件，$L$ 上就存在合适的 $\operatorname{U}(1)$ 丛与联络。给定一个这样的选择，所有满足曲率条件的丛及联络，模去规范等价后，由 $H^1(L,\mathbb R/\mathbb Z)$ 参数化；固定底层丛的拓扑类型时，应取其中相应的分量。

\SourceParagraph{ch06-6.3.2-022}{479}
利用目前的工具，分析这一整性条件有些困难。因此，先理解一个较简单的条件：$\mathbf B|_{L(B',g)}$ 代表 $H^2(L(B',g),\mathbb R)$ 中的零元。也就是说，$L(B',g)$ 上存在整体定义的1形式 $A$，满足 $\mathrm dA=\mathbf B|_{L(B',g)}$。这样的 $A$，当然可以看作归一化联络1形式；取 $\nabla=\mathrm d-2\pi\mathrm i A$，便得到曲率为 $-2\pi\mathrm i\,\mathbf B|_{L(B',g)}$ 的联络。希望在相差恰当形式的意义下分类这些 $A$；这些恰当形式给出与恒等规范变换相连的规范改变。

\SourceParagraph{ch06-6.3.2-023}{479}
\BookStatementNumber{71}
\begin{theorem}\label{thm:6.71}
$\mathbf B|_{L(B',g)}$ 在 $L(B',g)$ 上为恰当形式，当且仅当在正合列
\[
H^1(B',\Lambda_{\mathbb R}')\longrightarrow H^1(B',\Lambda_{\mathbb R}|_{B'})
\longrightarrow H^1(B',\Lambda_{\mathbb R}|_{B'}/\Lambda_{\mathbb R}')
\]
中，存在 $\mathbf B|_{B'}\in H^1(B',\Lambda_{\mathbb R}|_{B'})$ 到 $\mathbf B'\in H^1(B',\Lambda_{\mathbb R}')$ 的提升，满足相对于自然杯积
\[
H^1(B',\check\Lambda')\times H^1(B',\Lambda_{\mathbb R}')
\longrightarrow H^2(B',\mathbb R),
\]
有 $[g]\cdot[\mathbf B']=0$。

\SourceParagraph{ch06-6.3.2-024}{479}
此外，若存在这样的 $\mathbf B'$，则 $L(B',g)$ 上满足 $\mathrm dA=\mathbf B|_{L(B',g)}$ 的1形式 $A$，模去恰当形式后，与下列数据一一对应：
\begin{enumerate}[label=(\arabic*)]
\item 选择 $\mathbf B|_{B'}$ 的提升 $\mathbf B'$，以及
\[
\alpha\in\Gamma\bigl(B',C^\infty(B')\otimes_{\mathbb R}
(\Lambda_{\mathbb R}|_{B'}/\Lambda_{\mathbb R}')\bigr),
\]
使得对于 $\alpha$ 到 $\widetilde\alpha\in\Gamma(B',C^\infty(B')\otimes_{\mathbb R}\Lambda_{\mathbb R}|_{B'})$ 的任意提升，$\mathbf B|_{B'}-\nabla\widetilde\alpha$ 代表 $\mathbf B'$。这里，$\nabla$ 是 $C^\infty(B')\otimes_{\mathbb R}\Lambda_{\mathbb R}|_{B'}$ 上的自然联络，其平坦截面就是 $\Lambda_{\mathbb R}|_{B'}$ 的截面。对于给定的 $\mathbf B'$，这样的 $\alpha$ 的选择集合，是群
\[
\operatorname{im}\!\left(H^0(B',\Lambda_{\mathbb R}|_{B'})
\longrightarrow H^0(B',\Lambda_{\mathbb R}|_{B'}/\Lambda_{\mathbb R}')\right)
\]
的主齐性空间。
\item 选择 $[g]\in H^1(B',\check\Lambda')$ 到 $H^1(B',\mathcal P^{\mathbb R}_{B',\mathbf B'})$ 的提升，连同它与固定的 $g$ 所给出的第二分量代表的识别。在这一相对提升中，只模去第一分量的恰当1形式。其中，$\mathcal P^{\mathbb R}_{B',\mathbf B'}$ 是由扩张
\[
0\longrightarrow\mathbb R\longrightarrow\mathcal P^{\mathbb R}_{B',\mathbf B'}
\longrightarrow\check\Lambda'\longrightarrow0
\]
定义的层，扩张类为 $\mathbf B'\in H^1(B',\Lambda_{\mathbb R}')$。
\end{enumerate}
\end{theorem}

\begin{proof}
\SourceParagraph{ch06-6.3.2-025}{480}
由\BookRef{lem:6.67}{引理6.67}看到，$\mathbf B|_{L(B',g)}$ 局部地可以写为 $(1,1)$ 型与 $(2,0)$ 型半平坦微分形式之和。$(2,0)$ 部分并不良定义，因为改变 $g$ 的代表，即把 $g$ 换成另一个与之模 $\check\Lambda|_{B'}+(TB')^\perp$ 同余的 $T^*B|_{B'}$ 截面，会改变 $(2,0)$ 部分。不过，$(1,1)$ 部分不依赖于选择，并且，像\BookSectionRef{sec:6.2.2}{6.2.2}中关联于局部系统的形式一样，在局部系统 $\Lambda_{\mathbb R}|_{B'}/\Lambda_{\mathbb R}'$ 的上同调中，给出良定义的元素 $\beta\in H^1(B',\Lambda_{\mathbb R}|_{B'}/\Lambda_{\mathbb R}')$。注意，$\beta$ 是 $\mathbf B|_{B'}\in H^1(B',\Lambda_{\mathbb R}|_{B'})$ 在自然映射
\[
H^1(B',\Lambda_{\mathbb R}|_{B'})
\longrightarrow H^1(B',\Lambda_{\mathbb R}|_{B'}/\Lambda_{\mathbb R}')
\]
下的像。以下也用 $\beta$ 表示由给定形式得到的 de Rham 代表。

\SourceParagraph{ch06-6.3.2-026}{480}
由正合列
\[
0\longrightarrow\Lambda_{\mathbb R}'\longrightarrow\Lambda_{\mathbb R}|_{B'}
\longrightarrow\Lambda_{\mathbb R}|_{B'}/\Lambda_{\mathbb R}'\longrightarrow0,
\]
看到，$\mathbf B|_{B'}$ 提升到 $\mathbf B'\in H^1(B',\Lambda_{\mathbb R}')$，当且仅当 $\beta$ 的上同调类为零；这又等价于存在截面
\[
\alpha\in\Gamma\bigl(B',C^\infty(B')\otimes_{\mathbb R}
(\Lambda_{\mathbb R}|_{B'}/\Lambda_{\mathbb R}')\bigr),
\]
满足 $\nabla\alpha=\beta$，其中，
\[
\nabla\colon C^\infty(B')\otimes_{\mathbb R}
(\Lambda_{\mathbb R}|_{B'}/\Lambda_{\mathbb R}')
\longrightarrow T^*B'\otimes
(\Lambda_{\mathbb R}|_{B'}/\Lambda_{\mathbb R}')
\]
是\BookSectionRef{sec:6.2.2}{6.2.2}中定义的联络。注意，若 $\mathbf B|_{L(B',g)}=\mathrm dA$，则在相差恰当形式的意义下，$A$ 可以取为 $(0,1)$ 型与 $(1,0)$ 型半平坦形式之和。按照所采用的楔积次序，$(0,1)$ 部分的负系数，确定上述满足 $\nabla\alpha=\beta$ 的 $\alpha$。因此，若 $\mathbf B|_{L(B',g)}$ 恰当，则 $\mathbf B|_{B'}$ 可以提升到 $H^1(B',\Lambda_{\mathbb R}')$。

\SourceParagraph{ch06-6.3.2-027}{480}
无论如何，提升 $\mathbf B'$ 由 $\alpha$ 确定：将
\[
\alpha\in\Gamma\bigl(B',C^\infty(B')\otimes_{\mathbb R}
(\Lambda_{\mathbb R}|_{B'}/\Lambda_{\mathbb R}')\bigr)
\]
任意提升为 $\widetilde\alpha\in\Gamma(B',C^\infty(B')\otimes_{\mathbb R}\Lambda_{\mathbb R}|_{B'})$，然后，$\mathbf B'=\mathbf B|_{B'}-\nabla\widetilde\alpha$ 定义 $\Lambda_{\mathbb R}'\otimes T^*B'$ 的截面，确定元素 $\mathbf B'\in H^1(B',\Lambda_{\mathbb R}')$。注意，$\alpha$ 的任意两个提升 $\widetilde\alpha_1,\widetilde\alpha_2$，都满足 $\widetilde\alpha_1-\widetilde\alpha_2\in\Gamma(B',C^\infty(B')\otimes_{\mathbb R}\Lambda_{\mathbb R}')$，所以，$\nabla(\widetilde\alpha_1-\widetilde\alpha_2)$ 代表 $H^1(B',\Lambda_{\mathbb R}')$ 中的零上同调类。因此，$\widetilde\alpha_1,\widetilde\alpha_2$ 确定的提升 $\mathbf B_1',\mathbf B_2'$ 同调，$\mathbf B'$ 的上同调类完全由 $\alpha$ 确定。另一方面，给 $\alpha$ 加上
\[
\alpha'\in\Gamma\bigl(B',C^\infty(B')\otimes_{\mathbb R}
(\Lambda_{\mathbb R}|_{B'}/\Lambda_{\mathbb R}')\bigr),
\qquad\nabla\alpha'=0,
\]
会使 $\mathbf B'$ 的类，改变一个等于 $\alpha'\in H^0(B',\Lambda_{\mathbb R}|_{B'}/\Lambda_{\mathbb R}')$ 在 $H^1(B',\Lambda_{\mathbb R}')$ 中的像的负值的元素。

\SourceParagraph{ch06-6.3.2-028}{481}
具体地，在局部坐标中，若 $\widetilde\alpha=\sum_{j=1}^n\widetilde\alpha_j\partial/\partial y_j$，则相应的半平坦 $(0,1)$ 形式为 $a=-\sum_{j=1}^n\widetilde\alpha_j\,\mathrm d\check x_j$，并且
\begin{equation}\label{eq:6.4}
\begin{aligned}
\mathbf B'
&=\sum_{i=1}^p\sum_{j=1}^n
\left(b_{ij}-\frac{\partial\widetilde\alpha_j}{\partial y_i}\right)
\mathrm d\check x_j\wedge\mathrm dy_i\\
&=\sum_{i,j=1}^p
\left(b_{ij}-\frac{\partial\widetilde\alpha_j}{\partial y_i}\right)
\mathrm d\check x_j\wedge\mathrm dy_i.
\end{aligned}
\end{equation}
后一等式由 $\mathbf B|_{B'}-\nabla\widetilde\alpha$ 为 $\Lambda_{\mathbb R}'\otimes T^*B'$ 的截面这一假设得到。现在，$\widetilde\alpha$ 给出 $\check f^{-1}(B')$ 上的半平坦 $(0,1)$ 形式 $a$；于是，可以把 $\mathbf B'$ 看作 $\check f^{-1}(B')$ 上的半平坦形式，即 $\mathbf B|_{\check f^{-1}(B')}-\mathrm da$，它当然与 $\mathbf B|_{\check f^{-1}(B')}$ 同调。进一步将其限制到 $L(B',g)$，由引理，并用 $g$ 的局部代表 $\sum_{j=1}^pg^j\,\mathrm dy_j$，得到 $(2,0)$ 型半平坦2形式
\[
\mathbf B|_{L(B',g)}-\mathrm da|_{L(B',g)}
=\sum_{i,j=1}^p
\left(b_{ij}-\frac{\partial\widetilde\alpha_j}{\partial y_i}\right)
\mathrm dg^j\wedge\mathrm dy_i.
\]
它现在在 $B'$ 上整体良定义。若这个形式恰当，即 $\mathbf B|_{L(B',g)}-\mathrm da|_{L(B',g)}=\mathrm d\widetilde\alpha'$，则 $A=a|_{L(B',g)}+\widetilde\alpha'$ 满足 $\mathbf B|_{L(B',g)}=\mathrm dA$。反过来，若存在这样的 $A$，则 $\mathbf B|_{L(B',g)}-\mathrm da|_{L(B',g)}$ 显然恰当。

\SourceParagraph{ch06-6.3.2-029}{481}
现在，理解 $\mathbf B|_{L(B',g)}-\mathrm da|_{L(B',g)}$ 的上同调类。类 $[g]\in H^1(B',\check\Lambda')$ 有如下 de Rham 代表。局部地，$g$ 由 $T^*B'=C^\infty(B')\otimes_{\mathbb R}\check\Lambda_{\mathbb R}'$ 的截面代表，
\[
g=\sum_{i=1}^pg^i\,\mathrm dy_i
=\sum_{i=1}^pg^i\frac{\partial}{\partial\check x_i},
\]
这里，把 $\check\Lambda_{\mathbb R,b}'$ 的元素，看作 $\check X(B')\to B'$ 的纤维的切向量。两个代表 $g,g'$ 相差 $(g')^i=g^i+n_i$，其中，$n_i\in\mathbb Z$。局部地施加自然联络，见\BookSectionRef{sec:6.2.2}{6.2.2}，
\[
\nabla\colon C^\infty(B')\otimes_{\mathbb R}\check\Lambda_{\mathbb R}'
\longrightarrow T^*B'\otimes\check\Lambda_{\mathbb R}',
\]
得到 $\nabla g=\sum_{i=1}^p\mathrm dg^i\otimes\partial/\partial\check x_i$。它不再依赖于代表的选择，给出 $B'$ 上的 $T^*B'\otimes\check\Lambda_{\mathbb R}'$ 截面。容易看到，这个截面代表 $[g]$ 的实像。杯积
\[
H^1(B',\check\Lambda')\times H^1(B',\Lambda_{\mathbb R}')
\longrightarrow H^2(B',\mathbb R)
\]
局部地由
\[
\begin{aligned}
\left(\sum_{i,j}\alpha_{ij}\,\mathrm dy_i\otimes
\frac{\partial}{\partial\check x_j}\right)
\cdot\left(\sum_{k,l}\beta_{kl}\,\mathrm dy_k\otimes
\frac{\partial}{\partial y_l}\right)
\longmapsto\sum_{i,j,k,l}\alpha_{ij}\beta_{kl}\delta_{jl}\,
\mathrm dy_i\wedge\mathrm dy_k
\end{aligned}
\]
代表。因此，$[g]$ 与 $\mathbf B'$ 的杯积，局部地由
\[
\sum_{i,j=1}^p
\left(b_{ij}-\frac{\partial\widetilde\alpha_j}{\partial y_i}\right)
\mathrm dg^j\wedge\mathrm dy_i
\]
代表，正是形式 $\mathbf B|_{L(U,g)}-\mathrm da|_{L(U,g)}$。这说明，$\mathbf B|_{L(B',g)}$ 恰当，当且仅当 $\mathbf B|_{B'}$ 在 $H^1(B',\Lambda_{\mathbb R}|_{B'}/\Lambda_{\mathbb R}')$ 中的像为零，并且对于 $\mathbf B|_{B'}$ 的某个提升 $\mathbf B'$，在 $H^2(B',\mathbb R)$ 中有杯积 $[g]\cdot\mathbf B'=0$。

\SourceParagraph{ch06-6.3.2-030}{481}
只剩下描述 $L(B',g)$ 上满足 $\mathrm dA=\mathbf B|_{L(B',g)}$ 的1形式 $A$，模去恰当形式后的选择。若 $\mathrm dA=\mathbf B|_{L(B',g)}$，则由于 $\mathbf B$ 半平坦，不难说明，$A$ 在相差恰当形式的意义下，可以取为 $L(B',g)$ 上的半平坦1形式。因此，可以假定 $A$ 半平坦。

\SourceParagraph{ch06-6.3.2-031}{482}
首先，如前所述，选择 $A$ 的 $(0,1)$ 部分，等价于选择满足 $\nabla\alpha=\beta$ 的
\[
\alpha\in\Gamma\bigl(B',C^\infty(B')\otimes_{\mathbb R}
(\Lambda_{\mathbb R}|_{B'}/\Lambda_{\mathbb R}')\bigr).
\]
截面 $\alpha$ 在 $L(B',g)$ 上，指定模去 $(1,0)$ 型形式后的半平坦 $(0,1)$ 型形式；用上述负系数的识别，这个形式的微分，在模去 $(2,0)$ 型半平坦形式的意义下，为 $\mathbf B|_{L(B',g)}$。选择 $\alpha$ 的提升 $\widetilde\alpha$，就是为这一形式选择一个具体代表 $a$，并且，特别地，给出 $\mathbf B'$ 的选择。

\SourceParagraph{ch06-6.3.2-032}{482}
给定 $\widetilde\alpha$，现在希望说明，满足
\begin{equation}\label{eq:6.5}
\mathrm d\widetilde\alpha'
=\mathbf B|_{L(B',g)}-\mathrm da|_{L(B',g)}
\end{equation}
的半平坦 $\widetilde\alpha'$，模去恰当形式后的选择，与 $[g]$ 到 $H^1(B',\mathcal P^{\mathbb R}_{B',\mathbf B'})$ 的、固定第二分量代表的相对提升一一对应。事实上，扩张 $\mathcal P^{\mathbb R}_{B',\mathbf B'}$ 可以如下表达。取这一扩张的实线性化，所得光滑向量丛分裂为
\[
C^\infty(B')\oplus
\bigl(C^\infty(B')\otimes_{\mathbb R}\check\Lambda_{\mathbb R}'\bigr),
\]
带有平坦联络 $\nabla_{\mathcal P}$：
\[
\nabla_{\mathcal P}(f,\zeta)
=\bigl(\mathrm df-\mathbf B'(\zeta),\nabla_{\check\Lambda'}\zeta\bigr).
\]
这里，$\nabla_{\check\Lambda'}$ 是 $C^\infty(B')\otimes_{\mathbb R}\check\Lambda_{\mathbb R}'$ 上的平坦联络，平坦截面就是 $\check\Lambda_{\mathbb R}'$ 的截面；而
\[
\mathbf B'\in T^*B'\otimes\Lambda_{\mathbb R}'
=\operatorname{Hom}_{C^\infty(B')}
\bigl(\check\Lambda_{\mathbb R}'\otimes_{\mathbb R}C^\infty(B'),T^*B'\bigr)
\]
在 $C^\infty(B')\otimes_{\mathbb R}\check\Lambda_{\mathbb R}'$ 的截面上取值，给出 $B'$ 上的1形式。一个初等练习是：若定义 $\mathcal P^{\mathbb R}_{B',\mathbf B'}$ 为上述直和中相对于这一联络平坦、并且投影到第二分量后属于整平坦截面 $\check\Lambda'$ 的截面层，则其扩张类确实为 $\mathbf B'$。可以利用\BookRef{prop:6.59}{命题6.59}中使用的扩张类的 Čech 代表，完成这一证明。

\SourceParagraph{ch06-6.3.2-033}{482}
现在，给定类 $[g]\in H^1(B',\check\Lambda')$，由 $\widetilde g\in T^*B'\otimes\check\Lambda_{\mathbb R}'$ 代表，局部地写为 $\widetilde g=\sum_{i=1}^p\mathrm dg^i\otimes\partial/\partial\check x_i$，满足 $\nabla_{\check\Lambda'}\widetilde g=0$。选择它的上述相对提升，就是给出第二分量固定的配对
\[
(-\widetilde\alpha',\widetilde g)
\in T^*B'\oplus(T^*B'\otimes\check\Lambda_{\mathbb R}'),
\]
满足
\[
\begin{aligned}
0=\nabla_{\mathcal P}(-\widetilde\alpha',\widetilde g)
&=\bigl(-\mathrm d\widetilde\alpha'-\mathbf B'(\widetilde g),
\nabla_{\check\Lambda'}\widetilde g\bigr)\\
&=\bigl(-\mathrm d\widetilde\alpha'-\mathbf B'(\widetilde g),0\bigr).
\end{aligned}
\]
不过，局部地有
\[
\mathbf B'(\widetilde g)
=\sum_{i,j=1}^p
\left(b_{ij}-\frac{\partial\widetilde\alpha_j}{\partial y_i}\right)
\mathrm dy_i\wedge\mathrm dg^j,
\]
所以，$\nabla_{\mathcal P}(-\widetilde\alpha',\widetilde g)=0$，正是\BookEqRef{eq:6.5}{6.5}的陈述。因此，固定投影代表、只模去第一分量中恰当形式的提升选择，等价于 $\widetilde\alpha'$ 模去恰当1形式后的选择。结论得证。
\end{proof}

\SourceParagraph{ch06-6.3.2-034}{483}
现在，不给出证明地陈述所需的完整结论，用来把半平坦 A 膜，与前一小节的半平坦 B 膜比较。需要理解，$\mathbf B|_{L(B',g)}$ 何时为整类，而不只是恰当形式。这里省略证明，因为论证需要更深入地理解上同调类的 Čech 代表，超出了这里希望展开的范围。

\SourceParagraph{ch06-6.3.2-035}{483}
\BookStatementNumber{72}
\begin{theorem}\label{thm:6.72}
$\mathbf B|_{L(B',g)}$ 在 $L(B',g)$ 上为整类，当且仅当在正合列
\[
\begin{aligned}
H^1(B',\Lambda_{\mathbb R}'/\Lambda')
&\longrightarrow H^1(B',(\Lambda_{\mathbb R}/\Lambda)|_{B'})\\
&\longrightarrow H^1(B',\Lambda_{\mathbb R}|_{B'}/(\Lambda|_{B'}+\Lambda_{\mathbb R}'))
\end{aligned}
\]
中，存在 $\mathbf B|_{B'}\in H^1(B',(\Lambda_{\mathbb R}/\Lambda)|_{B'})$ 到 $\mathbf B'\in H^1(B',\Lambda_{\mathbb R}'/\Lambda')$ 的提升，满足相对于自然杯积
\[
H^1(B',\check\Lambda')\times H^1(B',\Lambda_{\mathbb R}'/\Lambda')
\longrightarrow H^2(B',\mathbb R/\mathbb Z),
\]
有 $[g]\cdot[\mathbf B']=0$。

\SourceParagraph{ch06-6.3.2-036}{483}
此外，若 $\mathbf B|_{L(B',g)}$ 为整类，且 $\mathcal L$ 为 $L(B',g)$ 上的 $\operatorname{U}(1)$ 丛，第一 Chern 类取为这一实类的指定整提升，则 $\mathcal L$ 上曲率为 $-2\pi\mathrm i\,\mathbf B|_{L(B',g)}$ 的半平坦联络，模去规范等价后，与下列数据中对应于这一拓扑类型的部分一一对应。若允许具有这一曲率的全部底层丛，则取全部相应数据。
\begin{enumerate}[label=(\arabic*)]
\item 选择 $\mathbf B|_{B'}$ 的提升 $\mathbf B'$，以及取值于
\[
\Lambda_{\mathbb R}|_{B'}/(\Lambda|_{B'}+\Lambda_{\mathbb R}')
\]
的 Čech 0上链 $(h_i)$，使得对于 $h_i$ 到 $(\Lambda_{\mathbb R}/\Lambda)|_{B'}$ 的任意提升 $\widetilde h_i$，$(b_{ij}+\widetilde h_i-\widetilde h_j)_{ij}$ 代表 $\mathbf B'$。对于给定的 $\mathbf B'$，$(h_i)$ 的选择集合，是群
\[
\operatorname{im}\!\left(H^0(B',(\Lambda_{\mathbb R}/\Lambda)|_{B'})
\longrightarrow H^0(B',\Lambda_{\mathbb R}|_{B'}/(\Lambda|_{B'}+\Lambda_{\mathbb R}'))\right)
\]
的主齐性空间。
\item 选择 $[g]\in H^1(B',\check\Lambda')$ 到 $H^1(B',\mathcal P^{\mathbb R/\mathbb Z}_{B',\mathbf B'})$ 的提升，连同其投影到固定 $g$ 的 Čech 代表的识别；等价通过 $\mathbb R/\mathbb Z$ 值的0上链给出的规范改变取商。其中，$\mathcal P^{\mathbb R/\mathbb Z}_{B',\mathbf B'}$ 为由扩张
\[
0\longrightarrow\mathbb R/\mathbb Z
\longrightarrow\mathcal P^{\mathbb R/\mathbb Z}_{B',\mathbf B'}
\longrightarrow\check\Lambda'\longrightarrow0
\]
定义的层，扩张类为 $\mathbf B'\in H^1(B',\Lambda_{\mathbb R}'/\Lambda')$。
\end{enumerate}
\end{theorem}

\SourceParagraph{ch06-6.3.2-037}{483}
把这一结论与\BookRef{thm:6.63}{定理6.63}比较，使几何提升与联络提升的投影识别彼此相容，并同时取商，就看到，半平坦 A 膜与 B 膜的模空间对应。这里有一点细微之处需要指出：这种同构并非典范。事实上，为了得到同构，正在对主齐性空间作非典范的识别。具有给定 Chern 类的线丛集合，是零 Chern 类线丛群的主齐性空间；也就是说，零 Chern 类线丛群，通过张量积自由且传递地作用于给定 Chern 类的线丛集合。不过，零 Chern 类的线丛，与给定 Chern 类的线丛之间，没有典范双射。选定一个具有该 Chern 类的线丛以后，才能实现这样的识别。

\SourceParagraph{ch06-6.3.2-038}{484}
同样，给定曲率的 $\operatorname{U}(1)$ 丛及联络的规范等价类空间上，有群 $H^1(X,\mathbb R/\mathbb Z)$ 的作用；不过，这个空间并不典范地同构于 $H^1(X,\mathbb R/\mathbb Z)$。只有选定其中一个丛及联络以后，才能得到同构。因此，这个空间也是 $H^1(X,\mathbb R/\mathbb Z)$ 的主齐性空间。若固定线丛 $\mathcal L$ 的拓扑类型，应限制到相应分量，此时平移群为 $H^1(X,\mathbb R)/\operatorname{im}H^1(X,\mathbb Z)$。两侧都存在这种非典范性，意味着按照这里的表述，并未展示半平坦 A 膜与 B 膜模空间之间的典范同构。

\SourceParagraph{ch06-6.3.2-039}{484}
\BookStatementNumber{73}
\begin{example}\label{ex:6.73}
考虑 $B=\mathbb R/\tau_2\mathbb Z$ 的情形。采用与\BookRef{ex:6.27}{例题6.27}相同的归一化，令
\[
\mathbf B=\frac{\tau_1}{\tau_2}\,\mathrm d\check x\wedge\mathrm dy,
\]
使得 $X(B,\mathbf B)\cong\mathbb C/\langle1,\tau\rangle$，为周期1与 $\tau=\tau_1+\mathrm i\tau_2$ 的椭圆曲线，并取 $B'=B$。那么，在 $B'$ 上纤维化的半平坦 B 膜集合，就是 $X(B,\mathbf B)$ 的 Picard 群。

\SourceParagraph{ch06-6.3.2-040}{484}
在另一侧，$\check X(B)$ 只是一个辛环面，而 $S^1$ 纤维化 $\check X(B)\to B$ 的任意截面都是 Lagrange 的。可以选择类 $[g]\in H^1(B,\check\Lambda)=\mathbb Z$。将 $\check X(B)$ 看作 $\mathbb R^2/\langle(1,0),(0,\tau_2)\rangle$，任意截面都与斜率为 $[g]/\tau_2$ 的直线 $x=[g]y/\tau_2+b$ 给出的截面 Hamilton 同痕；其中，$[g]$ 看作整数，$b\in\mathbb R$ 在模 $\mathbb Z$ 的意义下定义。注意，由\BookRef{prop:6.69}{命题6.69}的正合列，按周期 $\tau_2$ 将常数部分归一化后，$H^1(B,\operatorname{Aff}_{\mathbb Z}(B,\mathbb R))$ 位于正合列
\[
0\longrightarrow\mathbb R/\mathbb Z
\longrightarrow H^1(B,\operatorname{Aff}_{\mathbb Z}(B,\mathbb R))
\longrightarrow\mathbb Z\longrightarrow0
\]
中。事实上，二元组 $(b,[g])$ 定义 $H^1(B,\operatorname{Aff}_{\mathbb Z}(B,\mathbb R))$ 的一个元素。当然，在这个情形中，$\mathbf B|_{L(B,g)}=0$，$L(B,g)$ 上的平坦 $\operatorname{U}(1)$ 丛及联络，模去规范等价后的空间，就是 $H^1(L(B,g),\mathbb R/\mathbb Z)=\mathbb R/\mathbb Z$。因此，在 A 膜与 B 膜的对应中，$X(B,\mathbf B)$ 上次数为 $d$ 的线丛，对应于 $\check X(B)$ 上斜率为 $d/\tau_2$ 的直线；B 膜的同构类，随后确定 $L(B,g)$ 的平移部分与 $L(B,g)$ 上的联络。
\end{example}

\section{紧化}\label{sec:6.4}

\subsection{带奇点的仿射流形}\label{sec:6.4.1}

\SourceParagraph{ch06-6.4.1-001}{484}
虽然前几节描述的半平坦图景很简单，并且对镜像对称的许多方面给出了优美解释，但它实际上并没有提供特别令人满意的一系列例子，基本上只给出复环面的镜像对称。原因当然是，环面丛的拓扑相当受限；特别地，它们的拓扑 Euler 示性数始终为零。因此，利用半平坦图景，甚至无法解释 K3 曲面的镜像对称。要克服这一点，需要把出现奇异纤维的可能性纳入考虑。

\SourceParagraph{ch06-6.4.1-002}{484}
为考察这一问题，将从一些比前面讨论过的更有意思的仿射流形例子开始。

\SourceParagraph{ch06-6.4.1-003}{484}
与通常一样，令 $M=\mathbb Z^n$、$M_{\mathbb R}=M\otimes_{\mathbb Z}\mathbb R$、$N=\operatorname{Hom}(M,\mathbb Z)$、$N_{\mathbb R}=N\otimes_{\mathbb Z}\mathbb R$。令 $\Xi\subseteq M_{\mathbb R}$ 为任意 $n$ 维凸多胞体，内部包含 $0\in M_{\mathbb R}$。进一步，选择 $\Xi$ 的多面体分解 $\mathcal P$，其意义如下：$\mathcal P$ 是包含在 $\Xi$ 中并覆盖 $\Xi$ 的一族闭多胞体，满足
\begin{itemize}
\item 若 $\sigma\in\mathcal P$，则 $\sigma$ 的每个面都属于 $\mathcal P$；
\item 若 $\sigma_1,\sigma_2\in\mathcal P$，则 $\sigma_1\cap\sigma_2$ 是 $\sigma_1$ 与 $\sigma_2$ 的面。
\end{itemize}

\SourceParagraph{ch06-6.4.1-004}{485}
回顾一下，若 $\sigma\subset M_{\mathbb R}$ 为多胞体，则这里所用的重心 $\operatorname{Bar}(\sigma)$，为 $\sigma$ 的顶点的平均值。$\sigma$ 的第一次重心细分，是由以下全部单形组成的三角剖分：取 $\sigma$ 的面的一条按包含关系递增的链，再取链中各个面的重心所张成的单形。因此，给定多面体分解 $\mathcal P$，可以定义其第一次重心细分 $\operatorname{Bar}(\mathcal P)$，为所有 $\sigma\in\mathcal P$ 的第一次重心细分中的全部单形组成的三角剖分。

\SourceParagraph{ch06-6.4.1-005}{485}
现在，令 $B=\partial\Xi$，因此，$B$ 为 $(n-1)$ 维球面。定义 $\Delta\subseteq B$，为 $\operatorname{Bar}(\mathcal P)$ 中以下单形的并：它们包含于 $\Xi$ 的 $(n-2)$ 维面中，并且不包含 $\mathcal P$ 的任何顶点。

\SourceParagraph{ch06-6.4.1-006}{485}
\BookStatementNumber{74}
\begin{example}\label{ex:6.74}
\begin{enumerate}[label=(\arabic*)]
\item 取 $\Xi\subseteq\mathbb R^3$ 为以下点的凸包：
\[
\begin{aligned}
P_0&=(-1,-1,-1),&P_1&=(3,-1,-1),\\
P_2&=(-1,3,-1),&P_3&=(-1,-1,3).
\end{aligned}
\]
不指定 $\mathcal P$ 本身，只规定它在每个 $(n-2)$ 维面，在此例中即每条边，上的限制。换言之，对于每个 $(n-2)$ 维面 $\sigma$，描述分解
\[
\mathcal P_\sigma
=\{\tau\cap\sigma\mid\tau\in\mathcal P,\ \tau\cap\sigma\ne\varnothing\}.
\]
因此，在此例中，选择沿每条边的分解，使其分为具有整顶点的四个区间，如\BookRef{fig:6.1}{图1}所示。
\begin{figure}[H]
\centering
\BookImage{01_Numbered_Figures/Chapter_06/Ch06_Fig01_PDFp485_Bookp475.pdf}
\BookNumberOnlyCaption\label{fig:6.1}
\end{figure}
集合 $\Delta$，就是这些区间的重心所构成的集合，共给出24个点。
\end{enumerate}

\SourceParagraph{ch06-6.4.1-007}{485}
\begin{enumerate}[label=(\arabic*),start=2]
\item 提高一维，取 $\Xi$ 为以下点的凸包：
\[
\begin{aligned}
P_0&=(-1,-1,-1,-1),\\
P_1&=(4,-1,-1,-1),\\
P_2&=(-1,4,-1,-1),\\
P_3&=(-1,-1,4,-1),\\
P_4&=(-1,-1,-1,4).
\end{aligned}
\]
\begin{figure}[H]
\centering
\BookImage{01_Numbered_Figures/Chapter_06/Ch06_Fig02_PDFp486_Bookp476.pdf}
\BookNumberOnlyCaption\label{fig:6.2}
\end{figure}
取 $\mathcal P$ 在每个二维面上的限制，如\BookRef{fig:6.2}{图2}所示，使得这个面上的 $\Delta$ 如\BookRef{fig:6.3}{图3}所示。
\begin{figure}[H]
\centering
\BookImage{01_Numbered_Figures/Chapter_06/Ch06_Fig03_PDFp486_Bookp476.pdf}
\BookNumberOnlyCaption\label{fig:6.3}
\end{figure}
\end{enumerate}
\end{example}

\SourceParagraph{ch06-6.4.1-008}{486}
继续这一构造。现在，将在 $B_0:=B\setminus\Delta$ 上构造仿射结构。为此，对于 $\Xi$ 的每个 $(n-1)$ 维面 $\sigma$，令 $W_\sigma=\operatorname{Int}(\sigma)$，其中，$\operatorname{Int}$ 表示相对内部。其次，对于包含在 $\Xi$ 的 $(n-2)$ 维面中的每个 $\mathcal P$ 的顶点 $v$，即 $\{v\}\in\mathcal P$，定义开集
\[
W_v=\bigcup_{\substack{\sigma\in\operatorname{Bar}(\mathcal P)\\v\in\sigma\subseteq B}}
\operatorname{Int}(\sigma).
\]

\SourceParagraph{ch06-6.4.1-009}{486}
\BookStatementNumber{75}
\begin{exercise}\label{exer:6.75}
证明
\[
\{W_\sigma\mid\sigma\text{ 为 }\Xi\text{ 的 }(n-1)\text{ 维面}\}
\cup\{W_v\mid v\text{ 为上述 }\mathcal P\text{ 的顶点}\}
\]
构成 $B_0$ 的开覆盖，并且 $W_{\sigma_1}\cap W_{\sigma_2}\ne\varnothing$ 当且仅当 $\sigma_1=\sigma_2$，$W_{v_1}\cap W_{v_2}\ne\varnothing$ 当且仅当 $v_1=v_2$。
\end{exercise}

\SourceParagraph{ch06-6.4.1-010}{487}
现在，在 $W_\sigma$ 与 $W_v$ 上定义仿射坐标图。当然，$\sigma$ 包含于 $M_{\mathbb R}$ 中的仿射超平面 $A_\sigma$，因此得到典范嵌入
\[
\psi_\sigma\colon\operatorname{Int}(\sigma)\longrightarrow A_\sigma.
\]
其次，考虑投影
\[
\psi_v\colon W_v\subseteq M_{\mathbb R}\longrightarrow M_{\mathbb R}/v\mathbb R.
\]

\SourceParagraph{ch06-6.4.1-011}{487}
\BookStatementNumber{76}
\begin{exercise}\label{exer:6.76}
证明，$\psi_v$ 将 $W_v$ 拓扑嵌入 $M_{\mathbb R}/v\mathbb R$。
\end{exercise}

\SourceParagraph{ch06-6.4.1-012}{487}
显然，转移映射 $\psi_v\circ\psi_\sigma^{-1}$ 为仿射变换。因此，在 $B_0$ 上得到仿射结构。

\SourceParagraph{ch06-6.4.1-013}{487}
\BookStatementNumber{77}
\begin{definition}\label{def:6.77}
带奇点的仿射流形，是拓扑流形 $B$，连同集合 $\Delta\subseteq B$；该集合局部地为有限个局部闭子流形的并，每个子流形的余维至少为2。此外，$B_0=B\setminus\Delta$ 带有仿射流形结构。若 $B_0$ 上的仿射结构为整仿射结构，则称这个带奇点的仿射流形为整的。始终用 $i\colon B_0\to B$ 表示包含映射。
\end{definition}

\SourceParagraph{ch06-6.4.1-014}{487}
在目前的例子中，的确只在 $B_0$ 上有仿射结构，而 $\Delta$ 为奇点集。

\SourceParagraph{ch06-6.4.1-015}{487}
要研究仿射结构的奇点的性质，考察仿射结构的单值变换会很有用。给定以点 $b$ 为基点的环路 $\gamma\in\pi_1(B_0,b)$，可以利用 $TB_0$ 上的平坦联络，沿 $\gamma$ 平行移动 $\Lambda_{\mathbb R,b}=T_bB_0$ 中的向量。这样给出线性映射 $T_\gamma\colon\Lambda_{\mathbb R,b}\to\Lambda_{\mathbb R,b}$，称为绕 $\gamma$ 的单值变换。因此，通过 $\gamma\mapsto T_{\gamma^{-1}}$，得到表示 $\pi_1(B_0,b)\to\operatorname{Aut}(\Lambda_{\mathbb R,b})$，即实切向量局部系统 $\Lambda_{\mathbb R}$ 的单值表示。

\SourceParagraph{ch06-6.4.1-016}{487}
\BookStatementNumber{78}
\begin{exercise}\label{exer:6.78}
证明，$T_\gamma$ 为 $\rho(\gamma^{-1})$ 的线性部分，其中，$\rho\colon\pi_1(B_0,b)\to\operatorname{Aff}(\Lambda_{\mathbb R,b})$ 是\BookRef{def:6.6}{定义6.6}中的完整表示。
\end{exercise}

\SourceParagraph{ch06-6.4.1-017}{487}
当然，若 $\gamma$ 为绕判别集 $\Delta$ 的一条小环路，而 $T_\gamma$ 非平凡，则不可能在 $\gamma$ 附近将仿射结构延拓过 $\Delta$。因此，单值变换可以探测奇点。

\SourceParagraph{ch06-6.4.1-018}{487}
\BookStatementNumber{79}
\begin{proposition}\label{prop:6.79}
令 $\gamma$ 为 $B_0$ 中的环路，基点为 $\mathcal P$ 的一个顶点 $v_1$，该顶点包含于 $\Xi$ 的面 $\tau$ 中。环路先进入包含 $\tau$ 的 $\Xi$ 的 $(n-1)$ 维面 $\sigma_1$，经过包含于 $\tau$ 的 $\mathcal P$ 的顶点 $v_2$，再进入包含 $\tau$ 的 $\Xi$ 的 $(n-1)$ 维面 $\sigma_2$，然后返回 $v_1$，见\BookRef{fig:6.4}{图4}。令 $w_1,w_2\in N_{\mathbb R}$ 为唯一满足 $w_i$ 在 $\sigma_i$ 上取值为1的元素。令 $T_\gamma\colon\Lambda_{\mathbb R,v_1}\to\Lambda_{\mathbb R,v_1}$ 表示绕 $\gamma$ 的平行移动。注意，通过坐标图 $\psi_{v_1}\colon W_{v_1}\to M_{\mathbb R}/v_1\mathbb R$ 的微分，可以把 $\Lambda_{\mathbb R,v_1}$ 作为向量空间与 $M_{\mathbb R}/v_1\mathbb R$ 识别。在这一识别下，$T_\gamma$ 由下式给出：
\[
T_\gamma(m)=m+\langle w_1-w_2,m\rangle(v_2-v_1)
\quad\bmod v_1\mathbb R.
\]
\end{proposition}
\begin{figure}[H]
\centering
\BookImage{01_Numbered_Figures/Chapter_06/Ch06_Fig04_PDFp488_Bookp478.pdf}
\BookNumberOnlyCaption\label{fig:6.4}
\end{figure}
\begin{proof}
\SourceParagraph{ch06-6.4.1-019}{488}
$\sigma_i$ 的切空间，典范地与空间 $w_i^\perp\subseteq M_{\mathbb R}$ 识别。将切向量从 $W_{\sigma_i}$ 平行移动到 $W_{v_j}$，由投影诱导的同构 $w_i^\perp\to M_{\mathbb R}/v_j\mathbb R$ 给出。于是，$T_\gamma$ 通过按顺时针方向绕下面的图一周得到：
\[
\BookImage{02_Unnumbered_Diagrams/Chapter_06/Ch06_Diagram_PDFp488_01_Bookp478.pdf}
\]
因此，若 $m\in M_{\mathbb R}$ 代表 $M_{\mathbb R}/v_1\mathbb R$ 中的一个元素，则
\[
\begin{aligned}
T_\gamma(m)
&=\psi_4\bigl(\psi_3^{-1}(\psi_2(\psi_1^{-1}(m)))\bigr)\\
&=\psi_4\bigl(\psi_3^{-1}(\psi_2(m-\langle w_1,m\rangle v_1))\bigr)\\
&=\psi_4\bigl(\psi_3^{-1}(m-\langle w_1,m\rangle v_1
\bmod v_2\mathbb R)\bigr)\\
&=\psi_4\bigl(m-\langle w_1,m\rangle v_1
-\langle w_2,m-\langle w_1,m\rangle v_1\rangle v_2\bigr)\\
&=m-\langle w_2,m\rangle v_2+\langle w_1,m\rangle v_2
\quad\bmod v_1\mathbb R\\
&=m+\langle w_1-w_2,m\rangle(v_2-v_1)
\quad\bmod v_1\mathbb R.
\end{aligned}
\]
\end{proof}

\SourceParagraph{ch06-6.4.1-020}{488}
注意，对于一般的 $\Xi$，这一变换未必属于 $\operatorname{GL}_{n-1}(\mathbb Z)$。不过，有一种重要情形，确实属于这个群。

\SourceParagraph{ch06-6.4.1-021}{489}
\BookStatementNumber{80}
\begin{definition}\label{def:6.80}
若整格多胞体 $\Xi$ 的内部包含0，并且 $\Xi$ 的任意 $(n-1)$ 维面所张成的仿射超平面，由某个 $y\in N$ 给出的仿射线性方程 $\langle y,\cdot\rangle=1$ 定义，则称 $\Xi$ 为反身多胞体。
\end{definition}

\SourceParagraph{ch06-6.4.1-022}{489}
\BookStatementNumber{81}
\begin{proposition}\label{prop:6.81}
若 $\Xi$ 为反身多胞体，则用 $\Xi$ 的一个仅有整顶点的多面体分解 $\mathcal P$，按上述方式构造的 $B_0$，为整仿射流形。因此，在这种情形中，$B$ 是带奇点的整仿射流形。
\end{proposition}
\begin{proof}
\SourceParagraph{ch06-6.4.1-023}{489}
若 $v$ 为顶点，$\sigma$ 为包含 $v$ 的 $\Xi$ 的 $(n-1)$ 维面，则存在 $w\in N$，在 $\sigma$ 上取值为1。于是，转移映射 $\psi_v\circ\psi_\sigma^{-1}$ 的线性部分，是投影 $w^\perp\to M_{\mathbb R}/v\mathbb R$，它显然限制为整同构
\[
w^\perp\cap M\longrightarrow M/v\mathbb Z.
\]
在 $A_\sigma$ 中选一个整顶点为仿射原点后，平移部分也在相应的整格中。因此，转移映射为整仿射映射。
\end{proof}

\SourceParagraph{ch06-6.4.1-024}{489}
反身多胞体的数量非常庞大。在四维中，Kreuzer 与 Skarke 得到了 $473{,}800{,}776$ 个反身多胞体\BookCite{ref315}。因此，这提供了大量非紧整仿射流形的例子。实际上，还有一种更一般的构造，允许从反身多胞体及一些附加数据产生仿射流形，因而可以得到无穷多个例子。见 Haase 与 Zharkov\BookCite{ref216}，以及 Gross\BookCite{ref203}。

\SourceParagraph{ch06-6.4.1-025}{489}
当然，反身条件并非凭空而来！此时，有必要提醒读者，反身多胞体在最基本的镜像对称构造中，具有根本作用；更多细节，见\BookBibText{ref242}{MS1}第7.10节，或 Batyrev\BookCite{ref032}。反身多胞体 $\Xi$，定义 Gorenstein 射影环面簇 $(\mathbb P_\Xi,\mathcal O_{\mathbb P_\Xi}(1))$。$\mathbb P_\Xi$ 的超平面截面，为 Calabi--Yau 簇。它通常有奇点，不过，可以找到极大的部分保典射影消解，消解余维小于4的奇点。于是，可以在 $\mathbb P_\Xi$ 与 $\mathbb P_{\Xi^*}$ 的超平面截面中找到 Calabi--Yau 镜像对，其中，$\Xi^*$ 为 $\Xi$ 的对偶多胞体，或极多胞体，定义为
\[
\Xi^*=\{\nu\in N_{\mathbb R}\mid
\langle\nu,m\rangle\geq-1\text{ 对所有 }m\in\Xi\text{ 成立}\}.
\]
以下讨论中，读者应当记住的基本思想是：由反身多胞体 $\Xi$ 产生的整仿射流形 $B$，应当同时给出 $\mathbb P_\Xi$ 的超平面截面，以及它的镜像。更准确地说，$B$ 带有开子集 $B_0\subseteq B$，其上具有整仿射结构。这产生 $X(B_0)$ 与 $\check X(B_0)$。基本思想是，$X(B_0)$ 应当能够紧化为 $\mathbb P_{\Xi^*}$ 的超平面截面，$\check X(B_0)$ 应当能够紧化为 $\mathbb P_\Xi$ 的超平面截面。不过，这些只是拓扑陈述。细节见 Gross 与 Siebert\BookCite{ref204}，或 Gross\BookCite{ref202}。

\SourceParagraph{ch06-6.4.1-026}{490}
\BookStatementNumber{82}
\begin{example}\label{ex:6.82}
\BookRef{ex:6.74}{例题6.74}中的多胞体，为反身多胞体。来考察这些例子中的单值变换。
\begin{enumerate}[label=(\arabic*)]
\item 选择沿一条边的任意两个相邻顶点。任意选择都会给出相同结果，因此，取 $v_1=(-1,-1,-1)$、$v_2=(0,-1,-1)$。那么，对于同时包含 $v_1,v_2$ 的两个二维面 $\sigma_1,\sigma_2$，可以取 $w_1=(0,0,-1)$、$w_2=(0,-1,0)$。因此，
\[
T_\gamma(m)=m+\langle(0,1,-1),m\rangle(1,0,0)
\quad\bmod v_1\mathbb R.
\]
若取 $(1,0,0)\bmod v_1$ 与 $(0,1,0)\bmod v_1$ 构成 $M/v_1\mathbb Z$ 的一组基，则 $T_\gamma$ 的矩阵为 $\begin{pmatrix}1&1\\0&1\end{pmatrix}$。
\end{enumerate}

\SourceParagraph{ch06-6.4.1-027}{490}
\begin{enumerate}[label=(\arabic*),start=2]
\item 在这一情形中，值得研究判别集 $\Delta$ 的每个顶点的邻域中的单值变换行为。存在两种不同类型的顶点：一种位于二维面的内部，另一种位于一维面的内部。详细考察这两种情形。首先，考虑顶点为
\[
v_1=(-1,-1,-1,-1),\quad
v_2=(0,-1,-1,-1),\quad
v_3=(-1,0,-1,-1)
\]
的三角形。它是 $\mathcal P$ 的一个二维胞腔，包含于 $\Xi$ 的两个三维面 $\sigma_1,\sigma_2$ 中，分别有 $w_1=(0,0,-1,0)$、$w_2=(0,0,0,-1)$。令 $\gamma_i$ 为\BookRef{prop:6.79}{命题6.79}中涉及 $v_i,v_{i+1},\sigma_1,\sigma_2$ 的环路，指标按模3计算。那么，由\BookRef{prop:6.79}{命题6.79}，得到
\[
\begin{aligned}
T_{\gamma_1}(m)&=m+\langle(0,0,-1,1),m\rangle(1,0,0,0)
\quad\bmod v_1\mathbb R,\\
T_{\gamma_2}(m)&=m+\langle(0,0,-1,1),m\rangle(-1,1,0,0)
\quad\bmod v_2\mathbb R,\\
T_{\gamma_3}(m)&=m+\langle(0,0,-1,1),m\rangle(0,-1,0,0)
\quad\bmod v_3\mathbb R.
\end{aligned}
\]
\end{enumerate}

\SourceParagraph{ch06-6.4.1-028}{490}
现在，$T_{\gamma_i}\colon M_{\mathbb R}/v_i\mathbb R\to M_{\mathbb R}/v_i\mathbb R$。因此，要比较这三个单值变换，将切空间 $\Lambda_{\mathbb R,v_i}$ 平行移动到例如 $\sigma_1$ 内部的切空间，会很方便；后者典范地与 $w_1^\perp$ 识别。取 $w_1^\perp$ 的基为向量 $e_1=(1,0,0,0)$、$e_2=(0,-1,0,0)$、$e_3=(0,0,0,1)$。在这组基中，$T_{\gamma_1},T_{\gamma_2},T_{\gamma_3}$ 的矩阵分别为
\[
\begin{pmatrix}1&0&1\\0&1&0\\0&0&1\end{pmatrix},\qquad
\begin{pmatrix}1&0&-1\\0&1&-1\\0&0&1\end{pmatrix},\qquad
\begin{pmatrix}1&0&0\\0&1&1\\0&0&1\end{pmatrix}.
\]
注意，$T_{\gamma_3}\circ T_{\gamma_2}\circ T_{\gamma_1}$ 为恒等变换，正如应当预期的。

\SourceParagraph{ch06-6.4.1-029}{490}
现在，看看 $\Delta$ 位于一条边上的顶点附近发生什么。比如，取 $v_1=(-1,-1,-1,-1)$、$v_2=(0,-1,-1,-1)$。它们是一条边所包含的一个区间的端点，包含于 $\Xi$ 的三个三维面 $\sigma_1,\sigma_2,\sigma_3$ 中，分别有 $w_1=(0,0,-1,0)$、$w_2=(0,0,0,-1)$、$w_3=(0,-1,0,0)$。若 $\gamma_i'$ 为\BookRef{prop:6.79}{命题6.79}中涉及 $v_1,v_2,w_i,w_{i+1}$ 的环路，则
\[
\begin{aligned}
T_{\gamma_1'}(m)&=m+\langle(0,0,-1,1),m\rangle(1,0,0,0)
\quad\bmod v_1\mathbb R,\\
T_{\gamma_2'}(m)&=m+\langle(0,1,0,-1),m\rangle(1,0,0,0)
\quad\bmod v_1\mathbb R,\\
T_{\gamma_3'}(m)&=m+\langle(0,-1,1,0),m\rangle(1,0,0,0)
\quad\bmod v_1\mathbb R.
\end{aligned}
\]
在 $M_{\mathbb R}/v_1\mathbb R$ 的基 $e_1'=(1,0,0,0)$、$e_2'=(0,1,0,0)$、$e_3'=(0,0,-1,0)$ 中，这些单值变换的矩阵分别为
\[
\begin{pmatrix}1&0&1\\0&1&0\\0&0&1\end{pmatrix},\qquad
\begin{pmatrix}1&1&0\\0&1&0\\0&0&1\end{pmatrix},\qquad
\begin{pmatrix}1&-1&-1\\0&1&0\\0&0&1\end{pmatrix}.
\]
这两种行为有根本区别。在第一种情形中，单值变换共同逐点固定 $w_1^\perp$ 的一个二维子空间；在第二种情形中，共同逐点固定的子空间只有一维。这些点将具有很不相同的性质。不过，注意，绕 $\Delta$ 的一条边的简单环路的单值变换，始终具有相同的形式；更换基以后，它总能写成 $\begin{pmatrix}1&1&0\\0&1&0\\0&0&1\end{pmatrix}$，与二维情形非常相似。\BookQED
\end{example}

\SourceParagraph{ch06-6.4.1-030}{491}
\BookStatementNumber{83}
\begin{exercise}\label{exer:6.83}
令 $\tau\in\mathcal P$，$\tau\subseteq B$。证明：对于 $\tau$ 的任意两个顶点 $v_1,v_2$，以及 $\Xi$ 中包含 $\tau$、确定环路 $\gamma$ 的任意两个 $(n-1)$ 维面 $\sigma_1,\sigma_2$，单值变换 $T_\gamma\colon M_{\mathbb R}/v_1\mathbb R\to M_{\mathbb R}/v_1\mathbb R$，固定 $\tau$ 所张成仿射空间的切空间。这样，就把上述例子中的行为，放在一般背景下理解。
\end{exercise}

\subsection{紧化的局部模型}\label{sec:6.4.2}

\SourceParagraph{ch06-6.4.2-001}{491}
若 $B$ 为带奇点的整仿射流形，则得到对偶环面丛 $X(B_0)\to B_0$ 与 $\check X(B_0)\to B_0$。如果愿意，可以像\BookSectionRef{sec:6.2.1}{6.2.1}与\BookSectionRef{sec:6.2.3}{6.2.3}中一样，为它们添加附加数据，不过，这与下面的讨论并无特别关联。这里，$X(B_0)$ 为复流形，$\check X(B_0)$ 为辛流形。于是，可以问，能否在以下意义下紧化 $X(B_0)$ 或 $\check X(B_0)$：若 $B$ 紧，是否存在紧的 $X(B)$，以及交换图
\[
\BookImage{02_Unnumbered_Diagrams/Chapter_06/Ch06_Diagram_PDFp491_01_Bookp481.pdf}
\]
对于 $\check X(B)$，也可以提出类似的图式。这个问题最简单的形式，是一个拓扑问题，只要求 $X(B)$ 或 $\check X(B)$ 分别为 $X(B_0)$ 或 $\check X(B_0)$ 的拓扑紧化。不过，还可以进一步细化问题，要求 $X(B)$ 在复流形的意义下紧化 $X(B_0)$，并要求 $\check X(B)$ 在辛流形的意义下紧化 $\check X(B_0)$。

\SourceParagraph{ch06-6.4.2-002}{492}
\BookStatementNumber{84}
\begin{note}[numbered]\label{note:6.84}
这里有一个非常重要的要点。在所关心的情形中，并不预期 $X(B_0)$ 作为复流形能够直接紧化，但预期 $\check X(B_0)$ 作为辛流形能够紧化。在三维情形中，Castaño-Bernard 与 Matessi\BookCite{ref090}已经得到关于辛紧化的结果。稍后，在\BookSectionRef{sec:7.3.7}{7.3.7}中将看到，紧化以前，需要先形变复结构。人们普遍认为，正是辛的一侧，即 A 模型，与复的一侧，即 B 模型，之间的这种差别，产生了瞬子修正。后面还将进一步考察这一点。\BookQED
\end{note}

\SourceParagraph{ch06-6.4.2-003}{492}
与\BookRef{def:6.17}{定义6.17}中一样，若构造出 $X(B)$ 与 $\check X(B)$，则可以称它们互为 SYZ 对偶。Strominger--Yau--Zaslow 猜想，进而提示，互为 SYZ 对偶的流形，也互为镜像对偶。

\SourceParagraph{ch06-6.4.2-004}{492}
将在二维与三维情形中，非常详细地讨论这些紧化。在三维情形中研究这一问题的一项重要益处，是镜像对称中的某些基本现象，例如 Euler 示性数改变符号，在底空间 $B$ 的纯局部层面就变得清楚。本小节末尾会看到这一点。本小节将讨论这些紧化的局部模型。首先，研究\BookRef{ex:6.74}{例题6.74}(1)中的二维情形；然后，研究\BookRef{ex:6.74}{例题6.74}(2)中两种不同顶点附近的三维情形。

\SourceParagraph{ch06-6.4.2-005}{492}
最基本的构件，是 NUT 局部模型，即 nut\footnote{术语 nut 与 bolt 的讨论，见\BookSectionRef{sec:7.1.4}{7.1.4}。}。对于这里的目的，可以如下看待它。令 $\overline X=\mathbb C^2$、$X=\mathbb C^2\setminus\{(0,0)\}$、$\overline Y=\mathbb R^3$、$Y=\mathbb R^3\setminus\{(0,0,0)\}$。考虑由下式给出的映射 $p\colon X\to Y$：
\begin{equation}\label{eq:6.6}
p(z_1,z_2)=\bigl(2\operatorname{Re}(z_1z_2),\,
2\operatorname{Im}(z_1z_2),\,|z_1|^2-|z_2|^2\bigr).
\end{equation}
它将 $X$ 表示为 $Y$ 上的 $S^1$ 丛，其纤维为 $S^1$ 作用
\[
(z_1,z_2)\longmapsto(e^{\mathrm i\vartheta}z_1,e^{-\mathrm i\vartheta}z_2)
\]
的轨道。可以如下计算这一丛的第一 Chern 类。选择丛上的联络1形式；可以取
\[
\theta=\mathrm i\,
\frac{\operatorname{Im}(\bar z_1\,\mathrm dz_1-\bar z_2\,\mathrm dz_2)}
{|z_1|^2+|z_2|^2}.
\]
那么，按照 $c_1=[\mathrm i F/(2\pi)]$ 的约定，这个丛的第一 Chern 类，由曲率2形式
\[
\frac{\mathrm i\,\mathrm d\theta}{2\pi}
=\frac{u_1\,\mathrm du_2\wedge\mathrm du_3
+u_2\,\mathrm du_3\wedge\mathrm du_1
+u_3\,\mathrm du_1\wedge\mathrm du_2}
{4\pi(u_1^2+u_2^2+u_3^2)^{3/2}}
\]
代表，其中，$u_1,u_2,u_3$ 为 $\mathbb R^3$ 上的坐标。要确定这个2形式所代表的上同调类，只需在 $\mathbb R^3$ 中按外法向定向的单位二维球面上积分，得到
\[
\int_{S^2}\frac{\mathrm i\,\mathrm d\theta}{2\pi}=1.
\]
因此，丛的第一 Chern 类为 $1\in\mathbb Z\cong H^2(Y,\mathbb Z)$。这里，$\overline X$ 就是 NUT 模型；与 $p$ 由同一公式给出的 $\overline p\colon\overline X\to\overline Y$，是 $S^1$ 丛 $X\to Y$ 的部分紧化。一般而言，它按如下方式使用。

\SourceParagraph{ch06-6.4.2-006}{493}
令 $\overline Y$ 为任意三维流形，不必紧，令 $Y\subseteq\overline Y$ 为开集，$\overline Y\setminus Y$ 为离散点集。设给定 $S^1$ 丛 $q\colon X\to Y$，具有如下性质：对于每个 $y\in\overline Y\setminus Y$，$y$ 在 $\overline Y$ 中有开邻域 $U_y$，使得 $U_y\setminus\{y\}\cong\mathbb R^3\setminus\{0\}$，并且 $q^{-1}(U_y\setminus\{y\})\to U_y\setminus\{y\}$ 是第一 Chern 类为 $\pm1$ 的 $S^1$ 丛。那么，可以把 NUT 模型粘接到 $q^{-1}(U_y\setminus\{y\})$ 中：有交换图
\[
\BookImage{02_Unnumbered_Diagrams/Chapter_06/Ch06_Diagram_PDFp493_01_Bookp483.pdf}
\]
因此，可以粘入 $\mathbb C^2$，得到 $\overline q\colon\overline X\to\overline Y$。

\SourceParagraph{ch06-6.4.2-007}{493}
下面的练习说明，所得空间 $\overline X$ 不依赖于具体的粘合方式。

\SourceParagraph{ch06-6.4.2-008}{493}
\BookStatementNumber{85}
\begin{exercise}\label{exer:6.85}
令 $Y,\overline Y$ 与 $q\colon X\to Y$ 如上。证明：集合 $\overline X=X\sqcup(\overline Y\setminus Y)$ 上有唯一的拓扑，延伸 $X$ 上的拓扑，使得存在连续映射的交换图
\[
\BookImage{02_Unnumbered_Diagrams/Chapter_06/Ch06_Diagram_PDFp493_02_Bookp483.pdf}
\]
其中，$\overline q$ 为固有映射，即紧集的逆像为紧集，并且 $\overline q|_{\overline Y\setminus Y}$ 是 $\overline Y\setminus Y$ 上的恒等映射。此外，$\overline X$ 为流形。
\end{exercise}

\SourceParagraph{ch06-6.4.2-009}{493}
现在，可以构造关键的二维例子。

\SourceParagraph{ch06-6.4.2-010}{493}
\BookStatementNumber{86}
\begin{construction}\label{constr:6.86}
令 $\overline Y=S^1\times D$，其中，$D$ 为二维开圆盘，选定一点 $0\in D$。还选择点 $p\in S^1$，并令 $Y=\overline Y\setminus\{(p,0)\}$。现在，可以计算 $H^2(Y,\mathbb Z)$，从而确定 $Y$ 上 $S^1$ 丛的 Chern 类。注意，$H^2(\overline Y,\mathbb Z)=0$。由相对上同调长正合列，得到
\[
0=H^2(\overline Y,\mathbb Z)\longrightarrow H^2(Y,\mathbb Z)
\longrightarrow H^3(\overline Y,Y,\mathbb Z)
\longrightarrow H^3(\overline Y,\mathbb Z)=0.
\]
可以如下计算 $H^3(\overline Y,Y,\mathbb Z)$。$(p,0)$ 有邻域 $U$，同胚于 $\mathbb R^3$，并且 $U\setminus\{(p,0)\}$ 同胚于 $\mathbb R^3\setminus\{0\}$。由切除性，
\[
H^3(\overline Y,Y,\mathbb Z)\cong H^3(U,U\setminus\{(p,0)\},\mathbb Z).
\]
不过，相对上同调长正合列又给出
\[
H^3(U,U\setminus\{(p,0)\},\mathbb Z)
\cong H^2(U\setminus\{(p,0)\},\mathbb Z)\cong\mathbb Z,
\]
因为 $\mathbb R^3\setminus\{0\}$ 形变收缩到 $S^2$。因此，得到结论 $H^2(Y,\mathbb Z)\cong\mathbb Z$。进一步，取 $S^1$ 丛 $q\colon X\to Y$，第一 Chern 类为1；这里的生成元与下面的粘合约定相容。于是，可以添加一个 NUT 模型，得到 $\overline q\colon\overline X\to\overline Y$。再将 $\overline q$ 与投影 $\overline Y\to D$ 复合，得到 $f\colon\overline X\to D$。

\SourceParagraph{ch06-6.4.2-011}{494}
来考察 $f$ 的纤维。若 $x\in D$、$x\ne0$，则 $f^{-1}(x)$ 为 $S^1\times\{x\}$ 上的 $S^1$ 丛，即二维环面。不过，$f^{-1}(0)$ 看起来像 $T^2$ 的一个圆圈被夹缩为一点，见\BookRef{fig:6.5}{图5}。
\begin{figure}[H]
\centering
\BookImage{01_Numbered_Figures/Chapter_06/Ch06_Fig05_PDFp494_Bookp484.pdf}
\BookNumberOnlyCaption\label{fig:6.5}
\end{figure}
\BookRef{fig:6.5}{图5}中的阴影圆盘，按如下方式得到。在 $D$ 中取一条闭线段，一端为0，另一端为 $b\in D$。把这条线段提升为 $\overline Y$ 中的一条线段，其中一端为 $(p,0)$。然后，在 $\overline X$ 中取这条线段的逆像，得到一个圆盘，其边界为包含于 $f^{-1}(b)$ 中的 $S^1$。这个 $S^1$ 给出 $H_1(f^{-1}(b),\mathbb Z)$ 中良定义的闭链，通常称为消灭闭链。在 $\overline X$ 中，它同调于零，因此得名。

\SourceParagraph{ch06-6.4.2-012}{495}
此时，很有必要理解单值变换在这里的作用。环面丛 $f\colon f^{-1}(D\setminus\{0\})\to D\setminus\{0\}$ 并非平凡丛。若 $b\in D\setminus\{0\}$ 为基点，则如下得到单值变换 $T\colon H_1(f^{-1}(b),\mathbb Z)\to H_1(f^{-1}(b),\mathbb Z)$。令 $\gamma\colon[0,1]\to D\setminus\{0\}$ 为以 $b$ 为基点的简单逆时针环路。沿 $\gamma$ 拉回纤维化 $f\colon\overline X\to D$，得到族 $X'\to[0,1]$。由于 $[0,1]$ 可缩，可以将其平凡化，即存在与映射 $X'\to[0,1]$ 相容的同构 $\psi\colon X'\cong[0,1]\times T^2$。利用 $X_0'=X_1'$，得到 $T^2\cong X_0'$ 的微分同胚 $\varphi$，由复合
\[
X_0'\xrightarrow{\psi}\{0\}\times T^2
=\{1\}\times T^2\xrightarrow{\psi^{-1}}X_1'
\]
给出。若 $S^1\subseteq D\setminus\{0\}$ 为 $\gamma$ 的像，则为了在拓扑上恢复纤维化 $f^{-1}(S^1)\to S^1$，通过 $\varphi$ 将 $\{0\}\times T^2$ 与 $\{1\}\times T^2$ 粘合。$\varphi$ 通常称为单值变换微分同胚。由于平凡化并不唯一，单值变换微分同胚只在相差同痕的意义下确定。不过，$\varphi$ 在同调上诱导唯一映射 $\varphi_*\colon H_1(f^{-1}(b),\mathbb Z)\to H_1(f^{-1}(b),\mathbb Z)$，这就是单值变换。还可以把它简单地看作这样一个映射：取 $f^{-1}(b)$ 中的闭链，沿环路前进时连续地追踪它。绕行以后，未必回到原来的闭链。

\SourceParagraph{ch06-6.4.2-013}{495}
在这个例子中，可以非常明确地看到这一点。考虑 $C\subseteq D\setminus\{0\}$，绕0一周的圆，并令 $Z=f^{-1}(C)$。按照构造，$Z$ 为 $T^2=C\times S^1\subseteq Y$ 上的 $S^1$ 丛，其 $c_1$ 是 $1\in H^2(Y,\mathbb Z)$ 在 $H^2(C\times S^1,\mathbb Z)$ 中的限制。容易看到，在相应生成元下，限制映射为恒等映射。沿覆盖 $C\times\mathbb R\to C\times S^1$，将这一 $S^1$ 丛拉回到 $C\times\mathbb R$ 后，它变得平凡，即为投影 $C\times\mathbb R\times S^1\to C\times\mathbb R$。对 $C\times\mathbb R\times S^1$ 按 $\mathbb Z$ 作用取商
\[
(e^{2\pi\mathrm i\theta},\rho,e^{2\pi\mathrm i\tau})
\longmapsto(e^{2\pi\mathrm i\theta},\rho+1,e^{2\pi\mathrm i(\theta+\tau)}),
\]
就可以恢复 $S^1$ 丛 $Z\to C\times S^1$。这里，将 $C$ 与 $S^1$ 和 $\operatorname{U}(1)$ 识别。于是，容易看到，投影 $(C\times\mathbb R\times S^1)/\mathbb Z\to C\times S^1$，是 $c_1=1$ 的 $S^1$ 丛；在本段中，选取底环面的定向为 $\mathrm d\rho\wedge\mathrm d\theta$。将这一点留作练习。现在，考虑映射 $\gamma\colon[0,1]\to C\subseteq D$，如前面一样，以点 $y\in D$ 为基点参数化环路 $C$，并沿 $\gamma$ 拉回 $f\colon\overline X\to D$，得到 $X'\to[0,1]$。由上述 $Z$ 的描述，看到 $X'\cong[0,1]\times S^1\times S^1$，具体地，通过同构
\[
([0,1]\times\mathbb R\times S^1)/\mathbb Z
\longrightarrow[0,1]\times S^1\times S^1,
\qquad
(\theta,\rho,e^{2\pi\mathrm i\tau})
\longmapsto(\theta,e^{2\pi\mathrm i\rho},e^{2\pi\mathrm i(\tau-\theta\rho)})
\]
给出。因此，单值变换微分同胚为
\[
(e^{2\pi\mathrm i\rho},e^{2\pi\mathrm i\tau})
\longmapsto(e^{2\pi\mathrm i\rho},e^{2\pi\mathrm i(\rho+\tau)}).
\]

\SourceParagraph{ch06-6.4.2-014}{496}
这是 Dehn 扭转的一个例子，可以如下理解。取环面，沿一个圆圈切开，如图所示。将所得圆柱扭转 $360$ 度，然后重新粘合，如\BookRef{fig:6.6}{图6}所示。
\begin{figure}[H]
\centering
\BookImage{01_Numbered_Figures/Chapter_06/Ch06_Fig06_PDFp496_Bookp486.pdf}
\BookNumberOnlyCaption\label{fig:6.6}
\end{figure}
若 $e_1,e_2$ 分别为 $\rho=\text{常数}$ 与 $\tau=\text{常数}$ 给出的 $H_1(f^{-1}(y),\mathbb Z)$ 生成元，则对应的同调变换 $T\colon H_1(f^{-1}(y),\mathbb Z)\to H_1(f^{-1}(y),\mathbb Z)$，在这组基中由矩阵 $\begin{pmatrix}1&1\\0&1\end{pmatrix}$ 给出。注意，这里被单值变换固定的闭链，就是消灭闭链。
\end{construction}

\SourceParagraph{ch06-6.4.2-015}{496}
\BookStatementNumber{87}
\begin{construction}\label{constr:6.87}
还有另一种方式，构造\BookRef{constr:6.86}{构造6.86}中的 $T^2$ 纤维化，而且能够很好地推广到更高维。考虑 $\overline X\subseteq\mathbb C^2\times\mathbb C^*$，其中，$z_1,z_2$ 为 $\mathbb C^2$ 上的坐标，$w$ 为 $\mathbb C^*$ 上的坐标，$\overline X$ 由方程 $z_1z_2=w-1$ 定义。定义 $f\colon\overline X\to\mathbb R^2$ 为
\[
f(z_1,z_2,w)=(|z_1|^2-|z_2|^2,\ln|w|).
\]
它通过映射 $\overline q\colon\overline X\to\overline Y=\mathbb R\times\mathbb C^*\cong D\times S^1$ 分解，其中，$\overline q(z_1,z_2,w)=(|z_1|^2-|z_2|^2,w)$。容易证明，若 $X=\overline X\setminus\{(0,0,1)\}$、$Y=\overline Y\setminus\{(0,1)\}$，则 $q=\overline q|_X\colon X\to Y$ 为 $S^1$ 丛，$S^1$ 在纤维上的作用为
\[
(z_1,z_2,w)\longmapsto(e^{\mathrm i\theta}z_1,e^{-\mathrm i\theta}z_2,w).
\]
显然，$\overline X$ 是通过 NUT 模型对 $X$ 作的紧化，而纤维化 $f\colon\overline X\to\mathbb R^2\cong D$，在拓扑上等价于\BookRef{constr:6.86}{构造6.86}中的纤维化。
\end{construction}

\SourceParagraph{ch06-6.4.2-016}{496}
接下来，考察三维紧化所需的局部模型。正是在三维情形中，才能在这一拓扑图景里真正看到镜像对称。

\SourceParagraph{ch06-6.4.2-017}{496}
\BookStatementNumber{88}
\begin{construction}\label{constr:6.88}
令 $f\colon\overline X\to D$ 为\BookRef{constr:6.86}{构造6.86}中的 $T^2$ 纤维化。那么，由 $(x,p,r)\mapsto(f(x),r)$，得到 $T^3$ 纤维化
\[
\overline X\times S^1\times(0,1)\longrightarrow D\times(0,1).
\]
奇异纤维位于 $\{0\}\times(0,1)$ 上，形如 $f^{-1}(0)\times S^1$，其中，$f^{-1}(0)$ 为夹缩环面。绕 $\Delta=\{0\}\times(0,1)$ 的简单环路，在适当的基中，诱导单值变换 $\begin{pmatrix}1&1&0\\0&1&0\\0&0&1\end{pmatrix}$。或者，像\BookRef{constr:6.87}{构造6.87}中一样，可以描述 $\overline X\subseteq\mathbb C^2\times(\mathbb C^*)^2$，其中，$z_1,z_2$ 为 $\mathbb C^2$ 上的坐标，$w_1,w_2$ 为 $(\mathbb C^*)^2$ 上的坐标；$\overline X$ 由方程 $z_1z_2=w_1-1$ 定义，$\overline X\to\mathbb R^3$ 则由 $(|z_1|^2-|z_2|^2,\ln|w_1|,\ln|w_2|)$ 定义。
\end{construction}

\SourceParagraph{ch06-6.4.2-018}{497}
\BookStatementNumber{89}
\begin{construction}\label{constr:6.89}
可以如下将 NUT 模型推广到更高维。令 $\overline Y$ 为流形，$Y\subseteq\overline Y$ 为开子集，$\overline Y\setminus Y=S$ 为余维三的子流形。考虑 $S^1$ 丛 $q\colon X\to Y$。假定在 $S$ 附近局部地，这个 $S^1$ 丛的第一 Chern 类为1。这里的意义如下：对于每个点 $y\in S$，$y$ 在 $\overline Y$ 中有开邻域 $U$，同胚于 $B^{n-3}\times\mathbb R^3$；其中，$B^{n-3}$ 为 $(n-3)$ 维开球，并且 $S\cap U=B^{n-3}\times\{0\}$。要求 $S^1$ 丛在 $U\setminus(S\cap U)$ 上的限制，具有第一 Chern 类
\[
1\in H^2(B^{n-3}\times(\mathbb R^3\setminus\{0\}),\mathbb Z)=\mathbb Z.
\]
给定这些数据，$\overline X=X\sqcup S$ 上有唯一的拓扑，延伸 $X$ 上的拓扑，使得 $\overline X$ 为拓扑流形，并且存在连续映射的交换图
\[
\BookImage{02_Unnumbered_Diagrams/Chapter_06/Ch06_Diagram_PDFp497_01_Bookp487.pdf}
\]
其中，$\overline q$ 为固有映射，$\overline q|_S$ 为恒等映射。为说明这一点，注意，取 $X$ 中的开集，以及 $\overline Y$ 中开集在 $\overline q$ 下的逆像，作为 $\overline X$ 拓扑的一组基，就容易看到，这是使 $\overline q$ 固有的唯一拓扑。

\SourceParagraph{ch06-6.4.2-019}{497}
现在，若 $y\in S$ 为一点，令 $U\subseteq\overline Y$ 为上述 $y$ 的开邻域，$U\cong B^{n-3}\times\mathbb R^3$。那么，对于 $B^{n-3}$ 中的任意点 $y'$，$q$ 在 $\{y'\}\times(\mathbb R^3\setminus\{0\})$ 上的限制，是第一 Chern 类为 $\pm1$ 的 $S^1$ 丛。因此，利用 NUT 模型，有交换图
\[
\BookImage{02_Unnumbered_Diagrams/Chapter_06/Ch06_Diagram_PDFp497_02_Bookp487.pdf}
\]
其中，$\overline q'$ 在第一个因子上为恒等映射，在第二个因子上为\BookEqRef{eq:6.6}{6.6}中的映射 $\overline p$。因此，由紧化的唯一性，$\overline q'$ 与 $\overline q^{-1}(U)\to U$ 相同。特别地，$\overline X$ 为流形。
\end{construction}

\SourceParagraph{ch06-6.4.2-020}{497}
\BookStatementNumber{90}
\begin{construction}\label{constr:6.90}
令 $B$ 为三维球，令 $\Delta\subseteq B$，$\Delta=l_1\cup l_2\cup l_3\cup\{0\}$；其中，$l_1,l_2,l_3$ 为开线段，如\BookRef{fig:6.7}{图7}所示。令 $\overline Y=B\times(\mathbb R^2/\mathbb Z^2)$。将描述一个特定曲面 $S\subseteq\overline Y$。令 $C_1$ 为圆圈 $(1,0)\mathbb R/(1,0)\mathbb Z\subseteq\mathbb R^2/\mathbb Z^2$，令 $C_2$ 为圆圈 $(0,1)\mathbb R/(0,1)\mathbb Z\subseteq\mathbb R^2/\mathbb Z^2$。令 $S_1=\overline l_1\times C_1$、$S_2=\overline l_2\times C_2$，为包含于 $\overline Y$、分别位于 $l_1,l_2$ 的闭包 $\overline l_1,\overline l_2$ 上方的圆柱。
\begin{figure}[H]
\centering
\BookImage{01_Numbered_Figures/Chapter_06/Ch06_Fig07_PDFp498_Bookp488.pdf}
\BookNumberOnlyCaption\label{fig:6.7}
\end{figure}

\SourceParagraph{ch06-6.4.2-021}{498}
如下构造 $S_3$。在圆圈 $C_3=(1,-1)\mathbb R/(1,-1)\mathbb Z$ 与 $C_1-C_2$ 之间，选择一个足够良好的同伦。也就是说，取连续映射 $H\colon S^1\times[0,1]\to\mathbb R^2/\mathbb Z^2$，满足 $H(S^1\times\{1\})=C_3$、$H(S^1\times\{0\})=C_1\cup C_2$。所谓足够良好，是指：若取参数化 $\alpha\colon[0,1)\to\overline l_3$，$\alpha(0)=0$，则
\[
S_3=\{(\alpha(t),H(\theta,t))\in\overline Y
\mid\theta\in S^1,\ 0\leq t<1\}\subseteq\overline Y
\]
在 $\{0\}\times(\mathbb R^2/\mathbb Z^2)$ 以外为光滑曲面。令 $S=S_1\cup S_2\cup S_3$。这三个曲面拼接在一起，给出在 $\Delta$ 上纤维化的裤形曲面，如\BookRef{fig:6.8}{图8}所示。注意，无法使 $S$ 成为光滑曲面，因此，这将是一个纯拓扑构造。

\SourceParagraph{ch06-6.4.2-022}{498}
现在，令 $Y=\overline Y\setminus S$。可以尝试如下计算 $H^2(Y,\mathbb Z)$。有相对上同调长正合列
\[
\begin{aligned}
H^2(\overline Y,Y,\mathbb Z)&\longrightarrow H^2(\overline Y,\mathbb Z)
\longrightarrow H^2(Y,\mathbb Z)\\
&\longrightarrow H^3(\overline Y,Y,\mathbb Z)
\longrightarrow H^3(\overline Y,\mathbb Z).
\end{aligned}
\]
$\overline Y$ 形变收缩到 $T^2=\mathbb R^2/\mathbb Z^2$，所以，$H^2(\overline Y,\mathbb Z)=H^2(T^2,\mathbb Z)$，$H^3(\overline Y,\mathbb Z)=0$。此外，由 Alexander 对偶，$H^p(\overline Y,Y,\mathbb Z)=H^{p-3}(S,\mathbb Z)$，所以
\[
H^2(\overline Y,Y,\mathbb Z)=0,
\qquad
H^3(\overline Y,Y,\mathbb Z)=H^0(S,\mathbb Z)=\mathbb Z.
\]
这给出短正合列
\[
0\longrightarrow H^2(\overline Y,\mathbb Z)
\longrightarrow H^2(Y,\mathbb Z)\longrightarrow\mathbb Z\longrightarrow0.
\]
最后，对于 $b\in B\setminus\Delta$，限制映射 $H^2(Y,\mathbb Z)\to H^2(\{b\}\times(\mathbb R^2/\mathbb Z^2),\mathbb Z)$ 分裂上述正合列。因此，得到典范同构 $H^2(Y,\mathbb Z)\cong H^2(\overline Y,\mathbb Z)\oplus\mathbb Z$。
\begin{figure}[H]
\centering
\BookImage{01_Numbered_Figures/Chapter_06/Ch06_Fig08_PDFp499_Bookp489.pdf}
\BookNumberOnlyCaption\label{fig:6.8}
\end{figure}

在这一识别下，考虑 $H^2(Y,\mathbb Z)$ 中由 $(0,1)$ 给出的元素 $c_1$。事实上，若 $y\in S$，$U$ 为\BookRef{constr:6.89}{构造6.89}中那样的 $y$ 的小邻域，则 $c_1$ 在 $U\setminus S$ 上的限制，为 $1\in H^2(U\setminus S,\mathbb Z)=\mathbb Z$。于是，若 $q\colon X\to Y$ 为第一 Chern 类为 $c_1$ 的 $S^1$ 丛，则 $q$ 满足\BookRef{constr:6.89}{构造6.89}的条件，可以紧化为六维流形 $\overline X$，得到映射 $\overline q\colon\overline X\to\overline Y$。再与到 $B$ 的投影复合，得到映射 $f\colon\overline X\to B$。

\SourceParagraph{ch06-6.4.2-023}{499}
这一构造得到了什么？由于对于 $b\in B\setminus\Delta$，$c_1$ 在每个 $\{b\}\times T^2$ 上都限制为零，$f^{-1}(b)$ 为 $T^2$ 上的平凡 $S^1$ 丛，即 $T^3$。因此，$f$ 为 $T^3$ 纤维化。此外，若 $b\in l_i$，则 $f^{-1}(b)$ 形如 $F\times S^1$，其中，$F$ 为夹缩环面。最后，$f^{-1}(0)$ 为一个退化程度更高的纤维，奇点集是“8”字形，因为 $(\{0\}\times(\mathbb R^2/\mathbb Z^2))\cap S$ 为“8”字形。注意，这个奇异纤维的拓扑 Euler 示性数为 $-1$。因此，称这个纤维为负型纤维。
\end{construction}

\SourceParagraph{ch06-6.4.2-024}{499}
\BookStatementNumber{91}
\begin{exercise}\label{exer:6.91}
在\BookRef{constr:6.90}{构造6.90}中，计算环面丛 $f^{-1}(B\setminus\Delta)\to B\setminus\Delta$ 的单值变换。将其与\BookRef{ex:6.82}{例题6.82}(2)中计算的仿射结构的单值变换比较。这个例子对应于哪一类顶点？
\end{exercise}

\SourceParagraph{ch06-6.4.2-025}{499}
\BookStatementNumber{92}
\begin{construction}\label{constr:6.92}
令 $\overline X\subseteq\mathbb C^3\times\mathbb C^*$，其中，$z_1,z_2,z_3$ 为 $\mathbb C^3$ 上的坐标，$w$ 为 $\mathbb C^*$ 上的坐标，$\overline X$ 由方程 $z_1z_2z_3=w-1$ 定义。考虑映射 $f\colon\overline X\to\mathbb R^3$，定义为
\[
f(z_1,z_2,z_3,w)=(|z_1|^2-|z_2|^2,\,|z_1|^2-|z_3|^2,\,\log|w|).
\]

\SourceParagraph{ch06-6.4.2-026}{500}
现在，分析这个映射。首先，注意，它分解为
\[
\overline X\xrightarrow{\overline q}\overline Y\xrightarrow{g}\mathbb R^3,
\]
其中，$\overline Y=\mathbb R^2\times\mathbb C^*$，
\[
\overline q(z_1,z_2,z_3,w)=(|z_1|^2-|z_2|^2,\,|z_1|^2-|z_3|^2,\,w),
\qquad
g(x_1,x_2,w)=(x_1,x_2,\log|w|).
\]

\SourceParagraph{ch06-6.4.2-027}{500}
映射 $g$，使得可以将 $\overline Y$ 与 $\mathbb R^3\times S^1$ 识别。同时，$\overline q$ 在 $T^2$ 作用
\[
(z_1,z_2,z_3,w)\longmapsto
(e^{\mathrm i\theta_1}z_1,e^{\mathrm i\theta_2}z_2,
 e^{-\mathrm i(\theta_1+\theta_2)}z_3,w)
\]
下不变。容易看到，$\overline q$ 的每个纤维，都是这个 $T^2$ 作用的一条轨道。不过，有一些退化轨道。令
\[
S=\{(z_1,z_2,z_3,w)\in\overline X\mid
z_1,z_2,z_3\text{ 中至少两个为零，且 }w=1\}.
\]
若 $y\in S\setminus\{(0,0,0,1)\}$，则 $y$ 的 $T^2$ 轨道为圆圈；若 $y=(0,0,0,1)$，则 $y$ 的 $T^2$ 轨道为一点。因此，特别地，若 $X=\overline X\setminus S$、$Y=\overline Y\setminus\overline q(S)$，则 $\overline q$ 在 $X$ 上的限制 $q\colon X\to Y$，为 $T^2$ 主丛。注意，
\[
\begin{aligned}
\overline q(S)
&=\{(0,x_2,1)\in\overline Y\mid x_2\leq0\}\\
&\qquad\cup\{(x_1,0,1)\in\overline Y\mid x_1\leq0\}\\
&\qquad\cup\{(x_1,x_2,1)\in\overline Y\mid x_1=x_2\geq0\}.
\end{aligned}
\]
令 $\Delta=f(S)$。可以写成 $\Delta=l_1\cup l_2\cup l_3\cup\{(0,0,0)\}$，其中，
\[
\begin{aligned}
l_1&=\{(0,x_2,0)\in\mathbb R^3\mid x_2<0\},\\
l_2&=\{(x_1,0,0)\in\mathbb R^3\mid x_1<0\},\\
l_3&=\{(x_1,x_2,0)\in\mathbb R^3\mid x_1=x_2>0\}.
\end{aligned}
\]

\SourceParagraph{ch06-6.4.2-028}{500}
对于点 $y\in\mathbb R^3\setminus\Delta$，$g^{-1}(y)=S^1$，所以，$f^{-1}(y)$ 是 $S^1$ 上的 $T^2$ 主丛，即三维环面。因此，$f$ 为 $T^3$ 纤维化，在 $\Delta$ 上有退化纤维。希望更仔细地分析一般奇异纤维附近的行为，也就是 $l_i$ 上的情形，$i=1,2,3$。为此，令 $U_i$ 为 $\mathbb R^3$ 的开集，同胚于 $l_i\times D$；其中，$D$ 为二维圆盘，要求 $l_i\subseteq U_i$、$(0,0,0)\notin U_i$，且 $i\ne j$ 时，$U_i\cap l_j=\varnothing$。见\BookRef{fig:6.9}{图9}。可以研究 $f^{-1}(U_i)\to U_i$。注意，
\[
f^{-1}(U_i)\cap S=
\begin{cases}
\{(z_1,z_2,z_3,1)\mid z_1=z_2=0,\ z_3\ne0\},&i=1,\\
\{(z_1,z_2,z_3,1)\mid z_1=z_3=0,\ z_2\ne0\},&i=2,\\
\{(z_1,z_2,z_3,1)\mid z_2=z_3=0,\ z_1\ne0\},&i=3.
\end{cases}
\]
例如，取 $i=1$。可以将商映射
\[
\overline q\colon f^{-1}(U_1)\longrightarrow
f^{-1}(U_1)/T^2=g^{-1}(U_1)=U_1\times S^1
\]
按以下方式分解。
\begin{figure}[H]
\centering
\BookImage{01_Numbered_Figures/Chapter_06/Ch06_Fig09_PDFp501_Bookp491.pdf}
\BookNumberOnlyCaption\label{fig:6.9}
\end{figure}

\SourceParagraph{ch06-6.4.2-029}{501}
先对 $f^{-1}(U_1)$ 按 $S^1$ 作用 $(z_1,z_2,z_3)\mapsto(e^{\mathrm i\theta}z_1,e^{-\mathrm i\theta}z_2,z_3)$ 取商，得到
\[
f^{-1}(U_1)\xrightarrow{\overline q_1}f^{-1}(U_1)/S^1
\xrightarrow{\overline q_2}U_1\times S^1.
\]
现在，$\overline q_2$ 为必然平凡的 $S^1$ 丛，因为剩余的 $T^2/S^1$ 作用在 $f^{-1}(U_1)/S^1$ 上为自由作用。因此，$f^{-1}(U_1)/S^1\cong U_1\times S^1\times S^1=U_1\times T^2$，而 $\overline q_1\colon f^{-1}(U_1)\to U_1\times T^2$，是沿包含于 $U_1\times T^2$ 的曲面 $\overline q_1(S\cap f^{-1}(U_1))$ 退化的 $S^1$ 丛；具体地，这个曲面为
\[
l_1\times\{\text{一点}\}\times S^1\subseteq U_1\times S^1\times S^1.
\]
在 $U_2$ 与 $U_3$ 上，也得到类似的图景。
\end{construction}

\SourceParagraph{ch06-6.4.2-030}{501}
\BookStatementNumber{93}
\begin{exercise}\label{exer:6.93}
利用上述描述，证明 $f^{-1}(U_i)\to U_i$ 在拓扑上与\BookRef{constr:6.88}{构造6.88}中的纤维化相同。利用这一点，计算 $T^3$ 丛 $f^{-1}(\mathbb R^3\setminus\Delta)\to\mathbb R^3\setminus\Delta$ 绕 $l_1,l_2,l_3$ 的环路的单值变换。选择适当的基，证明它是\BookRef{constr:6.90}{构造6.90}中纤维化的单值变换的转置的逆。
\end{exercise}

\SourceParagraph{ch06-6.4.2-031}{501}
\begin{definition*}[正型纤维]\label{def:6.positive-fibre}
最后，注意，$f^{-1}(0,0,0)$ 看起来像一个 $T^3$，其中的 $T^2$ 被夹缩为一点，如\BookRef{fig:6.10}{图10}所示。这个纤维的 Euler 示性数为 $+1$，因此，称它为正型纤维。
\end{definition*}
\begin{figure}[H]
\centering
\BookImage{01_Numbered_Figures/Chapter_06/Ch06_Fig10_PDFp502_Bookp492.pdf}
\BookNumberOnlyCaption\label{fig:6.10}
\end{figure}

\SourceParagraph{ch06-6.4.2-032}{501}
现在，看到 SYZ 对偶的一种简单局部形式：正型与负型纤维互为 SYZ 对偶。其意义是，在\BookRef{constr:6.90}{构造6.90}与\BookRef{constr:6.92}{构造6.92}中，已经构造出纤维化 $X(B)\to B$ 与 $\check X(B)\to B$，其中，$B$ 为例如\BookRef{ex:6.74}{例题6.74}(2)中某个三价顶点的邻域。这些纤维化互为 SYZ 对偶。这种对偶的一个特定结果，是奇异纤维的 Euler 示性数符号相反。将此看作镜像对称的一种局部表现。特别地，这一观察，将解释互为 SYZ 对偶的环面纤维化的 Euler 示性数，在整体上改变符号的现象。

\subsection{紧化}\label{sec:6.4.3}

\SourceParagraph{ch06-6.4.3-001}{502}
可以利用前一小节描述的局部模型，在拓扑上紧化具有适当单值变换的二维与三维环面纤维化。这里不会像前一小节那样给出那么多细节，因为这些细节主要是技术性的，将它们留作练习。

\SourceParagraph{ch06-6.4.3-002}{502}
首先，令 $B$ 为带奇点的二维整仿射流形。假定对于绕 $B$ 的奇点的任意简单环路 $\gamma$，其基点 $b$ 位于奇点附近，单值变换 $T_\gamma\colon\Lambda_b\to\Lambda_b$ 在某组基中形如 $\begin{pmatrix}1&1\\0&1\end{pmatrix}$。例如，\BookRef{ex:6.82}{例题6.82}(1)中就是如此。现在，考虑环面丛 $f_0\colon X(B_0)\to B_0$。一般而言，$\Lambda_b\cong H_1(f_0^{-1}(b),\mathbb Z)$，而 $T_\gamma$ 与单值变换 $H_1(f_0^{-1}(b),\mathbb Z)\to H_1(f_0^{-1}(b),\mathbb Z)$ 相同。因此，对仿射流形的单值变换所作的限制，保证了环面丛在奇点附近的局部单值变换，与\BookRef{constr:6.86}{构造6.86}中的相同。

\SourceParagraph{ch06-6.4.3-003}{502}
\BookRef{constr:6.86}{构造6.86}给出了 $T^2$ 纤维化 $\overline X\to D$。那里，$\overline X$ 与 $D$ 都是开流形。将 $D$ 换成包含于 $D$ 内部的闭圆盘 $\overline D$，并将 $\overline X$ 相应地换成 $\overline D$ 的逆像，就得到带边界流形之间的 $T^2$ 纤维化 $\overline f\colon\overline X\to\overline D$，具有一个奇异纤维。于是，$\overline f\colon\partial\overline X\to\partial\overline D$ 是圆周上的 $T^2$ 丛，前面已经明确描述过它。若将 $\overline D$ 与 $B$ 中包含一个奇点 $b_0\in B$ 的小闭圆盘 $\overline D'\subseteq B$ 识别，就得到 $T^2$ 丛 $f_0^{-1}(\partial\overline D')\to\partial\overline D'$。

\SourceParagraph{ch06-6.4.3-004}{503}
\BookStatementNumber{94}
\begin{exercise}\label{exer:6.94}
将 $T^2$ 丛 $\overline f\colon\partial\overline X\to\partial\overline D$，描述为 $\mathcal E/\Gamma\to\partial\overline D$，其中，$\mathcal E$ 为 $\partial\overline D$ 上必然平凡的实秩2向量丛，$\Gamma$ 为 $\mathcal E$ 中的一族格。当然，按照构造，$T^2$ 丛 $f_0^{-1}(\partial\overline D')\to\partial\overline D'$ 也具有同样的形式，即
\[
TB_0|_{\partial\overline D'}/\Lambda|_{\partial\overline D'}
\longrightarrow\partial\overline D'.
\]
然后，证明：由于局部系统 $\Lambda|_{\partial\overline D'}$ 与 $\Gamma$ 具有相同的单值变换，存在同胚 $\partial\overline X\to f_0^{-1}(\partial\overline D')$，保持纤维结构，并且在每个纤维上为线性映射。
\end{exercise}

\SourceParagraph{ch06-6.4.3-005}{503}
由这一练习，可以利用上述同胚，沿共同边界，将 $\overline f\colon\overline X\to\overline D$ 与 $f_0^{-1}(B_0\setminus(\overline D'\cap B_0))\to B_0\setminus(\overline D'\cap B_0)$ 粘合。在每个奇点处作这样的粘合，就产生一个紧化后的拓扑纤维化 $X(B)\to B$。可以安排这一过程，使得所得 $X(B)$ 为可微流形，且 $X(B)\to B$ 为 $C^\infty$ 映射。

\SourceParagraph{ch06-6.4.3-006}{503}
对于 $\check X(B_0)\to B_0$，也同样可以进行这一过程，因为对偶化不改变单值变换的共轭类。事实上，在这种情形中，可以在辛范畴内完成紧化，见文献\BookCite{ref435,ref090}，不过，这里不给出细节。

\SourceParagraph{ch06-6.4.3-007}{503}
转向三维情形，仍须对纤维化的单值变换作一个假设。将为这一限制给出一个比较专门的定义。这是文献\BookCite{ref206}第1.5节中一个更一般定义的特殊情形。

\SourceParagraph{ch06-6.4.3-008}{503}
\BookStatementNumber{95}
\begin{definition}\label{def:6.95}
令 $B$ 为带奇点的三维整仿射流形，判别集 $\Delta$ 为图。称 $B$ 为简单的，是指：
\begin{enumerate}[label=(\arabic*)]
\item $\Lambda$ 绕图 $\Delta$ 的每条边的单值变换，在适当的基中具有形式 $\begin{pmatrix}1&1&0\\0&1&0\\0&0&1\end{pmatrix}$；
\item $\Delta$ 的每个顶点都为三价。若 $b\in B_0$ 位于 $\Delta$ 的顶点 $b_0$ 附近，$\gamma_1,\gamma_2,\gamma_3$ 为以 $b$ 为基点、分别绕从 $b_0$ 出发的三条边的简单环路，选择其定向使得 $\gamma_1\gamma_2\gamma_3=1$ 于 $\pi_1(B_0,b)$ 中，则 $T_{\gamma_i}\colon\Lambda_b\to\Lambda_b$，$i=1,2,3$，在适当的基中分别具有形式
\[
\begin{pmatrix}1&0&1\\0&1&0\\0&0&1\end{pmatrix},\qquad
\begin{pmatrix}1&0&-1\\0&1&-1\\0&0&1\end{pmatrix},\qquad
\begin{pmatrix}1&0&0\\0&1&1\\0&0&1\end{pmatrix},
\]
或者
\[
\begin{pmatrix}1&0&1\\0&1&0\\0&0&1\end{pmatrix},\qquad
\begin{pmatrix}1&1&0\\0&1&0\\0&0&1\end{pmatrix},\qquad
\begin{pmatrix}1&-1&-1\\0&1&0\\0&0&1\end{pmatrix}.
\]
\end{enumerate}
\end{definition}

\SourceParagraph{ch06-6.4.3-009}{503}
目标是：对于简单、紧、带奇点的整仿射流形 $B$，紧化 $X(B_0)$ 或 $\check X(B_0)$。

\SourceParagraph{ch06-6.4.3-010}{504}
先从一个特殊情形开始。假定 $\Delta$ 的连通分支 $C$ 没有三价顶点。那么，$C$ 为圆周，可以说明，如何把 $f\colon X(B_0)\to B_0$ 延拓过 $C$。令 $B'$ 为 $C$ 在 $B$ 中的一个管状邻域，满足 $B'\cap\Delta=C$，其边界 $\partial B'$ 为二维环面，如\BookRef{fig:6.11}{图11}所示。可以将 $T^3$ 丛 $f\colon X(B_0)\to B_0$ 限制到 $\partial B'$，并问，这一环面丛绕 $\pi_1(\partial B',b)$ 的生成元 $\gamma_1,\gamma_2$ 的单值变换，能够是什么；其中，$b\in\partial B'$ 为基点。
\begin{figure}[H]
\centering
\BookImage{01_Numbered_Figures/Chapter_06/Ch06_Fig11_PDFp504_Bookp494.pdf}
\BookNumberOnlyCaption\label{fig:6.11}
\end{figure}
按假设，$H_1(f^{-1}(b),\mathbb Z)$ 有一组基 $e_1,e_2,e_3$，使得 $T_{\gamma_1}=\begin{pmatrix}1&1&0\\0&1&0\\0&0&1\end{pmatrix}$。由于 $\pi_1(\partial B',b)$ 为 Abel 群，$T_{\gamma_1}$ 与 $T_{\gamma_2}$ 必须对易。若写成 $T_{\gamma_2}=(a_{ij})_{1\leq i,j\leq3}$，这一对易性保证
\[
a_{21}=a_{23}=a_{31}=0,\qquad a_{11}=a_{22},
\]
所以，
\[
T_{\gamma_2}=\begin{pmatrix}a_{11}&a_{12}&a_{13}\\0&a_{11}&0\\0&a_{32}&a_{33}\end{pmatrix}.
\]
由于 $\det T_{\gamma_2}=\pm1$，必有 $a_{11}=\pm1$、$a_{33}=\pm1$。因此，与 $T_{\gamma_1}$ 对易的矩阵所构成的群 $G\subseteq\operatorname{GL}_3(\mathbb Z)$，由下列矩阵生成：
\[
\begin{pmatrix}1&a&b\\0&1&0\\0&0&1\end{pmatrix},\ a,b\in\mathbb Z,\qquad
\begin{pmatrix}1&0&0\\0&1&0\\0&1&1\end{pmatrix},\qquad
\begin{pmatrix}1&0&0\\0&1&0\\0&0&-1\end{pmatrix},\qquad
\begin{pmatrix}-1&0&0\\0&-1&0\\0&0&1\end{pmatrix}.
\]

\SourceParagraph{ch06-6.4.3-011}{504}
现在，与二维情形中一样，令 $\overline h\colon\overline X\to\overline D$，为\BookRef{constr:6.86}{构造6.86}在闭圆盘 $\overline D$ 上的纤维化。那么，$\dim\overline X=4$，$\overline h$ 有一个奇异纤维，为夹缩环面。将在 $\overline B:=S^1\times\overline D$ 上构造环面纤维化：在 $[0,1]\times\overline X\times S^1$ 中，通过同胚 $\varphi\colon\overline X\times S^1\to\overline X\times S^1$，将 $\{0\}\times\overline X\times S^1$ 与 $\{1\}\times\overline X\times S^1$ 识别，要求 $\varphi$ 放入交换图
\[
\BookImage{02_Unnumbered_Diagrams/Chapter_06/Ch06_Diagram_PDFp505_01_Bookp495.pdf}
\]
其中，竖箭头为先投影到 $\overline X$，再与 $\overline h$ 复合。作这一识别，得到六维流形 $Z$，带有纤维化 $\overline f\colon Z\to\overline B$，奇异纤维位于 $S^1\times\{0\}$ 上。

\SourceParagraph{ch06-6.4.3-012}{505}
用这种方式得到的环面纤维化，其单值变换可如下描述。回顾一下，按照 $\overline X$ 的构造，有映射 $\overline X\to\overline Y=S^1\times\overline D$，纤维上有 $S^1$ 作用，且这个作用有一个不动点。固定 $b\in\overline D\setminus\{0\}$，选择 $H_1(\overline h^{-1}(b),\mathbb Z)$ 的生成元 $e_1,e_2$，使得 $e_1$ 是这个 $S^1$ 作用的轨道，绕0的单值变换具有形式 $\begin{pmatrix}1&1\\0&1\end{pmatrix}$。再将 $b$ 与纤维化 $\overline f$ 的底空间 $\overline B$ 中的点识别，即 $b\in\{0\}\times(\overline D\setminus\{0\})\subseteq S^1\times(\overline D\setminus\{0\})$；这里以 $\mathbb R/\mathbb Z$ 表示圆周。那么，$H_1(\overline f^{-1}(b),\mathbb Z)=H_1(\overline h^{-1}(b)\times S^1,\mathbb Z)$ 的一组基为
\[
e_1=e_1\times\{\mathrm{pt}\},\qquad
e_2=e_2\times\{\mathrm{pt}\},\qquad
e_3=\{\mathrm{pt}\}\times S^1.
\]
于是，绕环路 $\gamma_1$ 的单值变换，为所需的 $T_{\gamma_1}=\begin{pmatrix}1&1&0\\0&1&0\\0&0&1\end{pmatrix}$。由于 $\varphi\colon\{0\}\times\overline X\times S^1\to\{1\}\times\overline X\times S^1$ 被选择为保持纤维，$\varphi$ 作用于 $\overline f^{-1}(b)$，并且 $T_{\gamma_2}=\varphi_*\colon H_1(\overline f^{-1}(b),\mathbb Z)\to H_1(\overline f^{-1}(b),\mathbb Z)$。因此，为实现特定的 $T_{\gamma_2}$，需要选择一个在纤维上实现这一变换的微分同胚 $\varphi$。对于群 $G$ 的每个生成元，分别说明如何做到。

\SourceParagraph{ch06-6.4.3-013}{505}
对于 $\begin{pmatrix}1&a&b\\0&1&0\\0&0&1\end{pmatrix}$，有 $S^1$ 纤维化 $\overline q\colon\overline X\to\overline Y=S^1\times\overline D$，因而有诱导映射
\[
\overline X\times S^1\xrightarrow{\overline q\times\operatorname{id}}
(S^1\times\overline D)\times S^1.
\]
将 $S^1$ 与 $\mathbb C$ 中的单位圆识别，则有映射 $m\colon S^1\times S^1\to S^1$，由 $(z_1,z_2)\mapsto z_1^a z_2^b$ 给出。令 $\pi\colon S^1\times\overline D\times S^1\to S^1\times S^1$ 为投影。回顾一下，$S^1$ 作用于 $\overline X$，因而也作用于 $\overline X\times S^1$，在第二个因子上的作用平凡。于是，可以定义
\[
\varphi(x)=\bigl(m\circ\pi\circ(\overline q\times\operatorname{id})(x)\bigr)\cdot x.
\]
容易看到，$\varphi$ 保持 $\overline X\times S^1\to\overline D$ 的纤维，在每个纤维上的作用，在同调上诱导矩阵 $\begin{pmatrix}1&a&b\\0&1&0\\0&0&1\end{pmatrix}$。

\SourceParagraph{ch06-6.4.3-014}{505}
对于 $\begin{pmatrix}1&0&0\\0&1&0\\0&1&1\end{pmatrix}$，在 $(S^1\times\overline D)\times S^1$ 上通过 $(z_1,b,z_2)\mapsto(z_1,b,z_1z_2)$ 作用。注意，$\overline X\times S^1\to(S^1\times\overline D)\times S^1$，限制到 $((S^1\times\overline D)\setminus\{(p,0)\})\times S^1$ 上，是 $S^1$ 丛。容易证明，这一作用保持该 $S^1$ 丛的 Chern 类。因此，这个同胚提升为 $\overline X\times S^1$ 的同胚。

\SourceParagraph{ch06-6.4.3-015}{506}
对于 $\begin{pmatrix}1&0&0\\0&1&0\\0&0&-1\end{pmatrix}$，只需在 $\overline X\times S^1$ 上通过 $(x,z)\mapsto(x,\bar z)$ 作用。

\SourceParagraph{ch06-6.4.3-016}{506}
对于 $\begin{pmatrix}-1&0&0\\0&-1&0\\0&0&1\end{pmatrix}$，把 $\overline X\to\overline D$ 的光滑纤维看作 Abel 群，通过取负作用，并将这一作用延伸到夹缩纤维。或者，利用\BookRef{constr:6.87}{构造6.87}中对 $\overline X\to\overline D$ 的描述，通过 $(z_1,z_2,w)\mapsto(\bar z_1,\bar z_2,\bar w)$ 作用于 $\overline X$。

\SourceParagraph{ch06-6.4.3-017}{506}
因此，给定 $T_{\gamma_1}$ 的形式，就能构造具有所需任意相容单值变换的 $T^3$ 纤维化 $\overline f\colon Z\to\overline B$。然后，可以将其粘入 $X(B_0)$。

\SourceParagraph{ch06-6.4.3-018}{506}
\BookStatementNumber{96}
\begin{exercise}[粘合的细节]\label{exer:6.96}
将 $C\cong S^1$ 的管状邻域的闭包 $\overline B'$ 与 $\overline B$ 识别。将 $T^3$ 丛 $\overline f\colon\overline f^{-1}(\partial\overline B)\to\partial\overline B$，描述为 $\mathcal E/\Gamma\to\partial\overline B$，其中，$\mathcal E$ 为 $\partial\overline B$ 上的实秩3向量丛，$\Gamma$ 为其中的一族格；单值变换保持定向时，这里可取 $\mathcal E$ 为平凡丛。像\BookRef{exer:6.94}{练习6.94}中一样，证明，利用这一描述，可以沿共同边界粘合纤维化 $\overline f$ 与 $f\colon X(B_0\setminus\overline B')\to B_0\setminus\overline B'$。
\end{exercise}

\SourceParagraph{ch06-6.4.3-019}{506}
\BookStatementNumber{97}
\begin{note}[numbered]\label{note:6.97}
对于截面的存在性，粘合时还需要格外小心。当然，$X(B_0)\to B_0$ 有截面，由 $TB_0$ 的零截面给出。希望保证，所构造的紧化 $X(B)\to B$ 也有截面。为此，可以为每个局部模型找到截面，并确保粘合时，$X(B_0)\to B_0$ 的零截面，与所选局部模型的截面能够识别。

\SourceParagraph{ch06-6.4.3-020}{506}
例如，先考虑上述 $T^2$ 纤维化 $\overline X\to\overline D$。有 $S^1$ 纤维化 $\overline q\colon\overline X\to\overline Y=S^1\times\overline D$，在 $(p,0)\in S^1\times\overline D$ 上退化。选择点 $p'\in S^1$、$p'\ne p$。那么，$\overline q^{-1}(\{p'\}\times\overline D)\to\{p'\}\times\overline D$，是圆盘上的 $S^1$ 丛，因而必然平凡。可以选择这个 $S^1$ 丛的截面，通过复合
\[
\overline D\xrightarrow{\ \cong\ }\{p'\}\times\overline D
\longrightarrow\overline q^{-1}(\{p'\}\times\overline D)\subseteq\overline X,
\]
定义截面 $\sigma\colon\overline D\to\overline X$；其中，第二个箭头为所选 $S^1$ 丛的截面。将如下检查留给读者：在\BookRef{exer:6.94}{练习6.94}中，可以将 $T^2$ 丛 $\partial\overline X\to\partial\overline D$ 与 $\mathcal E/\Gamma\to\partial\overline D$ 识别，使得 $\sigma|_{\partial\overline D}$ 与零截面识别。这样，就能按\BookRef{exer:6.94}{练习6.94}进行粘合，使截面彼此匹配，从而得到具有截面的紧化。

\SourceParagraph{ch06-6.4.3-021}{506}
在刚才考虑的三维情形中，只要上述构造的 $\overline f\colon Z\to S^1\times\overline D$ 始终有截面，就可以采用同一过程。将如下检查留给读者：对于上述每个 $T_{\gamma_2}$ 的选择，确实如此，并且\BookRef{exer:6.96}{练习6.96}中的粘合，可以安排为使截面粘合在一起。
\end{note}

\SourceParagraph{ch06-6.4.3-022}{507}
现在，已经可以完成三维紧化。首先，若 $\Delta$ 有不含顶点的连通分支，则如前所述，把纤维丛 $X(B_0)\to B_0$ 延拓过这些圆周。

\SourceParagraph{ch06-6.4.3-023}{507}
分两步完成剩下的紧化。令 $\overline B'$ 为 $\Delta$ 的其余分支的管状邻域。令 $l$ 为包含于 $\Delta$、连接两个顶点的线段，令 $l'\subseteq l$ 为位于 $l$ 内部的较小闭区间。选择 $l'$ 的闭管状邻域 $\overline U\cong\overline D\times l'$，如\BookRef{fig:6.12}{图12}所示，使得 $\partial\overline D\times l'\subseteq\partial\overline B'$。然后，与前面一样，可以粘入纤维化 $\overline X\times S^1\times l'\to\overline U$。对每条这样的支边，都进行这一操作。
\begin{figure}[H]
\centering
\BookImage{01_Numbered_Figures/Chapter_06/Ch06_Fig12_PDFp507_Bookp497.pdf}
\BookNumberOnlyCaption\label{fig:6.12}
\end{figure}

\SourceParagraph{ch06-6.4.3-024}{507}
注意，与\BookRef{note:6.97}{笔记6.97}中一样，可以安排粘合，使 $X(B\setminus\overline B')\to B\setminus\overline B'$ 的零截面，与 $\overline X\times S^1\times l'\to\overline U$ 的所选截面匹配。

\SourceParagraph{ch06-6.4.3-025}{507}
还需要保持另一种相容性。$\overline X\times S^1\times l'$ 与 $X(\overline U\cap(B\setminus\overline B'))$，都具有自然的 $T^2=S^1\times S^1$ 作用。对于 $\overline X\times S^1\times l'$，第一个 $S^1$ 因子按照构造作用于 $\overline X$，第二个因子作用于 $S^1$ 因子。另一方面，若绕支边的单值变换，在基 $e_1,e_2,e_3$ 中由矩阵 $\begin{pmatrix}1&1&0\\0&1&0\\0&0&1\end{pmatrix}$ 给出，则在 $\overline U\cap(B\setminus\overline B')$ 上，$e_1,e_3$ 局部地给出 $\Lambda$ 的良定义截面，$T^2=(\mathbb Re_1/\mathbb Ze_1)\times(\mathbb Re_3/\mathbb Ze_3)$ 自然地通过平移作用于 $X(\overline U\cap(B\setminus\overline B'))$。要求安排粘合，使这两个作用匹配。不难看到，这可以做到。

\SourceParagraph{ch06-6.4.3-026}{507}
进行这一粘合以后，得到纤维化 $f\colon X(B_0')\to B_0'$，其中，$B_0'\subseteq B$ 为 $B\setminus\overline B'$ 与上述各个集合 $\overline U$ 的并。现在，只剩下若干三价顶点的邻域，如\BookRef{fig:6.13}{图13}所示。令 $U$ 为图中那样的三价顶点的开邻域，令 $\overline V=U\cap(B\setminus B_0')$，于是，有纤维化 $f\colon X(U\setminus\overline V)\to U\setminus\overline V$。
\begin{figure}[H]
\centering
\BookImage{01_Numbered_Figures/Chapter_06/Ch06_Fig13_PDFp508_Bookp498.pdf}
\BookNumberOnlyCaption\label{fig:6.13}
\end{figure}
这个三价顶点，在\BookSectionRef{sec:6.4.2}{6.4.2}的意义下，要么为正型顶点，要么为负型顶点。

\SourceParagraph{ch06-6.4.3-027}{508}
先考虑负型顶点的情形。$\Lambda$ 在顶点附近的单值变换，有一个不变量：在\BookRef{def:6.95}{定义6.95}(2)中所选的基 $e_1,e_2,e_3$ 中，$e_1$ 不变。于是，$S^1=\mathbb Re_1/\mathbb Ze_1$ 通过平移作用于 $X(U\setminus\Delta)$。由于已经要求粘合保持这一作用，实际上在 $f^{-1}(\partial\overline V)$ 上得到 $S^1$ 作用。

\SourceParagraph{ch06-6.4.3-028}{508}
\BookStatementNumber{98}
\begin{exercise}\label{exer:6.98}
证明 $f^{-1}(\partial\overline V)/S^1\cong T^2\times\partial\overline V$。利用这一点，证明 $f^{-1}(\partial\overline V)\to\partial\overline V$，与\BookRef{constr:6.90}{构造6.90}中的映射 $f^{-1}(\partial\overline B)\to\partial\overline B$ 相同，因此，可以将这一纤维化粘入。注意，可以安排这一操作，使所选截面与 $f\colon X(B_0')\to B_0'$ 的零截面粘合，并使 $S^1$ 作用匹配，从而将 $S^1$ 作用延拓到紧化。
\end{exercise}

\SourceParagraph{ch06-6.4.3-029}{508}
同样，对于正型顶点，$\Lambda$ 有一个秩2的单值变换不变子群，因而有 $T^2$ 作用。此时，$f^{-1}(\partial\overline V)/T^2\cong S^1\times\partial\overline V$。与前面一样，可以证明，能够将\BookRef{constr:6.92}{构造6.92}的正型顶点纤维化粘入。仍然可以安排这一操作，使其保持截面与局部 $T^2$ 作用。

\SourceParagraph{ch06-6.4.3-030}{508}
按上述步骤完成紧化后，得到环面纤维化 $f\colon X(B)\to B$，具有以下性质。
\begin{property*}\label{pty:6.compactification}
\begin{enumerate}[label=(\arabic*)]
\item $f|_{f^{-1}(B\setminus\overline B')}$，与 $X(B\setminus\overline B')\to B\setminus\overline B'$ 相同；
\item $f$ 有截面，延伸 $X(B\setminus\overline B')\to B\setminus\overline B'$ 的零截面；
\item 对于任意开集 $U\subseteq B$，$f^{-1}(U)\to U$ 上存在沿纤维的环面平移作用，由群 $\Gamma(U\setminus\Delta,\Lambda_{\mathbb R}/\Lambda)$ 给出。
\end{enumerate}
\end{property*}
这些性质，由局部地保持了截面与作用这一事实得到。

\SourceParagraph{ch06-6.4.3-031}{509}
\BookStatementNumber{99}
\begin{note}[numbered]\label{note:6.99}
这一作用，使得可以对 $f\colon X(B)\to B$ 进行“拓扑扭曲”。事实上，给定元素
\[
\mathbf B\in H^1(B,i_*(\Lambda_{\mathbb R}/\Lambda))\text{，}
\]
由 Čech 1上闭链 $(b_{ij})$ 代表，其中，$b_{ij}\in\Gamma((U_i\cap U_j)\setminus\Delta,\Lambda_{\mathbb R}/\Lambda)$，就可以像\BookSectionRef{sec:6.2.3}{6.2.3}中一样，重新粘合纤维化 $f\colon X(B)\to B$，得到 $X(B,\mathbf B)\to B$。不过，现在这一操作发生在纯拓扑层面。若新的纤维化 $X(B,\mathbf B)\to B$ 没有截面，可能得到一个不同的拓扑流形。
\end{note}

\SourceParagraph{ch06-6.4.3-032}{509}
\BookStatementNumber{100}
\begin{example}\label{ex:6.100}
考虑\BookRef{ex:6.74}{例题6.74}中的 $B$。在每个例子中，现在可以构造 $X(B_0)\to B_0$ 的紧化 $X(B)\to B$，以及 $\check X(B_0)\to B_0$ 的紧化 $\check X(B)\to B$。如\BookSectionRef{sec:6.4.1}{6.4.1}末尾所述，预期 $X(B)$ 在拓扑上同胚于 $\mathbb P_{\Xi^*}$ 的超平面截面的一个消解，而 $\check X(B)$ 同胚于 $\mathbb P_\Xi$ 的超平面截面的一个消解。在那里给出的例子中，的确如此。

\SourceParagraph{ch06-6.4.3-033}{509}
具体地，可以证明，在二维例子中，得到拓扑 K3 曲面。在三维例子中，Gross\BookCite{ref202}的主要结果是，$X(B)$ 同胚于镜像五次三维流形，而 $\check X(B)$ 同胚于五次三维流形。

\SourceParagraph{ch06-6.4.3-034}{509}
对于任意维数中更一般的反身多胞体选择，也可以证明一些类似结果。Gross\BookCite{ref203}与 Gross、Siebert\BookCite{ref204}对此作了研究。这里省略这些陈述的细节。
\end{example}

\subsection{SYZ 纤维化的上同调}\label{sec:6.4.4}

\SourceParagraph{ch06-6.4.4-001}{509}
令 $B$ 为带奇点的仿射流形，$i\colon B_0\hookrightarrow B$ 为包含映射。设有 $X(B_0)$ 与 $\check X(B_0)$ 的拓扑紧化，分别记为 $X(B)$ 与 $\check X(B)$，连同环面纤维化 $f\colon X(B)\to B$、$\check f\colon\check X(B)\to B$，分别延伸 $f_0\colon X(B_0)\to B_0$ 与 $\check f_0\colon\check X(B_0)\to B_0$。假定处于二维或三维情形，则只要对单值变换施加适当限制，即要求 $B$ 简单，这些紧化就存在。此外，若 $B$ 简单，则在三维情形中，$X(B)$ 与 $\check X(B)$ 的 Euler 示性数符号相反。事实上，$X(B)\to B$ 与 $\check X(B)\to B$ 都只有有限多个纤维具有非零 Euler 示性数，因此，$X(B)$ 或 $\check X(B)$ 的 Euler 示性数，只是这些纤维的 Euler 示性数之和。于是，
\[
\begin{aligned}
\chi(X(B))
&=\#\{X(B)\to B\text{ 的正型纤维}\}\\
&\qquad-\#\{X(B)\to B\text{ 的负型纤维}\}.
\end{aligned}
\]
不过，SYZ 对偶化过程交换正型与负型纤维，因此，显然有 $\chi(X(B))=-\chi(\check X(B))$。

\SourceParagraph{ch06-6.4.4-002}{509}
然而，镜像对称至少要求镜像对的 Hodge 数相互交换。为什么应当预期这一点会发生？这里将在三维情形中解释，不过，先给出一般设定。

\SourceParagraph{ch06-6.4.4-003}{510}
给定 $f\colon X(B)\to B$，可以考虑层 $R^qf_*\mathbb Q$。它是 $B$ 上的预层 $U\mapsto H^q(f^{-1}(U),\mathbb Q)$ 的关联层。同样，$B_0$ 上有层 $R^qf_{0*}\mathbb Q$。由于 $f_0$ 为环面丛，后一个层为局部系统，即局部地同构于以环面的上同调为系数的常值层。当然，$(R^qf_*\mathbb Q)|_{B_0}=R^qf_{0*}\mathbb Q$，不过，$R^qf_*\mathbb Q$ 沿 $\Delta=B\setminus B_0$ 一般不再是局部系统。在这里的固有纤维化中，$R^pf_*\mathbb Q$ 的茎由下式给出：
\[
(R^pf_*\mathbb Q)_b\cong H^p(f^{-1}(b),\mathbb Q).
\]
因此，若 $f^{-1}(b)$ 为奇异纤维，可以预期，这个上同调群与非奇异环面的上同调群不同。

\SourceParagraph{ch06-6.4.4-004}{510}
\BookStatementNumber{101}
\begin{definition}\label{def:6.101}
若对所有 $p$，自然限制比较映射给出同构
\[
R^pf_*\mathbb Q\cong i_*R^pf_{0*}\mathbb Q,
\]
则称 $f\colon X\to B$ 为 $\mathbb Q$-简单的。更一般地，若 $G$ 为任意 Abel 群，则当
\[
R^pf_*G\cong i_*R^pf_{0*}G
\]
对所有 $p$ 都通过自然比较映射成立时，称 $f$ 为 $G$-简单的。
\end{definition}

\SourceParagraph{ch06-6.4.4-005}{510}
注意，按定义，若 $U\subseteq B$ 为开集，则
\[
\Gamma(U,i_*R^pf_{0*}\mathbb Q)
=\Gamma(U\cap B_0,R^pf_{0*}\mathbb Q).
\]
因此，简单性本质上要求：若 $b\in B\setminus B_0$，则对于 $b$ 的充分小开邻域 $U$，有
\[
H^p(f^{-1}(b),\mathbb Q)
\cong\Gamma(U\cap B_0,R^pf_{0*}\mathbb Q).
\]

\SourceParagraph{ch06-6.4.4-006}{510}
这是相当强的假设。来看看它有什么帮助。首先，考虑一个环面 $V/\Gamma$ 的上同调，其中，$V$ 为 $n$ 维向量空间，$\Gamma$ 为秩 $n$ 的格。那么，
\[
H_1(V/\Gamma,\mathbb Z)\cong\Gamma,\qquad
H^1(V/\Gamma,\mathbb Z)\cong\Gamma^\vee,\qquad
H^p(V/\Gamma,\mathbb Z)\cong\bigwedge^p\Gamma^\vee.
\]
另一方面，把 $V^\vee/\Gamma^\vee$ 看作对偶环面，将 $\Gamma^\vee$ 与 $\{v\in V^\vee\mid v(\Gamma)\subseteq\mathbb Z\}$ 识别，就看到 $H^p(V^\vee/\Gamma^\vee,\mathbb Z)\cong\bigwedge^p\Gamma$。若选择 $\bigwedge^n\Gamma^\vee\cong\mathbb Z$ 的一个生成元，则自然配对
\[
\bigwedge^p\Gamma^\vee\otimes\bigwedge^{n-p}\Gamma^\vee
\longrightarrow\bigwedge^n\Gamma^\vee\xrightarrow{\ \cong\ }\mathbb Z
\]
诱导同构 $\bigwedge^p\Gamma^\vee\cong\bigwedge^{n-p}\Gamma$。因此，在选定 $\bigwedge^n\Gamma^\vee$ 的生成元后，有典范同构
\[
H^p(V/\Gamma,\mathbb Z)\cong H^{n-p}(V^\vee/\Gamma^\vee,\mathbb Z).
\]

\SourceParagraph{ch06-6.4.4-007}{511}
将其用于 $f_0\colon X(B_0)\to B_0$ 与 $\check f_0\colon\check X(B_0)\to B_0$，并记住 $X(B_0)=TB_0/\Lambda$、$\check X(B_0)=T^*B_0/\check\Lambda$，有
\[
R^1f_{0*}\mathbb Z=\check\Lambda,\qquad R^1\check f_{0*}\mathbb Z=\Lambda.
\]
因此，
\[
R^pf_{0*}\mathbb Z=\bigwedge^p\check\Lambda,\qquad
R^p\check f_{0*}\mathbb Z=\bigwedge^p\Lambda.
\]
与 $\mathbb Q$ 作张量积，得到
\[
R^pf_{0*}\mathbb Q=\bigwedge^p\check\Lambda_{\mathbb Q},\qquad
R^p\check f_{0*}\mathbb Q=\bigwedge^p\Lambda_{\mathbb Q},
\]
其中，$\Lambda_{\mathbb Q}=\Lambda\otimes_{\mathbb Z}\mathbb Q$、$\check\Lambda_{\mathbb Q}=\check\Lambda\otimes_{\mathbb Z}\mathbb Q$。现在，假定 $\bigwedge^n\Lambda\cong\mathbb Z$；若 $B$ 的转移映射属于 $\mathbb R^n\rtimes\operatorname{SL}_n(\mathbb Z)$，如\BookRef{exer:6.13}{练习6.13}中一样，就满足这一条件。那么，选定生成元后，有像该小节中一样的自然配对
\[
\bigwedge^p\Lambda\times\bigwedge^{n-p}\Lambda
\longrightarrow\bigwedge^n\Lambda\cong\mathbb Z,
\]
给出同构 $\bigwedge^p\check\Lambda\cong\bigwedge^{n-p}\Lambda$。因此，
\[
R^pf_{0*}\mathbb Q\cong R^{n-p}\check f_{0*}\mathbb Q,
\qquad
i_*R^pf_{0*}\mathbb Q\cong i_*R^{n-p}\check f_{0*}\mathbb Q.
\]
若两个纤维化都为 $\mathbb Q$-简单的，就得到
\begin{equation}\label{eq:6.7}
R^pf_*\mathbb Q\cong R^{n-p}\check f_*\mathbb Q.
\end{equation}
利用 Leray 谱序列，确定 $X(B)$ 与 $\check X(B)$ 的上同调之间的关系时，这将极其有用。

\SourceParagraph{ch06-6.4.4-008}{511}
已经两次使用数学中被过度使用的词“simple”，即“简单”。来看看，\BookRef{def:6.95}{定义6.95}中的简单性，蕴含\BookRef{def:6.101}{定义6.101}中的 $\mathbb Q$-简单性。

\SourceParagraph{ch06-6.4.4-009}{511}
\BookStatementNumber{102}
\begin{proposition}\label{prop:6.102}
若 $B$ 为三维、简单的，则 $f_0$ 与 $\check f_0$ 的紧化 $f\colon X(B)\to B$ 与 $\check f\colon\check X(B)\to B$，为 $\mathbb Q$-简单的。
\end{proposition}
\begin{proof}
\SourceParagraph{ch06-6.4.4-010}{511}
令 $b\in\Delta$。只讨论 $f$；$\check f$ 的情形相同。只需证明：对于 $b$ 的充分小邻域 $U$，满足 $\Gamma(U,R^pf_*\mathbb Q)\cong H^p(f^{-1}(b),\mathbb Q)$，限制映射
\[
\Gamma(U,R^pf_*\mathbb Q)\longrightarrow
\Gamma(U\setminus\Delta,R^pf_*\mathbb Q)
=\Gamma(U,i_*R^pf_{0*}\mathbb Q)
\]
为同构。现在，$\Gamma(U,i_*R^pf_{0*}\mathbb Q)$ 可以与 $H^p(f^{-1}(b_0),\mathbb Q)$ 中的单值变换不变元素的空间识别，其中，$b_0\in U\setminus\Delta$。特别地，只需证明，上同调类的限制
\[
H^p(f^{-1}(U),\mathbb Q)\longrightarrow H^p(f^{-1}(b_0),\mathbb Q)
\]
给出到上同调群 $H^p(f^{-1}(b_0),\mathbb Q)$ 的单值变换不变子空间 $W$ 的同构；对偶地，映射
\[
H_p(f^{-1}(b_0),\mathbb Q)\longrightarrow H_p(f^{-1}(U),\mathbb Q)
\]
在 $H_p(f^{-1}(b_0),\mathbb Q)/W^\perp$ 与 $H_p(f^{-1}(U),\mathbb Q)$ 之间诱导同构。通过将 $f^{-1}(U)$ 形变收缩到 $f^{-1}(b)$，后一个群可以与 $H_p(f^{-1}(b),\mathbb Q)$ 识别。因此，对于 $f^{-1}(b)$ 的三种不同可能，可以分别检查这个陈述：\BookRef{constr:6.88}{构造6.88}中的一般奇异纤维、\BookRef{constr:6.90}{构造6.90}中的负型纤维，以及\BookRef{constr:6.92}{构造6.92}中的正型纤维。

\SourceParagraph{ch06-6.4.4-011}{512}
考虑负型情形，其他情形留给读者。首先，$p=0,3$ 时，结论平凡。现在，$f^{-1}(b)$ 是 $T^2$ 上的 $S^1$ 纤维化，其中，位于“8”字形上的 $S^1$ 退化为点。容易看到，$H_1(f^{-1}(b),\mathbb Q)=\mathbb Q^2$，由映到 $H_1(T^2,\mathbb Q)$ 的两个生成元的圆圈生成。另一方面，在\BookRef{exer:6.91}{练习6.91}中，已经计算了这一纤维化的单值变换：绕 $\Delta$ 的支边的环路，在 $H_1(f^{-1}(b_0),\mathbb Z)$ 上的单值变换，由矩阵
\[
\begin{pmatrix}1&0&1\\0&1&0\\0&0&1\end{pmatrix},\qquad
\begin{pmatrix}1&1&0\\0&1&0\\0&0&1\end{pmatrix},\qquad
\begin{pmatrix}1&-1&-1\\0&1&0\\0&0&1\end{pmatrix}
\]
给出。这里，基为 $e_1,e_2,e_3$，其中，$e_1$ 为构造 $f^{-1}(U)$ 时使用的、$f^{-1}(U)$ 上 $S^1$ 作用的一条轨道的类；$e_2,e_3$ 是将上述 $H_1(f^{-1}(b),\mathbb Z)$ 的两个生成元，形变到 $f^{-1}(b_0)$ 的选择。这类形变的选择并不唯一，其不唯一性由单值变换编码。不过，从这一单值变换的描述可知，$W^\perp$ 由 $e_1$ 生成，因此，得到所需同构。对 $H_2$，类似的论证也成立：$H_2(f^{-1}(b),\mathbb Q)$ 由一个 $T^2$ 生成，它是投影 $f^{-1}(b)\to T^2$ 的截面。其余细节留给读者。
\end{proof}

\SourceParagraph{ch06-6.4.4-012}{512}
现在，详细解释 Leray 谱序列是什么，以及它能做什么。对于任意连续映射 $f\colon X\to B$，Leray 谱序列为
\[
E_2^{p,q}=H^p(B,R^qf_*\mathbb Q)\quad\Longrightarrow\quad H^{p+q}(X,\mathbb Q).
\]
这一记号包含大量信息。首先，可以把群 $H^p(B,R^qf_*\mathbb Q)$ 排成一个阵列，组成谱序列的所谓 $E_2$ 页，或 $E_2$ 项。
\[
\BookImage{02_Unnumbered_Diagrams/Chapter_06/Ch06_Diagram_PDFp513_01_Bookp503.pdf}
\]

\SourceParagraph{ch06-6.4.4-013}{513}
还有映射的附加数据：
\[
d_2\colon H^p(B,R^qf_*\mathbb Q)
\longrightarrow H^{p+2}(B,R^{q-1}f_*\mathbb Q).
\]
若 $p$ 或 $q-1$ 为负数，则这一映射为零。不过，还不止如此：两个这样的映射的复合始终为零。若令
\[
E_3^{p,q}=
\frac{\ker\bigl(d_2\colon H^p(B,R^qf_*\mathbb Q)
\to H^{p+2}(B,R^{q-1}f_*\mathbb Q)\bigr)}
{\operatorname{im}\bigl(d_2\colon H^{p-2}(B,R^{q+1}f_*\mathbb Q)
\to H^p(B,R^qf_*\mathbb Q)\bigr)},
\]
则还得到映射 $d_3\colon E_3^{p,q}\to E_3^{p+3,q-2}$。这里不解释这些映射从何而来；这是谱序列的神奇之处的一部分。若希望理解谱序列机制的来源，可以参阅任意同调代数教材，例如文献\BookCite{ref053}。可以这样继续，在第 $r$ 页，有群 $E_r^{p,q}$、映射
\[
d_r\colon E_r^{p,q}\longrightarrow E_r^{p+r,q-r+1},
\]
以及
\[
E_{r+1}^{p,q}=
\frac{\ker(d_r\colon E_r^{p,q}\to E_r^{p+r,q-r+1})}
{\operatorname{im}(d_r\colon E_r^{p-r,q+r-1}\to E_r^{p,q})}.
\]

\SourceParagraph{ch06-6.4.4-014}{513}
由于 $p<0$ 或 $q<0$ 时，$E_r^{p,q}=0$，所以对于给定的 $p,q$，存在充分大的 $r$，使得 $E_r^{p,q}=E_{r+1}^{p,q}=\cdots$。若将这个稳定的群记为 $E_\infty^{p,q}$，谱序列的神奇之处便给出以下结论：对于每个总次数 $n$，存在滤过
\[
H^n(X,\mathbb Q)=F^0\supseteq\cdots\supseteq F^n\supseteq F^{n+1}=0,
\]
满足 $F^p/F^{p+1}\cong E_\infty^{p,n-p}$。换言之，应当把 Leray 谱序列看作一台机器：输入是 $B$ 上某些层的上同调群，转动曲柄以后，就产生 $X$ 的上同调的一个滤过的分次片。

\SourceParagraph{ch06-6.4.4-015}{514}
一般的困难是，并不知道这些映射 $d_r$ 是什么，甚至最初的那些群，也可能很难计算。

\SourceParagraph{ch06-6.4.4-016}{514}
来看看，上述机制对映射 $f\colon X(B)\to B$ 与 $\check f\colon\check X(B)\to B$ 给出什么；这里，$\dim B=3$。首先，由于纤维连通，$R^0f_{0*}\mathbb Q=f_{0*}\mathbb Q=\mathbb Q$。还假定 $R^3f_{0*}\mathbb Q\cong\mathbb Q$；这就是 $\bigwedge^3\check\Lambda\cong\mathbb Z$ 的假设，等价于 $B_0$ 的完整表示的像包含于 $\mathbb R^3\rtimes\operatorname{SL}_3(\mathbb Z)$。这样，可以选择 $R^3f_{0*}\mathbb Q$ 的一个整体截面。由 $\mathbb Q$-简单性，对于 $p=0,3$，有
\[
R^pf_*\mathbb Q=i_*R^pf_{0*}\mathbb Q=i_*\mathbb Q=\mathbb Q.
\]
注意，$i_*\mathbb Q=\mathbb Q$ 来自 $B\setminus B_0$ 的余维为二这一事实。其次，进一步假定 $B$ 单连通，在这里的紧三维情形中为 $S^3$，则当 $q=0,3$ 时，知道 $H^p(B,R^qf_*\mathbb Q)$。于是，$f$ 与 $\check f$ 的 Leray 谱序列的 $E_2$ 页分别为
\[
\BookImage{02_Unnumbered_Diagrams/Chapter_06/Ch06_Diagram_PDFp514_01_Bookp504.pdf}
\]
与
\[
\BookImage{02_Unnumbered_Diagrams/Chapter_06/Ch06_Diagram_PDFp514_02_Bookp504.pdf}
\]

\SourceParagraph{ch06-6.4.4-017}{514}
这两个阵列具有对称性，即由\BookEqRef{eq:6.7}{6.7}产生自然同构
\[
H^p(B,R^qf_*\mathbb Q)\cong H^p(B,R^{3-q}\check f_*\mathbb Q),
\]
它反转各行的次序。这类似于 Calabi--Yau 三维流形镜像对的 Hodge 菱形的对称性。

\SourceParagraph{ch06-6.4.4-018}{515}
现在，考虑第一个图中的 $E_2^{0,1}$ 项，即 $H^0(B,R^1f_*\mathbb Q)$。以 $E_2^{0,1}$ 为定义域或值域的两个映射都为零，因此，$E_2^{0,1}=E_\infty^{0,1}$。此外，$E_\infty^{1,0}=0$。所以，$H^1(X(B),\mathbb Q)=F^0\supseteq F^1$ 有滤过，满足 $F^0/F^1=E_2^{0,1}$，$F^1=E_2^{1,0}=0$。

\SourceParagraph{ch06-6.4.4-019}{515}
假定 $H^1(X(B),\mathbb Q)=0$。例如，$X(B)$ 的基本群有限时，就满足这一条件；有意思的 Calabi--Yau 三维流形通常如此，本书所考虑的环面簇中的 Calabi--Yau 三维超曲面也包括在内。那么，得到结论 $H^0(B,R^1f_*\mathbb Q)=0$。类似的论证说明，由 Poincaré 对偶，$H^5(X(B),\mathbb Q)=0$，因此还有 $H^3(B,R^2f_*\mathbb Q)=0$。

\SourceParagraph{ch06-6.4.4-020}{515}
若进一步有 $H^1(\check X(B),\mathbb Q)=0$，则同样有
\[
H^0(B,R^1\check f_*\mathbb Q)=H^3(B,R^2\check f_*\mathbb Q)=0.
\]
结合 $f$ 与 $\check f$ 的各个群之间的对称性，看到它们的 $E_2$ 页现在形如
\[
\BookImage{02_Unnumbered_Diagrams/Chapter_06/Ch06_Diagram_PDFp515_01_Bookp505.pdf}
\]
与
\[
\BookImage{02_Unnumbered_Diagrams/Chapter_06/Ch06_Diagram_PDFp515_02_Bookp505.pdf}
\]

\SourceParagraph{ch06-6.4.4-021}{515}
图景几乎完整了。在第一个谱序列中，至多有两个非零映射：
\[
d_2\colon H^0(B,R^3f_*\mathbb Q)\longrightarrow H^2(B,R^2f_*\mathbb Q),
\qquad
d_2'\colon H^1(B,R^1f_*\mathbb Q)\longrightarrow H^3(B,f_*\mathbb Q).
\]
要证明它们为零，再引入一个假设：$f$ 有截面 $\sigma\colon B\to X(B)$。在\BookSectionRef{sec:6.4.3}{6.4.3}中，实际构造的紧化 $f\colon X(B)\to B$，的确带有截面。可能还有其他环面纤维化，是 $X(B_0)$ 的扭曲的紧化，不过，在这里的论证中不考虑它们。

\SourceParagraph{ch06-6.4.4-022}{516}
注意，$E_\infty^{3,0}$ 为 $E_2^{3,0}=H^3(B,\mathbb Q)$ 的商，并且 $E_\infty^{3,0}\subseteq H^3(X(B),\mathbb Q)$。Leray 谱序列的一个一般事实是：复合
\[
E_2^{3,0}\longrightarrow E_\infty^{3,0}\hookrightarrow H^3(X(B),\mathbb Q)
\]
就是拉回 $f^*\colon H^3(B,\mathbb Q)\to H^3(X(B),\mathbb Q)$。不过，复合
\[
H^3(B,\mathbb Q)\xrightarrow{f^*}H^3(X(B),\mathbb Q)
\xrightarrow{\sigma^*}H^3(B,\mathbb Q)
\]
为恒等映射，因为 $f\circ\sigma$ 为恒等映射。因此，$f^*$ 为单射，于是 $d_2'=0$。

\SourceParagraph{ch06-6.4.4-023}{516}
同样，$E_\infty^{0,3}$ 为 $E_2^{0,3}$ 的子群。复合映射
\[
H^3(X(B),\mathbb Q)\longrightarrow E_\infty^{0,3}\longrightarrow E_2^{0,3}
=H^0(B,R^3f_*\mathbb Q)
\]
是如下自然映射：$R^3f_*\mathbb Q$ 为预层 $U\mapsto H^3(f^{-1}(U),\mathbb Q)$ 的关联层，因此，特别地，有自然映射 $H^3(f^{-1}(B),\mathbb Q)\to\Gamma(B,R^3f_*\mathbb Q)$。显然，$\sigma(B)$ 在 $H^3(X(B),\mathbb Q)$ 中的 Poincaré 对偶类，限制为 $H^0(B_0,R^3f_{0*}\mathbb Q)$ 中的非零元素。因此，$d_2$ 也为零。所以，$E_2$ 页中的全部微分都为零，而当 $r>2$ 时，也没有容纳非零 $d_r$ 的位置。因此，$E_2^{p,q}=E_\infty^{p,q}$。在这种情形中，称谱序列在 $E_2$ 页退化。

\SourceParagraph{ch06-6.4.4-024}{516}
将上述结论合在一起，在对 $f$ 与 $\check f$ 所作的各项假设下，得到同构
\[
\begin{aligned}
H^0(B,\mathbb Q)&\cong H^0(X(B),\mathbb Q),\\
H^1(B,R^1f_*\mathbb Q)&\cong H^2(X(B),\mathbb Q),\\
H^2(B,R^2f_*\mathbb Q)&\cong H^4(X(B),\mathbb Q),\\
H^3(B,R^3f_*\mathbb Q)&\cong H^6(X(B),\mathbb Q),
\end{aligned}
\]
对于 $\check f\colon\check X(B)\to B$，也有类似的陈述。此外，有滤过
\[
H^3(X(B),\mathbb Q)=F^0\supseteq F^1\supseteq F^2\supseteq F^3\supseteq0,
\]
满足
\[
\begin{aligned}
F^0/F^1&\cong H^0(B,R^3f_*\mathbb Q),\\
F^1/F^2&\cong H^1(B,R^2f_*\mathbb Q),\\
F^2/F^3&\cong H^2(B,R^1f_*\mathbb Q),\\
F^3&\cong H^3(B,f_*\mathbb Q).
\end{aligned}
\]
由 Poincaré 对偶，$\dim H^2(\check X(B),\mathbb Q)=\dim H^4(\check X(B),\mathbb Q)$，所以
\[
\dim(F^1/F^2)=\dim(F^2/F^3).
\]
将这一切合在一起，得到如下定理。

\SourceParagraph{ch06-6.4.4-025}{517}
\BookStatementNumber{103}
\begin{theorem}\label{thm:6.103}
令 $B$ 为紧、连通、带奇点的三维简单仿射流形，完整表示的像包含于 $\mathbb R^3\rtimes\operatorname{SL}_3(\mathbb Z)$。令 $f\colon X(B)\to B$ 与 $\check f\colon\check X(B)\to B$，为\BookSectionRef{sec:6.4.3}{6.4.3}中给出的、带有截面的 $f_0\colon X(B_0)\to B_0$ 与 $\check f_0\colon\check X(B_0)\to B_0$ 的紧化。进一步假定 $H^1(X(B),\mathbb Q)=H^1(\check X(B),\mathbb Q)=0$。则
\[
\dim H^{2*}(X(B),\mathbb Q)=\dim H^3(\check X(B),\mathbb Q),
\]
并且
\[
\dim H^{2*}(\check X(B),\mathbb Q)=\dim H^3(X(B),\mathbb Q).
\]
这里，$H^{2*}$ 表示偶数次数的上同调。特别地，若 $X(B)$ 与 $\check X(B)$ 带有使其成为 Calabi--Yau 三维流形的复结构，则
\[
h^{1,1}(X(B))=h^{1,2}(\check X(B)),\qquad
h^{1,2}(X(B))=h^{1,1}(\check X(B)).
\]
\end{theorem}

\SourceParagraph{ch06-6.4.4-026}{517}
将 $\mathbb Q$ 换成 $\mathbb Z$，Leray 谱序列仍然有意义。不过，在这种情形中，情况更加微妙，因为谱序列的各个项可能出现挠元。这是一个相当有意思的现象，文献\BookCite{ref201}第3节对此作了解释。

\SourceParagraph{ch06-6.4.4-027}{517}
对 $X(B)$ 与 $\check X(B)$ 的上同调所作的分析，提出了一个有意思的问题：Leray 谱序列在 $H^3$ 上给出的滤过，究竟是什么？考察紧化族中的单值变换权滤过，可以给出答案。这里将概述这一点，不深入讨论细节。

\SourceParagraph{ch06-6.4.4-028}{517}
假定 $B_0$ 实际上具有整仿射结构，那么，可以像\BookSectionRef{sec:6.2.4}{6.2.4}中一样，在 $B_0'=B_0\times\mathbb R_{>0}$ 上构造整仿射结构，从而得到族 $X(B_0')\to X(\mathbb R_{>0})=D^*$。\BookRef{prop:6.32}{命题6.32}说明，这个族的单值变换，由截面给出的平移 $T_\delta\colon X(B_0)\to X(B_0)$ 描述；其中，截面由仿射结构确定，即本质上为展开映射 $\delta\colon\widetilde B_0\to M_{\mathbb R}$ 的图像。可以证明，将截面换成与之同伦的截面以后，这一平移可以延伸为映射 $T_\delta\colon X(B)\to X(B)$。

\SourceParagraph{ch06-6.4.4-029}{518}
这启发下面的定理。设处于\BookRef{thm:6.103}{定理6.103}的情形，给定截面 $\sigma\colon B\to X(B)$，诱导自同构 $T_\sigma\colon X(B)\to X(B)$。这始终涉及零截面 $\sigma_0\colon B\to X(B)$ 的选择。现在，考虑 $\sigma-\sigma_0$ 的 Poincaré 对偶上同调类 $[\sigma-\sigma_0]\in H^3(X(B),\mathbb Q)$。由于 $[\sigma-\sigma_0]$ 在 $X(B)\to B$ 的任意纤维上的限制都为零，所以，$[\sigma-\sigma_0]\in F^1$，进而在模 $F^2$ 的意义下定义类
\[
D_\sigma\in H^1(B,R^2f_*\mathbb Q)
\cong H^1(B,R^1\check f_*\mathbb Q)
\cong H^2(\check X(B),\mathbb Q).
\]

\SourceParagraph{ch06-6.4.4-030}{518}
\BookStatementNumber{104}
\begin{theorem}\label{thm:6.104}
令 $B$、$f\colon X(B)\to B$ 与 $\check f\colon\check X(B)\to B$ 满足\BookRef{thm:6.103}{定理6.103}的假设。设 $\sigma\colon B\to X(B)$ 为截面，诱导同胚 $T_\sigma\colon X(B)\to X(B)$，在 $X(B_0)$ 上由 $T_\sigma(v)=v+\sigma(f_0(v))$ 给出。那么，
\[
T_\sigma^*\colon H^3(X(B),\mathbb Q)\longrightarrow H^3(X(B),\mathbb Q)
\]
满足 $(T_\sigma^*-I)(F^i)\subseteq F^{i+1}$，诱导映射
\[
\begin{aligned}
T_\sigma^*-I\colon F^i/F^{i+1}
&\cong H^i(B,R^{3-i}f_*\mathbb Q)\\
&\longrightarrow F^{i+1}/F^{i+2}
\cong H^{i+1}(B,R^{2-i}f_*\mathbb Q).
\end{aligned}
\]
此外，下图在相差符号的意义下交换：
\[
\BookImage{02_Unnumbered_Diagrams/Chapter_06/Ch06_Diagram_PDFp518_01_Bookp508.pdf}
\]
\end{theorem}
\begin{proof}
\SourceParagraph{ch06-6.4.4-031}{518}
见文献\BookCite{ref200}的定理4.1。
\end{proof}

\SourceParagraph{ch06-6.4.4-032}{518}
利用这一定理，在所讨论的大复结构极限退化情形中，可以证明：按相应的重编号，滤过 $F^\bullet$ 与 $T_\delta$ 在 $H^3(X(B),\mathbb Q)$ 上诱导的单值变换权滤过相同。此外，可以利用定理中的交换图，证明镜像的 $(1,1)$ 与 $(1,2)$ 型 Yukawa 耦合，在零阶一致；也就是说，$\check X(B)$ 上的拓扑 $(1,1)$ 型耦合，与镜像 $X(B)$ 的复模空间上的 $(1,2)$ 型耦合所预测的结果一致。更多细节，见文献\BookCite{ref200}。
